\documentclass[10pt]{article}
\usepackage[letterpaper,margin=1.8cm]{geometry}
\usepackage[T1]{fontenc}
\usepackage{newtxtext,newtxmath}
\usepackage{booktabs}
\usepackage{graphicx}
\usepackage{amsmath}
\usepackage{xspace}
\usepackage{array}
\usepackage{longtable}
\usepackage{xcolor}
\usepackage{hyperref}
\usepackage[capitalise,nameinlink]{cleveref}
\providecommand{\checkmark}{$\surd$}

% method and benchmark names
\newcommand{\baro}{BARO\xspace}
\newcommand{\circa}{CIRCA\xspace}
\newcommand{\ediag}{$\epsilon$-Diagnosis\xspace}
\newcommand{\rcd}{RCD\xspace}
\newcommand{\causalrca}{CausalRCA\xspace}
\newcommand{\cfa}{counterfactual attribution\xspace}
\newcommand{\simplerca}{SimpleRCA\xspace}
\newcommand{\tracerca}{TraceRCA\xspace}
\newcommand{\microrank}{MicroRank\xspace}
\newcommand{\maxz}{max-$|Z|$\xspace}
\newcommand{\alertcount}{alert-count\xspace}
\newcommand{\cdmin}{CD-1min\xspace}
\newcommand{\rcaeval}{RCAEval\xspace}
\newcommand{\openrca}{OpenRCA\xspace}
\newcommand{\petshop}{PetShop\xspace}
\newcommand{\acc}{Acc@1\xspace}
\newcommand{\ci}[2]{$[#1, #2]$}
\newcommand{\todo}[1]{\textcolor{red}{[TODO: #1]}}

% generated by paper/scripts/make_numbers.py from paper/data/ -- do not edit
\newcommand{\NlongRows}{19,893\xspace}
\newcommand{\longShaShort}{78b24273\xspace}
\newcommand{\NpairAnalyses}{204\xspace}
\newcommand{\NcwDiff}{1\xspace}
\newcommand{\NReOnePubPairs}{10\xspace}
\newcommand{\NReOnePubRev}{1\xspace}
\newcommand{\NReOnePubRevCI}{0\xspace}
\newcommand{\NReOneAllPairs}{28\xspace}
\newcommand{\NReOneAllRev}{8\xspace}
\newcommand{\NReOneAllRevCI}{4\xspace}
\newcommand{\NOpenrcaPubPairs}{3\xspace}
\newcommand{\NOpenrcaPubRev}{1\xspace}
\newcommand{\NOpenrcaPubRevCI}{0\xspace}
\newcommand{\NOpenrcaPubPIn}{3\xspace}
\newcommand{\NOpenrcaPubPIcross}{3\xspace}
\newcommand{\NOpenrcaAllPairs}{15\xspace}
\newcommand{\NOpenrcaAllRev}{7\xspace}
\newcommand{\NOpenrcaAllRevCI}{0\xspace}
\newcommand{\NOpenrcaAllPIn}{15\xspace}
\newcommand{\NOpenrcaAllPIcross}{11\xspace}
\newcommand{\NPetshopPubPairs}{10\xspace}
\newcommand{\NPetshopPubRev}{5\xspace}
\newcommand{\NPetshopPubRevCI}{2\xspace}
\newcommand{\NPetshopPubPIn}{9\xspace}
\newcommand{\NPetshopPubPIcross}{8\xspace}
\newcommand{\NPetshopAllPairs}{55\xspace}
\newcommand{\NPetshopAllRev}{12\xspace}
\newcommand{\NPetshopAllRevCI}{4\xspace}
\newcommand{\NPetshopAllPIn}{47\xspace}
\newcommand{\NPetshopAllPIcross}{32\xspace}
\newcommand{\NReTwoPubPairs}{28\xspace}
\newcommand{\NReTwoPubRev}{10\xspace}
\newcommand{\NReTwoPubRevCI}{6\xspace}
\newcommand{\NReTwoAllPairs}{55\xspace}
\newcommand{\NReTwoAllRev}{21\xspace}
\newcommand{\NReTwoAllRevCI}{13\xspace}
\newcommand{\NMainPubPairs}{23\xspace}
\newcommand{\NMainPubRev}{7\xspace}
\newcommand{\NMainPubRevCI}{2\xspace}
\newcommand{\NMainAllPairs}{98\xspace}
\newcommand{\NMainAllRev}{27\xspace}
\newcommand{\NMainAllRevCI}{8\xspace}
\newcommand{\ReOneCrRcdPooled}{$-$0.030\xspace}
\newcommand{\ReOneCrRcdPooledCI}{[$-$0.077, 0.017]\xspace}
\newcommand{\ReOneCrRcdOB}{$+$0.075\xspace}
\newcommand{\ReOneCrRcdSS}{$-$0.213\xspace}
\newcommand{\ReOneCrRcdTT}{$+$0.048\xspace}
\newcommand{\ReOneCrRcdSeedZero}{$+$0.027\xspace}
\newcommand{\ReOneCrRcdSeedZeroCI}{1\xspace}
\newcommand{\ReOneCrRcdSeedOne}{$-$0.035\xspace}
\newcommand{\ReOneCrRcdSeedOneCI}{1\xspace}
\newcommand{\ReOneCrRcdSeedTwo}{$-$0.083\xspace}
\newcommand{\ReOneCrRcdSeedTwoCI}{0\xspace}
\newcommand{\ReOneBaroMinMargin}{0.147\xspace}
\newcommand{\ReOneBaroMaxMargin}{0.656\xspace}
\newcommand{\ExBelowNoRevMargin}{0.251\xspace}
\newcommand{\ExBelowNoRevSD}{0.293\xspace}
\newcommand{\ExAboveRevMargin}{0.236\xspace}
\newcommand{\ExAboveRevSD}{0.195\xspace}
\newcommand{\ExAboveRevHigh}{$-$0.038\xspace}
\newcommand{\OpenrcaBaroMean}{0.127\xspace}
\newcommand{\OpenrcaRcdMean}{0.100\xspace}
\newcommand{\OpenrcaEdiagMean}{0.041\xspace}
\newcommand{\OpenrcaBaroRcdPooled}{$+$0.027\xspace}
\newcommand{\OpenrcaBaroRcdPooledCI}{[0.002, 0.051]\xspace}
\newcommand{\OpenrcaBaroRcdTel}{$-$0.010\xspace}
\newcommand{\OpenrcaBaroRcdTelCI}{[$-$0.058, 0.035]\xspace}
\newcommand{\OpenrcaBaroRcdSeedsWithRev}{2\xspace}
\newcommand{\PetshopRcdHighHits}{0\xspace}
\newcommand{\PetshopHighN}{26\xspace}
\newcommand{\PetshopBaroRcdHigh}{$+$0.154\xspace}
\newcommand{\PetshopBaroRcdHighCI}{[0.038, 0.308]\xspace}
\newcommand{\PetshopBaroRcdPooled}{$-$0.135\xspace}
\newcommand{\PetshopBaroRcdPooledCI}{[$-$0.300, 0.041]\xspace}
\newcommand{\PetshopCfaRcdHigh}{$+$0.308\xspace}
\newcommand{\PetshopCfaRcdHighCI}{[0.154, 0.500]\xspace}
\newcommand{\PetshopCfaRcdPooled}{$-$0.127\xspace}
\newcommand{\PetshopCfaRcdPooledCI}{[$-$0.300, 0.058]\xspace}
\newcommand{\PetshopPubRevNmin}{3\xspace}
\newcommand{\PetshopPubRevCINmin}{1\xspace}
\newcommand{\PetshopAllRevNmin}{7\xspace}
\newcommand{\PetshopAllRevCINmin}{1\xspace}
\newcommand{\PetshopRankedCorrMean}{0.639\xspace}
\newcommand{\PetshopCircaMean}{0.445\xspace}
\newcommand{\NPubBelow}{8\xspace}
\newcommand{\NPubBelowRev}{6\xspace}
\newcommand{\NPubBelowRevCI}{2\xspace}
\newcommand{\NPubAbove}{15\xspace}
\newcommand{\NPubAboveRev}{1\xspace}
\newcommand{\NPubAboveRevCI}{0\xspace}
\newcommand{\NAllBelow}{34\xspace}
\newcommand{\NAllBelowRev}{25\xspace}
\newcommand{\NAllBelowRevCI}{8\xspace}
\newcommand{\NAllAbove}{64\xspace}
\newcommand{\NAllAboveRev}{2\xspace}
\newcommand{\NAllAboveRevCI}{0\xspace}
\newcommand{\NTabRcaevalPairs}{55\xspace}
\newcommand{\NTabRcaevalBelow}{23\xspace}
\newcommand{\NTabRcaevalBelowRev}{21\xspace}
\newcommand{\NTabRcaevalAbove}{32\xspace}
\newcommand{\NTabRcaevalAboveRev}{2\xspace}
\newcommand{\NTabPetshopPairs}{15\xspace}
\newcommand{\NTabPetshopBelow}{8\xspace}
\newcommand{\NTabPetshopBelowRev}{6\xspace}
\newcommand{\NTabPetshopAbove}{7\xspace}
\newcommand{\NTabPetshopAboveRev}{0\xspace}
\newcommand{\NTabBelow}{31\xspace}
\newcommand{\NTabBelowRev}{27\xspace}
\newcommand{\NTabAbove}{39\xspace}
\newcommand{\NTabAboveRev}{2\xspace}
\newcommand{\NCrossPub}{18\xspace}
\newcommand{\NCrossPubAgree}{12\xspace}
\newcommand{\NCrossPubCIopp}{2\xspace}
\newcommand{\NCrossAll}{63\xspace}
\newcommand{\NCrossAllAgree}{43\xspace}
\newcommand{\NCrossAllCIopp}{9\xspace}
\newcommand{\CrossBaroCircaSSOne}{$+$0.144\xspace}
\newcommand{\CrossBaroCircaSSOneCI}{[0.064, 0.232]\xspace}
\newcommand{\CrossBaroCircaSSTwo}{$-$0.711\xspace}
\newcommand{\CrossBaroCircaSSTwoCI}{[$-$0.800, $-$0.611]\xspace}
\newcommand{\CrossBaroCircaOBOne}{$+$0.064\xspace}
\newcommand{\CrossBaroCircaOBOneCI}{[$-$0.032, 0.160]\xspace}
\newcommand{\CrossBaroCircaOBTwo}{$-$0.522\xspace}
\newcommand{\CrossBaroCircaOBTwoCI}{[$-$0.644, $-$0.400]\xspace}
\newcommand{\CrossBaroCrSSOne}{$+$0.672\xspace}
\newcommand{\CrossBaroCrSSOneCI}{[0.592, 0.752]\xspace}
\newcommand{\CrossBaroCrSSTwo}{$-$0.081\xspace}
\newcommand{\CrossBaroCrSSTwoCI}{[$-$0.163, $-$0.004]\xspace}
\newcommand{\CrossBaroEdSSOne}{$+$0.704\xspace}
\newcommand{\CrossBaroEdSSOneCI}{[0.616, 0.784]\xspace}
\newcommand{\CrossBaroEdSSTwo}{$-$0.089\xspace}
\newcommand{\CrossBaroEdSSTwoCI}{[$-$0.189, 0.000]\xspace}
\newcommand{\NCrossPubAgreeReOneFaults}{12\xspace}
\newcommand{\RegReTwoWinnerPooled}{0.663\xspace}
\newcommand{\RegReTwoTT}{0.256\xspace}
\newcommand{\RegReTwoTTbest}{0.800\xspace}
\newcommand{\RegReTwoTTwinner}{0.544\xspace}
\newcommand{\RegOpenrcaTel}{0.010\xspace}
\newcommand{\RegPetshopHigh}{0.038\xspace}
\newcommand{\RegPetshopTmpOne}{0.125\xspace}
\newcommand{\RegReOneAllSS}{0.080\xspace}
\newcommand{\RegPetshopAllTmpTwo}{0.375\xspace}
\newcommand{\InSSBaroSimple}{0.848\xspace}
\newcommand{\InSSBaroRaw}{0.200\xspace}
\newcommand{\InSSCircaSimple}{0.704\xspace}
\newcommand{\InSSCircaRaw}{0.200\xspace}
\newcommand{\InSSRcdSimple}{0.389\xspace}
\newcommand{\InSSRcdRaw}{0.296\xspace}
\newcommand{\InSSCdminSimple}{0.928\xspace}
\newcommand{\InSSCdminRaw}{0.736\xspace}
\newcommand{\InSSAlertSimple}{0.376\xspace}
\newcommand{\InSSAlertRaw}{0.432\xspace}
\newcommand{\InSSMaxzSimple}{0.744\xspace}
\newcommand{\InSSMaxzRaw}{0.608\xspace}
\newcommand{\InSSEdiagSimple}{0.144\xspace}
\newcommand{\InSSEdiagRaw}{0.168\xspace}
\newcommand{\InTTBaroSimple}{0.560\xspace}
\newcommand{\InTTBaroRaw}{0.120\xspace}
\newcommand{\InTTCircaSimple}{0.328\xspace}
\newcommand{\InTTCircaRaw}{0.000\xspace}
\newcommand{\InTTRcdSimple}{0.088\xspace}
\newcommand{\InTTRcdRaw}{0.091\xspace}
\newcommand{\InTTCdminSimple}{0.280\xspace}
\newcommand{\InTTCdminRaw}{0.456\xspace}
\newcommand{\InTTAlertSimple}{0.048\xspace}
\newcommand{\InTTAlertRaw}{0.192\xspace}
\newcommand{\InTTMaxzSimple}{0.424\xspace}
\newcommand{\InTTMaxzRaw}{0.448\xspace}
\newcommand{\InTTEdiagSimple}{0.000\xspace}
\newcommand{\InTTEdiagRaw}{0.000\xspace}
\newcommand{\InReTwoCdminDropMin}{0.333\xspace}
\newcommand{\InReTwoCdminDropMax}{0.544\xspace}
\newcommand{\InReTwoTTMmbaroSimple}{0.689\xspace}
\newcommand{\InReTwoTTMmbaroRaw}{0.144\xspace}
\newcommand{\InReOnePubComp}{12\xspace}
\newcommand{\InReOnePubSign}{5\xspace}
\newcommand{\InReOnePubSignSS}{3\xspace}
\newcommand{\InReOnePubCompSS}{6\xspace}
\newcommand{\InReOnePubSignTT}{2\xspace}
\newcommand{\InReOnePubCompTT}{6\xspace}
\newcommand{\InReOnePubRevRaw}{2\xspace}
\newcommand{\InReOnePubRevCIRaw}{1\xspace}
\newcommand{\InReOnePubRevOk}{0\xspace}
\newcommand{\InReOnePubRevCIOk}{0\xspace}
\newcommand{\InReOnePubRevRawOnly}{2\xspace}
\newcommand{\InReOnePubRevOkOnly}{0\xspace}
\newcommand{\InReOneAllComp}{42\xspace}
\newcommand{\InReOneAllSign}{16\xspace}
\newcommand{\InReOneAllSignSS}{7\xspace}
\newcommand{\InReOneAllCompSS}{21\xspace}
\newcommand{\InReOneAllSignTT}{9\xspace}
\newcommand{\InReOneAllCompTT}{21\xspace}
\newcommand{\InReOneAllRevRaw}{10\xspace}
\newcommand{\InReOneAllRevCIRaw}{6\xspace}
\newcommand{\InReOneAllRevOk}{6\xspace}
\newcommand{\InReOneAllRevCIOk}{2\xspace}
\newcommand{\InReOneAllRevRawOnly}{9\xspace}
\newcommand{\InReOneAllRevOkOnly}{5\xspace}
\newcommand{\InReTwoComp}{18\xspace}
\newcommand{\InReTwoSign}{5\xspace}
\newcommand{\InCircaRcdSSOk}{$+$0.315\xspace}
\newcommand{\InCircaRcdSSOkCI}{[0.195, 0.432]\xspace}
\newcommand{\InCircaRcdSSRaw}{$-$0.096\xspace}
\newcommand{\InCircaRcdSSRawCI}{[$-$0.184, $-$0.005]\xspace}
\newcommand{\InCircaRcdTTOk}{$+$0.240\xspace}
\newcommand{\InCircaRcdTTOkCI}{[0.149, 0.328]\xspace}
\newcommand{\InCircaRcdTTRaw}{$-$0.091\xspace}
\newcommand{\InCircaRcdTTRawCI}{[$-$0.139, $-$0.048]\xspace}
\newcommand{\InTauSS}{0.183\xspace}
\newcommand{\DefPubComp}{30\xspace}
\newcommand{\DefPubSign}{2\xspace}
\newcommand{\DefPubRevShipOnly}{2\xspace}
\newcommand{\DefPubRevCIShipOnly}{1\xspace}
\newcommand{\DefPubPooledFlips}{0\xspace}
\newcommand{\DefAllComp}{84\xspace}
\newcommand{\DefAllSign}{3\xspace}
\newcommand{\DefAllRevShipOnly}{2\xspace}
\newcommand{\DefAllRevCIShipOnly}{1\xspace}
\newcommand{\DefAllPooledFlips}{1\xspace}
\newcommand{\DefCircaFixed}{0.672\xspace}
\newcommand{\DefCircaShipped}{0.432\xspace}
\newcommand{\DefCircaFixedAff}{0.820\xspace}
\newcommand{\DefCircaShippedAff}{0.220\xspace}
\newcommand{\DefCrCircaFixed}{$-$0.056\xspace}
\newcommand{\DefCrCircaFixedCI}{[$-$0.144, 0.035]\xspace}
\newcommand{\DefCrCircaShipped}{$+$0.184\xspace}
\newcommand{\DefCrCircaShippedCI}{[0.077, 0.291]\xspace}
\newcommand{\DefCircaRcdFixed}{$+$0.131\xspace}
\newcommand{\DefCircaRcdShipped}{$-$0.115\xspace}
\newcommand{\CovTraceValid}{73\xspace}
\newcommand{\CovTraceAll}{0.533\xspace}
\newcommand{\CovTraceValidAcc}{0.658\xspace}
\newcommand{\CovCrReOneTTSeedZero}{0.160\xspace}
\newcommand{\CovCrReOneTTSeedOne}{0.208\xspace}
\newcommand{\CovCrReOneTTSeedTwo}{0.040\xspace}
\newcommand{\CovCrReTwoTTSeedZero}{0.267\xspace}
\newcommand{\CovCrReTwoTTSeedOne}{0.133\xspace}
\newcommand{\CovCrReTwoTTSeedTwo}{0.222\xspace}
\newcommand{\CovCrReTwoTTMean}{0.207\xspace}
\newcommand{\CovMicroValid}{73\xspace}
\newcommand{\CovCrSeedOneValid}{47\xspace}
\newcommand{\CovEdTelValid}{19\xspace}
\newcommand{\CovEdTelN}{51\xspace}
\newcommand{\CovEdHighValid}{16\xspace}
\newcommand{\NCovBelowOne}{25\xspace}
\newcommand{\NCovCells}{149\xspace}
\newcommand{\DenReTwoPubComp}{71\xspace}
\newcommand{\DenReTwoPubSign}{1\xspace}
\newcommand{\DenReTwoAllComp}{146\xspace}
\newcommand{\DenReTwoAllSign}{2\xspace}
\newcommand{\DenOtherSign}{0\xspace}
\newcommand{\DenOtherComp}{364\xspace}
\newcommand{\DenCircaTraceAll}{$+$0.011\xspace}
\newcommand{\DenCircaTraceValid}{$-$0.113\xspace}
\newcommand{\FwReOnePubComp}{30\xspace}
\newcommand{\FwReOnePubSign}{0\xspace}
\newcommand{\FwReOnePubSignCI}{0\xspace}
\newcommand{\FwReOneAllComp}{84\xspace}
\newcommand{\FwReOneAllSign}{0\xspace}
\newcommand{\FwReOneAllSignCI}{0\xspace}
\newcommand{\FwReTwoPubComp}{71\xspace}
\newcommand{\FwReTwoPubSign}{0\xspace}
\newcommand{\FwReTwoPubSignCI}{0\xspace}
\newcommand{\FwReTwoAllComp}{146\xspace}
\newcommand{\FwReTwoAllSign}{0\xspace}
\newcommand{\FwReTwoAllSignCI}{0\xspace}
\newcommand{\FwOpenrcaPubComp}{12\xspace}
\newcommand{\FwOpenrcaPubSign}{3\xspace}
\newcommand{\FwOpenrcaPubSignCI}{0\xspace}
\newcommand{\FwOpenrcaAllComp}{60\xspace}
\newcommand{\FwOpenrcaAllSign}{13\xspace}
\newcommand{\FwOpenrcaAllSignCI}{0\xspace}
\newcommand{\FwPetshopPubComp}{40\xspace}
\newcommand{\FwPetshopPubSign}{1\xspace}
\newcommand{\FwPetshopPubSignCI}{0\xspace}
\newcommand{\FwPetshopAllComp}{220\xspace}
\newcommand{\FwPetshopAllSign}{2\xspace}
\newcommand{\FwPetshopAllSignCI}{0\xspace}
\newcommand{\ScStrictPubComp}{12\xspace}
\newcommand{\ScStrictPubSign}{4\xspace}
\newcommand{\ScStrictPubPooledFlips}{0\xspace}
\newcommand{\ScStrictAllComp}{60\xspace}
\newcommand{\ScStrictAllSign}{11\xspace}
\newcommand{\ScStrictAllPooledFlips}{2\xspace}
\newcommand{\ScHalfPubComp}{12\xspace}
\newcommand{\ScHalfPubSign}{0\xspace}
\newcommand{\ScHalfPubPooledFlips}{0\xspace}
\newcommand{\ScHalfAllComp}{60\xspace}
\newcommand{\ScHalfAllSign}{3\xspace}
\newcommand{\ScHalfAllPooledFlips}{0\xspace}
\newcommand{\ScStrictPubToZero}{4\xspace}
\newcommand{\ScMktOneBaroStrict}{0.000\xspace}
\newcommand{\ScMktOneBaroPartial}{0.057\xspace}
\newcommand{\ScMktOneRcdStrict}{0.000\xspace}
\newcommand{\ScMktOneRcdPartial}{0.051\xspace}
\newcommand{\ScMktOneEdiagStrict}{0.000\xspace}
\newcommand{\ScMktOneEdiagPartial}{0.038\xspace}
\newcommand{\LineCw}{$+$0.038\xspace}
\newcommand{\LineSe}{$+$0.027\xspace}
\newcommand{\LineSD}{0.036\xspace}
\newcommand{\LineWins}{3\xspace}
\newcommand{\LineLosses}{1\xspace}
\newcommand{\LineTies}{0\xspace}
\newcommand{\LineCovBaro}{1.00\xspace}
\newcommand{\LineCovRcd}{0.96\xspace}
\newcommand{\LineWorst}{$-$0.010\xspace}
\newcommand{\LineWorstCI}{[$-$0.058, 0.035]\xspace}
\newcommand{\NBonfComp}{82\xspace}
\newcommand{\NBonfRevCI}{1\xspace}
\newcommand{\NBonfCIExcl}{46\xspace}
\newcommand{\NBonfCIAdjExcl}{31\xspace}
\newcommand{\SurveyN}{23\xspace}
\newcommand{\SurveyMulti}{17\xspace}
\newcommand{\SurveySingle}{6\xspace}
\newcommand{\SurveyRelAny}{7\xspace}
\newcommand{\SurveyRelDemo}{3\xspace}
\newcommand{\SurveyRelOwn}{3\xspace}
\newcommand{\SurveyRelUndet}{2\xspace}
\newcommand{\SurveyRelFullCheckable}{21\xspace}
\newcommand{\IsqOpenrcaPubMin}{49\xspace}
\newcommand{\IsqOpenrcaPubMax}{74\xspace}
\newcommand{\IsqPetshopPubMin}{0\xspace}
\newcommand{\IsqPetshopPubMax}{82\xspace}
\newcommand{\IntReOnePubFit}{6\xspace}
\newcommand{\IntReOnePubSig}{4\xspace}
\newcommand{\IntReOnePubSep}{4\xspace}
\newcommand{\IntOpenrcaPubFit}{3\xspace}
\newcommand{\IntOpenrcaPubSig}{0\xspace}
\newcommand{\IntOpenrcaPubSep}{0\xspace}
\newcommand{\IntPetshopPubFit}{3\xspace}
\newcommand{\IntPetshopPubSig}{0\xspace}
\newcommand{\IntPetshopPubSep}{7\xspace}
\newcommand{\XNMethods}{6\xspace}
\newcommand{\XNPairs}{15\xspace}
\newcommand{\XK}{11\xspace}
\newcommand{\XNBothSigns}{10\xspace}
\newcommand{\XNOneSign}{5\xspace}
\newcommand{\XNRevCI}{6\xspace}
\newcommand{\XNBelow}{12\xspace}
\newcommand{\XNFit}{11\xspace}
\newcommand{\XIsqMin}{80\xspace}
\newcommand{\XIsqMax}{98\xspace}
\newcommand{\XNPIcross}{11\xspace}
\newcommand{\XNLRTsig}{14\xspace}
\newcommand{\XNSep}{14\xspace}
\newcommand{\XJointLRT}{264.4\xspace}
\newcommand{\XJointDf}{50\xspace}
\newcommand{\XLosoNAllRev}{10\xspace}
\newcommand{\XLosoMin}{2\xspace}
\newcommand{\XLosoMax}{10\xspace}
\newcommand{\XLosoMaxRegret}{0.539\xspace}
\newcommand{\XLosoMeanRegretMax}{0.182\xspace}
\newcommand{\XBaroMaxzPooled}{$+$0.004\xspace}
\newcommand{\XBaroMaxzLoso}{10\xspace}
\newcommand{\XLofoN}{30\xspace}
\newcommand{\XLofoBoth}{28\xspace}
\newcommand{\XLofoPIn}{28\xspace}
\newcommand{\XLofoPI}{28\xspace}
\newcommand{\XLksoN}{550\xspace}
\newcommand{\XLksoBoth}{549\xspace}
\newcommand{\XLksoPer}{55\xspace}
\newcommand{\XPoolFlipPairs}{2\xspace}
\newcommand{\XBootBothMin}{0.997\xspace}
\newcommand{\HNDec}{36\xspace}
\newcommand{\HErr}{8\xspace}
\newcommand{\HTies}{2\xspace}
\newcommand{\HStrict}{2\xspace}
\newcommand{\HErrCw}{10\xspace}
\newcommand{\HRec}{19\xspace}
\newcommand{\HRecErr}{4\xspace}
\newcommand{\HFlag}{17\xspace}
\newcommand{\HFlagErr}{4\xspace}
\newcommand{\HFlagTies}{2\xspace}
\newcommand{\HBudgetErr}{3\xspace}
\newcommand{\HBudgetCwErr}{3\xspace}
\newcommand{\HReOneRec}{9\xspace}
\newcommand{\HReOneRecErr}{1\xspace}
\newcommand{\HReOneFlag}{3\xspace}
\newcommand{\HReOneFlagErr}{2\xspace}
\newcommand{\HOpenrcaRec}{4\xspace}
\newcommand{\HOpenrcaRecErr}{1\xspace}
\newcommand{\HOpenrcaFlag}{4\xspace}
\newcommand{\HOpenrcaFlagErr}{0\xspace}
\newcommand{\HPetshopRec}{6\xspace}
\newcommand{\HPetshopRecErr}{2\xspace}
\newcommand{\HPetshopFlag}{10\xspace}
\newcommand{\HPetshopFlagErr}{2\xspace}
\newcommand{\LosoRegPubN}{3\xspace}
\newcommand{\LosoRegPubOf}{11\xspace}
\newcommand{\LosoRegPubMax}{0.125\xspace}
\newcommand{\LosoRegAllN}{9\xspace}
\newcommand{\LosoRegAllOf}{11\xspace}
\newcommand{\LosoRegAllMax}{0.375\xspace}
\newcommand{\NAudPubPairs}{20\xspace}
\newcommand{\NAudPubRev}{6\xspace}
\newcommand{\NAudPubRevCI}{2\xspace}
\newcommand{\NAudPubEdiagPairs}{8\xspace}
\newcommand{\NMainPubEdiagPairs}{10\xspace}
\newcommand{\NMainPubRevNmin}{5\xspace}
\newcommand{\NMainPubRevCINmin}{1\xspace}
\newcommand{\NAudPubRevNmin}{4\xspace}
\newcommand{\NAudPubRevCINmin}{1\xspace}
\newcommand{\HNonTie}{34\xspace}
\newcommand{\HRandErr}{17\xspace}
\newcommand{\RegRandReOne}{0.317\xspace}
\newcommand{\RegWinnerMeanReOne}{0.000\xspace}
\newcommand{\RegRandOpenrca}{0.040\xspace}
\newcommand{\RegWinnerMeanOpenrca}{0.002\xspace}
\newcommand{\RegRandPetshop}{0.241\xspace}
\newcommand{\RegWinnerMeanPetshop}{0.041\xspace}
\newcommand{\BaroReTwoOBTopMetricN}{77\xspace}
\newcommand{\BaroReTwoOBAccThree}{0.878\xspace}
\newcommand{\BaroReTwoSSTopMetricN}{81\xspace}
\newcommand{\BaroReTwoSSAccThree}{0.956\xspace}
\newcommand{\BaroAccReOneOB}{0.736\xspace}
\newcommand{\BaroAccReOneSS}{0.848\xspace}
\newcommand{\BaroAccReTwoOB}{0.144\xspace}
\newcommand{\BaroAccReTwoSS}{0.067\xspace}
\newcommand{\ReOneOBCols}{51\xspace}
\newcommand{\ReOneSSCols}{60\xspace}
\newcommand{\ReTwoOBCols}{73\xspace}
\newcommand{\ReTwoSSCols}{76\xspace}
\newcommand{\ReOneSSRows}{721\xspace}
\newcommand{\ReTwoRows}{1441\xspace}


\sloppy
\begin{document}

\title{Additional Tables and Diagnostics for ``What a Pooled Leaderboard Does Not Say: Auditing the Scope and Provenance of Method Comparisons in Microservice RCA Benchmarks''}
\author{Anonymous Author(s)}
\date{}
\maketitle

This document gives additional definitions, run details, diagnostics, and complete result tables for the paper.
Section numbers of the paper are written as ``paper \S''.
All tables in \cref{sec:full} are generated from the analysis outputs of the frozen per-case table; the diagnostic tables in \cref{sec:simplerca,sec:table6,sec:window} are generated from the stored outputs of the diagnostic runs named in their captions.
Accuracy is always computed over all cases of a subsystem, with invalid outputs scored 0, unless a column says otherwise.

\section{Probe Definitions}
\label{sec:probes}

All three probes receive the same input as the published methods on each family: a table of metric time series, the injection (or query) time $t_0$, and a mapping from each metric to its component (the service prefix of the metric name on \rcaeval, the node on \petshop, and the component parser of the \openrca wrapper in \cref{sec:openrca}).
Let $x_m(t)$ be metric $m$, $B$ the time points before $t_0$ in the analysis window, and $A$ the time points at or after $t_0$.
Non-numeric columns and the time column are ignored.

\textbf{\maxz.}
For each metric, $\mu_m$ and $\sigma_m$ are the mean and the population standard deviation of $x_m$ over $B$, and the score is $z_m = \max_{t\in A} |x_m(t)-\mu_m|/\sigma_m$.
Metrics with $\sigma_m = 0$, or with no finite value in $A$, score 0.
Metrics are ranked by $z_m$ in descending order, ties by name; the answer is the component of the first metric (on \openrca, of the first metric that belongs to a candidate component; \cref{sec:openrca}).

\textbf{\alertcount.}
With $z_m$ as above, a metric raises an alert if $z_m > 3$.
Components are ranked by their number of alerting metrics in descending order, ties by name; within a component, metrics are ordered by $z_m$.
The answer is the first component.

\textbf{\cdmin.}
For each metric with at least 10 finite baseline values, $\mu_m$ and $s_m$ are the mean and sample standard deviation over $B$ (metrics with $s_m < 10^{-8}$ are skipped).
The peak $p_m$ is the value in $A$ farthest from $\mu_m$ (the maximum or the minimum).
The metric is flagged if both $|p_m-\mu_m|/s_m \ge 3$ and $100\,|p_m-\mu_m|/(|\mu_m|+10^{-8}) \ge 5$.
Each component receives the largest relative change among its flagged metrics, and components are ranked by it (ties by the largest $z$, then by name).
If no metric is flagged, the output is empty and counts as invalid.
On \openrca, \cdmin runs on one-minute data with a baseline of 4 hours before a detection window that starts 30 minutes before the query midpoint and ends one minute after it; this is the resolution for which it was written.

\section{\openrca Wrapper and Invalid Outputs}
\label{sec:openrca}

\openrca asks for a root-cause component, a reason, and an occurrence time, while the audited methods rank metrics.
For each query, the wrapper loads the metrics that cover the query window (a 4-hour baseline and 30 minutes on either side of the window midpoint) and passes them to the method.
The answer component is the component of the first ranked metric that belongs to the system's candidate component list; the reason is mapped from the type of that metric; the time is the moment at which that metric deviates most from its baseline.
The official scorer then awards the fraction of the required items that match (component and reason by exact string, time within 60\,s).

When the wrapper cannot produce an answer (no data for the window, an exception in the method, an empty ranking, or no candidate component in the ranking), it returns a fixed answer per system.
Such fixed answers can earn partial credit, because every query is scored on each required item separately.
We therefore record the wrapper status of every query and treat every fixed answer as an invalid output that scores 0 (paper \S3.3).
This matters most for \ediag on Telecom: the method returns exactly three metrics per query, none of which belonged to a candidate component in 32 of the 51 queries.
The per-subsystem coverage of every method is listed in \cref{supp:s5-coverage}.
One earlier run of \alertcount on \openrca passed a wrongly named keyword argument, so every query fell back to the fixed answer; the main analysis uses the corrected run, and the earlier run is kept only as a diagnostic in the released table.

\section{Under-Specified Choices and Silent Failures in \rcaeval}
\label{sec:defects}

This section gives the details of the under-specified input choice and the two silent failures that paper \S3.3 and \S5.2 discuss; all refer to the audited repository version (commit \texttt{497505b}).

\textbf{Input file.}
For RE1-SS and RE1-TT, each case directory contains a raw export (\texttt{data.csv}) and a reduced table (\texttt{simple\_data.csv}).
The released runner globs for \texttt{data.csv} first and falls back to \texttt{simple\_metrics.csv} (the RE2 file name); the lines that would switch to \texttt{simple\_data.csv} are commented out.
RE1-OB has only \texttt{data.csv}, with 51--60 columns.
Accuracy of every method on both inputs is in \cref{supp:s6-input}.

\textbf{Dependency version.}
\causalrca ranks the nodes of its learned graph with the PageRank implementation of scikit-network.
The \texttt{fit\_transform} method it calls was removed in scikit-network 0.32 (as stated in the maintainers' later fix); the released requirements set only a lower bound.
Under a newer version, every call raises an exception that the wrapper catches, and the wrapper returns the input column order.
In a permutation test on 30 cases (10 each from RE1-OB, RE2-OB, and RE2-SS), the ranking equalled the input column order in 30 of 30 cases under scikit-network 0.33.0, and after permuting the columns it equalled the permuted order in 30 of 30; under 0.31.0 it equalled the column order in 0 of 30.
All \causalrca results in the paper use 0.31.0.

\textbf{Duplicated header.}
Fifty of the 125 RE1-OB tables (all 25 CPU and all 25 memory faults) contain a second column named \texttt{time}, which pandas reads as \texttt{time.1}.
No other \rcaeval input file we use has a repeated column name.
On the audited release, \circa returned the input column order in 49 of the 50 affected cases and in none of the 75 unaffected cases; \ediag ranked \texttt{time.1} first in 49 of the 50.
We reran \circa, \ediag, and \rcd (three seeds) on all 125 RE1-OB cases with the duplicated column dropped when the table is read, which is the same change as the upstream fix.
When \circa was rerun on all 125 RE1-OB cases, 1 of the 75 unaffected cases gave a different top-1 than the older run, which came from a different run script and environment; we kept the older values for those 75 cases and replaced only the 50 affected ones, so the fixed condition combines two runs that differ in one case.
On the 50 affected cases, \circa's top-1 accuracy rose from 0.220 to 0.820 (32 case scores changed); \ediag's scores did not change (0.000 before and after); \rcd changed 5, 4, and 3 case scores for its three seeds.
\baro and the three probes gave identical scores with and without the fix, so their original runs are kept.
The as-shipped outputs remain in the released table under a separate role and are compared with the fixed outputs in \cref{supp:s8-release-defect}.

\section{\causalrca Runs}
\label{sec:causalrca}

\textbf{Seeds.}
The \rcaeval version of \causalrca defines a seed but never sets it, and in the original implementation the seed calls are commented out; the original repository's scripts repeat every configuration ten times.
We set the Python, NumPy, and PyTorch seeds to 0, 1, and 2.
Rerunning 18 case--seed combinations reproduced the top-20 ranking of the main run in all 18, so the runs are reproducible given the seed, and the differences between seeds come from the random initialization.
Per-seed accuracy varies widely, for example 0.160, 0.208, and 0.040 on RE1-TT; per-seed pair results are in \cref{supp:s2-seeds}.

\textbf{Empty graphs.}
The learned adjacency matrix is thresholded at 0.3.
On RE2-TT with seed 1, 43 of the 90 thresholded graphs were empty; PageRank then raises an error and the wrapper returns the column order.
We count these 43 outputs as invalid (coverage in \cref{supp:s5-coverage}).
No other dataset or seed returned the column order.

\textbf{Difference from the original implementation.}
The \rcaeval version divides every metric by its column maximum before training; the original training scripts do not normalize.
Hyperparameters, network architecture, and the thresholding steps are identical.
On RE2-TT our three-seed mean top-1 accuracy is 0.207, against 0.22 in \rcaeval's Table 6.

\section{Authors' \simplerca Branch}
\label{sec:simplerca}

The default branch of the authors' repository is not the code behind the \rcaeval results in the \simplerca paper.
Two other branches are: one for RE2 and RE3 (commit \texttt{5fd1f7f}), and one that changes the preprocessing and \baro setting for RE3-SS (commit \texttt{7f7f295}).
Both run on a benchmark platform whose data converter must produce log tables with an error attribute that the code reads; the version pinned in the authors' lock file does not, and the first platform release that does was published about seven hours after the branch commit.
We copy the algorithm files unchanged, replace only the platform input and output layer, and convert the data with that platform release.
\Cref{supp:simplerca} shows that this reproduces every value the authors report for \simplerca and their \baro baseline on RE2 and RE3 to two decimals.

The branch ranks services by \baro's metric scores (the maximum per service), and then adds points from trace and log evidence to the top three services.
This differs from the procedure described in the paper's text, which counts threshold alerts per service.
RE1 contains only metrics, so on RE1 the branch adds no trace or log evidence; its RE1 results (bottom of \cref{supp:simplerca}) are therefore reported here and not as a separate method in the main analysis.
They are close to those of the authors' own \baro adapter on RE1-OB and RE1-SS and higher on RE1-TT.
The RE1-OB inputs of these runs include the 50 tables with the duplicated header; we did not rerun them, because they enter no comparison.

The branch's code differs from the method description in the \simplerca paper's Section 3.1.2 in details of the scoring step; we use the code as released because it reproduces the authors' reported values, and we do not attempt to reconcile the text.

\begin{table}[h]
\caption{Authors' \simplerca branch. Top: our runs on RE2 and RE3 with the values reported in the \simplerca paper in parentheses (Top@1, Top@3, Top@5, Avg@3, Avg@5, MRR with misses counted as 0); the last column is the largest absolute difference after rounding to two decimals. RE3-SS uses the second branch the authors provide. Bottom: the same code on RE1 (20-minute window; OB \texttt{data.csv}, SS/TT \texttt{simple\_data.csv}), scored by exact service-name match; RE1 contains only metrics, so no trace or log evidence is added. Source: \texttt{simplerca\_summary.txt}.}
\label{supp:simplerca}
\footnotesize
\setlength{\tabcolsep}{3pt}
\begin{tabular}{@{}llrrrrrrr@{}}
\toprule
Data & Method & Top@1 & Top@3 & Top@5 & Avg@3 & Avg@5 & MRR & Max diff. \\
\midrule
RE2-OB & \baro (authors') & 0.14 (0.14) & 0.93 (0.93) & 0.98 (0.98) & 0.64 (0.64) & 0.77 (0.77) & 0.54 (0.54) & 0.00 \\
RE2-OB & \simplerca & 0.47 (0.47) & 0.93 (0.93) & 1.00 (1.00) & 0.68 (0.68) & 0.80 (0.80) & 0.67 (0.67) & 0.00 \\
RE2-SS & \baro (authors') & 0.14 (0.14) & 0.82 (0.82) & 0.94 (0.94) & 0.50 (0.50) & 0.67 (0.67) & 0.47 (0.47) & 0.00 \\
RE2-SS & \simplerca & 0.24 (0.24) & 0.88 (0.88) & 0.97 (0.97) & 0.59 (0.59) & 0.74 (0.74) & 0.55 (0.55) & 0.00 \\
RE2-TT & \baro (authors') & 0.67 (0.67) & 0.84 (0.84) & 0.89 (0.89) & 0.77 (0.77) & 0.81 (0.81) & 0.76 (0.76) & 0.00 \\
RE2-TT & \simplerca & 0.80 (0.80) & 0.87 (0.87) & 0.93 (0.93) & 0.84 (0.84) & 0.87 (0.87) & 0.85 (0.85) & 0.00 \\
RE3-OB & \baro (authors') & 0.00 (0.00) & 0.90 (0.90) & 0.90 (0.90) & 0.59 (0.59) & 0.71 (0.71) & 0.45 (0.45) & 0.00 \\
RE3-OB & \simplerca & 0.30 (0.30) & 0.90 (0.90) & 0.90 (0.90) & 0.62 (0.62) & 0.73 (0.73) & 0.56 (0.56) & 0.00 \\
RE3-SS & \baro (authors') & 0.00 (0.00) & 0.83 (0.83) & 0.93 (0.93) & 0.46 (0.46) & 0.65 (0.65) & 0.40 (0.40) & 0.00 \\
RE3-SS & \simplerca & 0.83 (0.83) & 0.90 (0.90) & 0.90 (0.90) & 0.88 (0.88) & 0.89 (0.89) & 0.87 (0.87) & 0.00 \\
RE3-TT & \baro (authors') & 0.50 (0.50) & 0.93 (0.93) & 1.00 (1.00) & 0.77 (0.77) & 0.86 (0.86) & 0.72 (0.72) & 0.00 \\
RE3-TT & \simplerca & 0.83 (0.83) & 1.00 (1.00) & 1.00 (1.00) & 0.93 (0.93) & 0.96 (0.96) & 0.91 (0.91) & 0.00 \\
\midrule
RE1-OB & \baro (authors') & 0.776 & 0.928 & 0.968 & 0.867 & 0.904 & 0.860 & -- \\
RE1-OB & \simplerca & 0.784 & 0.936 & 0.976 & 0.875 & 0.912 & 0.864 & -- \\
RE1-SS & \baro (authors') & 0.368 & 0.920 & 0.992 & 0.680 & 0.800 & 0.634 & -- \\
RE1-SS & \simplerca & 0.368 & 0.928 & 0.992 & 0.683 & 0.803 & 0.634 & -- \\
RE1-TT & \baro (authors') & 0.464 & 0.672 & 0.728 & 0.581 & 0.634 & 0.581 & -- \\
RE1-TT & \simplerca & 0.568 & 0.856 & 0.952 & 0.733 & 0.813 & 0.721 & -- \\
\bottomrule
\end{tabular}
\end{table}


\section{Comparison with \rcaeval's Published RE2 Table}
\label{sec:table6}

\rcaeval's paper reports per-fault accuracy only for RE2-TT (its Table 6).
\Cref{supp:table6} lists our runs next to it.
For \tracerca and \microrank, the 18 per-fault cells computed over valid outputs only are within 0.01 of Table 6, while the same accuracy over all cases is lower by up to 0.41 on the memory and socket faults: both methods raise an error on the same 17 of the 90 cases, because a service--operation pair that appears only after the injection has no baseline entry.
The multi-source variant of \baro uses its trace input only if the dataset name passed to it is \texttt{mm-tt}; with the name that the released runner passes (\texttt{re2-tt}), its AC@1 is almost unchanged (0.68 against 0.69), and with the raw metric export instead of the reduced one it drops to 0.14.
\circa, run with the configuration shipped in the repository, has a higher AC@1 than Table 6 on five of the six faults and the same on the sixth; we could not identify the configuration behind the published values.

\begin{longtable}{@{}llll@{}}
\caption{RE2-TT per fault type: our runs over all cases (invalid outputs score 0), over valid outputs only, and the values in \rcaeval's Table 6. Cells are AC@1 / AC@3 / Avg@5. Source: \texttt{re2\_table6\_compare.txt}.}\label{supp:table6}\\
\toprule
Fault & All cases & Valid outputs only & Table 6 \\
\endfirsthead
\toprule
Fault & All cases & Valid outputs only & Table 6 \\
\midrule
\endhead
\bottomrule
\endfoot
\midrule
\multicolumn{4}{@{}l}{\circa: 90/90 valid outputs} \\
cpu & 0.40 / 0.60 / 0.53 ($n$=15) & 0.40 / 0.60 / 0.53 ($n$=15) & 0.27 / 0.27 / 0.28 \\
mem & 0.73 / 0.93 / 0.88 ($n$=15) & 0.73 / 0.93 / 0.88 ($n$=15) & 0.47 / 0.73 / 0.68 \\
disk & 0.80 / 0.80 / 0.80 ($n$=15) & 0.80 / 0.80 / 0.80 ($n$=15) & 0.53 / 0.67 / 0.64 \\
socket & 0.80 / 0.87 / 0.84 ($n$=15) & 0.80 / 0.87 / 0.84 ($n$=15) & 0.27 / 0.53 / 0.52 \\
delay & 0.20 / 0.40 / 0.36 ($n$=15) & 0.20 / 0.40 / 0.36 ($n$=15) & 0.2 / 0.27 / 0.28 \\
loss & 0.33 / 0.60 / 0.56 ($n$=15) & 0.33 / 0.60 / 0.56 ($n$=15) & 0.2 / 0.33 / 0.35 \\
all & 0.54 / 0.70 / 0.66 ($n$=90) & 0.54 / 0.70 / 0.66 ($n$=90) & 0.32 / 0.47 / 0.46 \\
\midrule
\multicolumn{4}{@{}l}{\microrank: 73/90 valid outputs} \\
cpu & 0.20 / 0.40 / 0.32 ($n$=15) & 0.21 / 0.43 / 0.34 ($n$=14) & 0.21 / 0.43 / 0.34 \\
mem & 0.13 / 0.20 / 0.17 ($n$=15) & 0.25 / 0.38 / 0.33 ($n$=8) & 0.25 / 0.38 / 0.33 \\
disk & 0.00 / 0.33 / 0.25 ($n$=15) & 0.00 / 0.36 / 0.27 ($n$=14) & 0 / 0.36 / 0.27 \\
socket & 0.20 / 0.27 / 0.24 ($n$=15) & 0.30 / 0.40 / 0.36 ($n$=10) & 0.3 / 0.4 / 0.36 \\
delay & 0.07 / 0.27 / 0.20 ($n$=15) & 0.08 / 0.31 / 0.23 ($n$=13) & 0.08 / 0.31 / 0.23 \\
loss & 0.13 / 0.33 / 0.28 ($n$=15) & 0.14 / 0.36 / 0.30 ($n$=14) & 0.14 / 0.36 / 0.3 \\
all & 0.12 / 0.30 / 0.24 ($n$=90) & 0.15 / 0.37 / 0.30 ($n$=73) & 0.16 / 0.37 / 0.31 \\
\midrule
\multicolumn{4}{@{}l}{\baro (multi-source), name \texttt{mm-tt}: 89/90 valid outputs} \\
cpu & 0.47 / 0.80 / 0.75 ($n$=15) & 0.47 / 0.80 / 0.75 ($n$=15) & 0.47 / 0.8 / 0.75 \\
mem & 0.93 / 1.00 / 0.99 ($n$=15) & 0.93 / 1.00 / 0.99 ($n$=15) & 0.93 / 1 / 0.99 \\
disk & 1.00 / 1.00 / 1.00 ($n$=15) & 1.00 / 1.00 / 1.00 ($n$=15) & 1 / 1 / 1 \\
socket & 0.60 / 0.73 / 0.73 ($n$=15) & 0.64 / 0.79 / 0.79 ($n$=14) & 0.6 / 0.8 / 0.79 \\
delay & 0.47 / 0.67 / 0.61 ($n$=15) & 0.47 / 0.67 / 0.61 ($n$=15) & 0.47 / 0.67 / 0.61 \\
loss & 0.67 / 0.67 / 0.71 ($n$=15) & 0.67 / 0.67 / 0.71 ($n$=15) & 0.67 / 0.67 / 0.71 \\
all & 0.69 / 0.81 / 0.80 ($n$=90) & 0.70 / 0.82 / 0.81 ($n$=89) & 0.69 / 0.82 / 0.81 \\
\midrule
\multicolumn{4}{@{}l}{\baro (multi-source), \texttt{metrics.csv}: 89/90 valid outputs} \\
cpu & 0.00 / 0.00 / 0.00 ($n$=15) & 0.00 / 0.00 / 0.00 ($n$=15) & 0.47 / 0.8 / 0.75 \\
mem & 0.00 / 0.07 / 0.16 ($n$=15) & 0.00 / 0.07 / 0.16 ($n$=15) & 0.93 / 1 / 0.99 \\
disk & 0.80 / 1.00 / 0.96 ($n$=15) & 0.80 / 1.00 / 0.96 ($n$=15) & 1 / 1 / 1 \\
socket & 0.00 / 0.00 / 0.00 ($n$=15) & 0.00 / 0.00 / 0.00 ($n$=14) & 0.6 / 0.8 / 0.79 \\
delay & 0.07 / 0.13 / 0.17 ($n$=15) & 0.07 / 0.13 / 0.17 ($n$=15) & 0.47 / 0.67 / 0.61 \\
loss & 0.00 / 0.07 / 0.04 ($n$=15) & 0.00 / 0.07 / 0.04 ($n$=15) & 0.67 / 0.67 / 0.71 \\
all & 0.14 / 0.21 / 0.22 ($n$=90) & 0.15 / 0.21 / 0.22 ($n$=89) & 0.69 / 0.82 / 0.81 \\
\midrule
\multicolumn{4}{@{}l}{\baro (multi-source), name \texttt{re2-tt}: 89/90 valid outputs} \\
cpu & 0.47 / 0.80 / 0.75 ($n$=15) & 0.47 / 0.80 / 0.75 ($n$=15) & 0.47 / 0.8 / 0.75 \\
mem & 0.93 / 1.00 / 0.99 ($n$=15) & 0.93 / 1.00 / 0.99 ($n$=15) & 0.93 / 1 / 0.99 \\
disk & 1.00 / 1.00 / 1.00 ($n$=15) & 1.00 / 1.00 / 1.00 ($n$=15) & 1 / 1 / 1 \\
socket & 0.53 / 0.73 / 0.72 ($n$=15) & 0.57 / 0.79 / 0.77 ($n$=14) & 0.6 / 0.8 / 0.79 \\
delay & 0.47 / 0.67 / 0.61 ($n$=15) & 0.47 / 0.67 / 0.61 ($n$=15) & 0.47 / 0.67 / 0.61 \\
loss & 0.67 / 0.67 / 0.71 ($n$=15) & 0.67 / 0.67 / 0.71 ($n$=15) & 0.67 / 0.67 / 0.71 \\
all & 0.68 / 0.81 / 0.80 ($n$=90) & 0.69 / 0.82 / 0.80 ($n$=89) & 0.69 / 0.82 / 0.81 \\
\midrule
\multicolumn{4}{@{}l}{\tracerca: 73/90 valid outputs} \\
cpu & 0.60 / 0.73 / 0.69 ($n$=15) & 0.64 / 0.79 / 0.74 ($n$=14) & 0.64 / 0.79 / 0.74 \\
mem & 0.33 / 0.47 / 0.44 ($n$=15) & 0.62 / 0.88 / 0.82 ($n$=8) & 0.63 / 0.88 / 0.83 \\
disk & 0.60 / 0.67 / 0.69 ($n$=15) & 0.64 / 0.71 / 0.74 ($n$=14) & 0.64 / 0.71 / 0.74 \\
socket & 0.40 / 0.53 / 0.51 ($n$=15) & 0.60 / 0.80 / 0.76 ($n$=10) & 0.6 / 0.8 / 0.76 \\
delay & 0.73 / 0.73 / 0.76 ($n$=15) & 0.85 / 0.85 / 0.88 ($n$=13) & 0.85 / 0.85 / 0.88 \\
loss & 0.53 / 0.67 / 0.63 ($n$=15) & 0.57 / 0.71 / 0.67 ($n$=14) & 0.57 / 0.71 / 0.67 \\
all & 0.53 / 0.63 / 0.62 ($n$=90) & 0.66 / 0.78 / 0.76 ($n$=73) & 0.66 / 0.79 / 0.77 \\
\end{longtable}


\section{Window Length on RE1}
\label{sec:window}

All RE1 and RE2 runs in the paper use \rcaeval's default window of 20 minutes (10 minutes before and after the injection).
\Cref{supp:window} compares it with a 10-minute window for \baro and the three probes.
The accuracy changes by up to 0.208 (\maxz on RE1-OB); we report this as a diagnostic and do not use the 10-minute runs in any comparison.

\begin{table}[h]
\caption{RE1 top-1 accuracy with a 10- and a 20-minute window (125 cases each). Lost / gained: cases correct at 10 minutes and wrong at 20, and the reverse. Source: \texttt{re1\_window\_compare.txt}.}
\label{supp:window}
\footnotesize
\begin{tabular}{@{}llrrrrr@{}}
\toprule
Method & Subsystem & 10 min & 20 min & Lost & Gained & Valid (10/20) \\
\midrule
\baro & RE1-OB & 0.728 & 0.736 & 3 & 4 & 125/125 \\
\baro & RE1-SS & 0.832 & 0.848 & 1 & 3 & 125/125 \\
\baro & RE1-TT & 0.592 & 0.560 & 9 & 5 & 125/125 \\
\maxz & RE1-OB & 0.712 & 0.504 & 32 & 6 & 125/125 \\
\maxz & RE1-SS & 0.840 & 0.744 & 16 & 4 & 125/125 \\
\maxz & RE1-TT & 0.448 & 0.424 & 9 & 6 & 125/125 \\
\alertcount & RE1-OB & 0.432 & 0.392 & 7 & 2 & 125/125 \\
\alertcount & RE1-SS & 0.320 & 0.376 & 3 & 10 & 125/125 \\
\alertcount & RE1-TT & 0.072 & 0.048 & 5 & 2 & 125/125 \\
\cdmin & RE1-OB & 0.792 & 0.728 & 8 & 0 & 125/125 \\
\cdmin & RE1-SS & 0.944 & 0.928 & 3 & 1 & 125/125 \\
\cdmin & RE1-TT & 0.304 & 0.280 & 4 & 1 & 125/125 \\
\bottomrule
\end{tabular}
\end{table}


\section{Margin and Spread}
\label{sec:margin}

Paper \S4.1 describes every pooled comparison by its margin $|\bar\Delta|$ and the spread of its subsystem differences, and gives the counts for our audit; this section gives the counts for the two published tables separately.
Among published pairs, \NPubBelow have a margin smaller than the spread; \NPubBelowRev of them reverse, \NPubBelowRevCI with CI support, whereas of the \NPubAbove whose margin exceeds the spread \NPubAboveRev reverses, without CI support (with probes and shipped baselines: \NAllBelowRev of \NAllBelow, and \NAllAboveRev of \NAllAbove).
In \rcaeval's RE2-TT accuracy by fault type (RCAEval paper, Table 6), \NTabRcaevalBelowRev of the \NTabRcaevalBelow pairs with margin below spread reverse on some fault type, against \NTabRcaevalAboveRev of \NTabRcaevalAbove above it; in \petshop's accuracy by scenario and fault type (PetShop paper, Table 8), \NTabPetshopBelowRev of \NTabPetshopBelow and \NTabPetshopAboveRev of \NTabPetshopAbove.
Margin and spread computed on the subsystems other than a held-out one do not single out the wrong choices of the held-out analysis (\cref{supp:heldout-sec}, $c = 1$).

\section{Held-Out Subsystems}
\label{supp:heldout-sec}
For each family, each subsystem is held out in turn; methods present on every subsystem of the family are ranked by their mean over the remaining subsystems (subsystems weighted equally, ties by name), and every adjacent pair of that ranking is one decision. Margin and spread are computed from the remaining subsystems only, so the spread rests on two differences for RE1 and three for \openrca and \petshop. \Cref{supp:heldout} gives the counts for several multiples $c$ of the spread, by layer and family.

\begin{table}[h]
\centering
\caption{Held-out subsystems: adjacent pairs of the leaderboard built on the other subsystems of a family. A decision is passed when the margin on the remaining subsystems is at least $c$ times their spread ($c = 1$ is the line of the main text; $c = 0$ passes every decision with a direction). Wrong: the higher-ranked method is less accurate on the held-out subsystem. The last two columns pass the same number of decisions by the largest system-equal and case-weighted pooled margins.}\label{supp:heldout}
\scriptsize
\begin{tabular}{@{}llrrrrrrrr@{}}
\toprule
Layer & Family & $c$ & Decisions & Passed & Wrong & Flagged & Wrong & \shortstack{Margin only\\(equal) wrong} & \shortstack{Margin only\\(case-wtd.) wrong} \\
\midrule
Published & All three & 0.0 & 36 & 36 & 8 & 0 & 0 & 8 & 10 \\
 &  & 0.5 & 36 & 26 & 6 & 10 & 2 & 5 & 7 \\
 &  & 1.0 & 36 & 19 & 4 & 17 & 4 & 3 & 3 \\
 &  & 1.5 & 36 & 10 & 2 & 26 & 6 & 2 & 2 \\
 &  & 2.0 & 36 & 7 & 1 & 29 & 7 & 2 & 1 \\
\addlinespace
 & \rcaeval RE1 & 0.0 & 12 & 12 & 3 & 0 & 0 & 3 & 3 \\
 &  & 0.5 & 12 & 10 & 1 & 2 & 2 & 1 & 1 \\
 &  & 1.0 & 12 & 9 & 1 & 3 & 2 & 0 & 0 \\
 &  & 1.5 & 12 & 5 & 1 & 7 & 2 & 0 & 0 \\
 &  & 2.0 & 12 & 3 & 1 & 9 & 2 & 0 & 0 \\
\addlinespace
 & \openrca & 0.0 & 8 & 8 & 1 & 0 & 0 & 1 & 1 \\
 &  & 0.5 & 8 & 7 & 1 & 1 & 0 & 1 & 1 \\
 &  & 1.0 & 8 & 4 & 1 & 4 & 0 & 1 & 1 \\
 &  & 1.5 & 8 & 0 & 0 & 8 & 1 & 0 & 0 \\
 &  & 2.0 & 8 & 0 & 0 & 8 & 1 & 0 & 0 \\
\addlinespace
 & \petshop & 0.0 & 16 & 16 & 4 & 0 & 0 & 4 & 6 \\
 &  & 0.5 & 16 & 9 & 4 & 7 & 0 & 2 & 2 \\
 &  & 1.0 & 16 & 6 & 2 & 10 & 2 & 2 & 2 \\
 &  & 1.5 & 16 & 5 & 1 & 11 & 3 & 2 & 2 \\
 &  & 2.0 & 16 & 4 & 0 & 12 & 4 & 2 & 2 \\
\addlinespace
All pairs & All three & 0.0 & 81 & 79 & 27 & 2 & 0 & 27 & 27 \\
 &  & 0.5 & 81 & 44 & 13 & 37 & 14 & 14 & 15 \\
 &  & 1.0 & 81 & 20 & 6 & 61 & 21 & 5 & 4 \\
 &  & 1.5 & 81 & 9 & 3 & 72 & 24 & 2 & 2 \\
 &  & 2.0 & 81 & 6 & 1 & 75 & 26 & 1 & 1 \\
\addlinespace
 & \rcaeval RE1 & 0.0 & 21 & 21 & 8 & 0 & 0 & 8 & 8 \\
 &  & 0.5 & 21 & 12 & 4 & 9 & 4 & 3 & 3 \\
 &  & 1.0 & 21 & 8 & 2 & 13 & 6 & 1 & 1 \\
 &  & 1.5 & 21 & 5 & 2 & 16 & 6 & 1 & 1 \\
 &  & 2.0 & 21 & 3 & 1 & 18 & 7 & 0 & 0 \\
\addlinespace
 & \openrca & 0.0 & 20 & 20 & 6 & 0 & 0 & 6 & 5 \\
 &  & 0.5 & 20 & 10 & 2 & 10 & 4 & 1 & 1 \\
 &  & 1.0 & 20 & 4 & 1 & 16 & 5 & 0 & 0 \\
 &  & 1.5 & 20 & 0 & 0 & 20 & 6 & 0 & 0 \\
 &  & 2.0 & 20 & 0 & 0 & 20 & 6 & 0 & 0 \\
\addlinespace
 & \petshop & 0.0 & 40 & 38 & 13 & 2 & 0 & 13 & 14 \\
 &  & 0.5 & 40 & 22 & 7 & 18 & 6 & 6 & 7 \\
 &  & 1.0 & 40 & 8 & 3 & 32 & 10 & 3 & 3 \\
 &  & 1.5 & 40 & 4 & 1 & 36 & 12 & 2 & 1 \\
 &  & 2.0 & 40 & 3 & 0 & 37 & 13 & 1 & 1 \\
\bottomrule
\end{tabular}
\end{table}


\section{Cross-Family Comparison Set}
\label{supp:cross-sec}
Six methods have outputs on all 11 subsystems of the main analysis: \baro, \rcd, \ediag, and the three probes (on \petshop, \rcd and \ediag are the implementations \petshop ships). \Cref{supp:cross-signs} gives, for each of the 15 pairs, the subsystems each method wins; \cref{supp:cross-loso} lists the held-out subsystems on which the pooled winner of the other ten is less accurate; \cref{supp:cross-robust} gives the robustness checks summarized in paper \S4.5.

\begin{table}[h]
\centering
\caption{Cross-family set: the six methods with outputs on all 11 subsystems. Scores are not averaged across families (\openrca partial credit, top-1 accuracy elsewhere). $+$/$-$: subsystems where $A$ / $B$ is more accurate; CI excl.\ 0: of these, subsystems whose 95\% CI excludes zero. Both signs: the number of leave-one-family-out subsets (of three) and leave-two-subsystems-out subsets (of \XLksoPer) on which the pair still has subsystem differences of both signs.}
\label{supp:cross-signs}
\small
\setlength{\tabcolsep}{5pt}
\begin{tabular}{@{}lcccc@{}}
\toprule
Pair $A$ vs $B$ & \shortstack{Subsystems\\$+$ / $-$} & \shortstack{CI excl.\ 0\\$+$ / $-$} & \shortstack{Both signs,\\family dropped} & \shortstack{Both signs,\\two subsystems dropped} \\
\midrule
Alert-count vs BARO & 4 / 6 & 3 / 4 & 2 of 3 & 55 of 55 \\
Alert-count vs CD-1min & 5 / 6 & 0 / 3 & 3 of 3 & 55 of 55 \\
Alert-count vs $\epsilon$-Diagnosis & 9 / 0 & 7 / 0 & 0 of 3 & 0 of 55 \\
Alert-count vs max-$|Z|$ & 2 / 9 & 0 / 4 & 2 of 3 & 54 of 55 \\
Alert-count vs RCD & 5 / 6 & 3 / 3 & 3 of 3 & 55 of 55 \\
BARO vs CD-1min & 7 / 4 & 2 / 3 & 3 of 3 & 55 of 55 \\
BARO vs $\epsilon$-Diagnosis & 11 / 0 & 8 / 0 & 0 of 3 & 0 of 55 \\
BARO vs max-$|Z|$ & 5 / 6 & 3 / 2 & 3 of 3 & 55 of 55 \\
BARO vs RCD & 7 / 4 & 5 / 0 & 3 of 3 & 55 of 55 \\
CD-1min vs $\epsilon$-Diagnosis & 9 / 0 & 7 / 0 & 0 of 3 & 0 of 55 \\
CD-1min vs max-$|Z|$ & 4 / 7 & 2 / 2 & 3 of 3 & 55 of 55 \\
CD-1min vs RCD & 6 / 4 & 6 / 2 & 3 of 3 & 55 of 55 \\
$\epsilon$-Diagnosis vs max-$|Z|$ & 0 / 11 & 0 / 9 & 0 of 3 & 0 of 55 \\
$\epsilon$-Diagnosis vs RCD & 0 / 10 & 0 / 7 & 0 of 3 & 0 of 55 \\
max-$|Z|$ vs RCD & 7 / 3 & 4 / 0 & 3 of 3 & 55 of 55 \\
\bottomrule
\end{tabular}

\end{table}

\begin{table}[h]
\centering
\caption{Cross-family set, held-out selection: for each of the 11 subsystems, the pooled winner of the other ten (subsystems weighted equally) is applied to it. Regret: accuracy lost on the held-out subsystem.}\label{supp:cross-loso}
\scriptsize
\begin{tabular}{@{}lrrrp{6.2cm}@{}}
\toprule
Pair & Reversals & Mean regret & Max regret & Held-out subsystems where the pooled winner loses \\
\midrule
\alertcount vs \baro & 4 of 11 & 0.023 & 0.091 & Bank, Market-1, Market-2, Telecom \\
\alertcount vs \cdmin & 5 of 11 & 0.023 & 0.125 & Bank, Market-2, Telecom, High, Temporal-2 \\
\alertcount vs \ediag & 0 of 11 & 0.000 & 0.000 &  \\
\alertcount vs \maxz & 2 of 11 & 0.005 & 0.045 & Bank, Telecom \\
\alertcount vs \rcd & 5 of 11 & 0.036 & 0.130 & Bank, Market-1, Market-2, Telecom, High \\
\baro vs \cdmin & 4 of 11 & 0.025 & 0.104 & RE1-SS, Bank, Market-1, Market-2 \\
\baro vs \ediag & 0 of 11 & 0.000 & 0.000 &  \\
\baro vs \maxz & 10 of 11 & 0.099 & 0.232 & RE1-OB, RE1-SS, RE1-TT, Market-1, Market-2, Telecom, High, Low, Temporal-1, Temporal-2 \\
\baro vs \rcd & 4 of 11 & 0.064 & 0.375 & Telecom, Low, Temporal-1, Temporal-2 \\
\cdmin vs \ediag & 0 of 11 & 0.000 & 0.000 &  \\
\cdmin vs \maxz & 4 of 11 & 0.041 & 0.224 & RE1-OB, RE1-SS, Bank, Market-1 \\
\cdmin vs \rcd & 7 of 11 & 0.182 & 0.539 & RE1-OB, RE1-SS, RE1-TT, Telecom, Low, Temporal-1, Temporal-2 \\
\ediag vs \maxz & 0 of 11 & 0.000 & 0.000 &  \\
\ediag vs \rcd & 0 of 11 & 0.000 & 0.000 &  \\
\maxz vs \rcd & 3 of 11 & 0.040 & 0.250 & RE1-OB, Low, Temporal-1 \\
\bottomrule
\end{tabular}
\end{table}

\begin{table}[h]
\centering
\caption{Cross-family set, robustness. Drop one family: subsystems where $A$ / $B$ is more accurate among the remaining ones. Drop two: cases (out of all 55 ways to drop two subsystems) in which both signs remain. Pooled difference under four pooling rules (trimmed: the largest and smallest subsystem differences removed). Bootstrap: share of 2{,}000 case resamples in which both signs remain.}\label{supp:cross-robust}
\scriptsize
\setlength{\tabcolsep}{3pt}
\begin{tabular}{@{}lp{4.6cm}rrrrrr@{}}
\toprule
Pair & Drop one family & Drop two & Mean & Median & Case-wtd. & Trimmed & Bootstrap \\
\midrule
\alertcount vs \baro & \rcaeval RE1: 4/3; \openrca: 0/6; \petshop: 4/3 & 55/55 & $-$0.145 & $-$0.115 & $-$0.202 & $-$0.131 & 1.000 \\
\alertcount vs \cdmin & \rcaeval RE1: 5/3; \openrca: 2/5; \petshop: 3/4 & 55/55 & $-$0.095 & $-$0.017 & $-$0.175 & $-$0.069 & 0.999 \\
\alertcount vs \ediag & \rcaeval RE1: 6/0; \openrca: 5/0; \petshop: 7/0 & 0/55 & $+$0.127 & $+$0.125 & $+$0.167 & $+$0.115 & 0.185 \\
\alertcount vs \maxz & \rcaeval RE1: 2/6; \openrca: 0/7; \petshop: 2/5 & 54/55 & $-$0.141 & $-$0.112 & $-$0.144 & $-$0.136 & 0.997 \\
\alertcount vs \rcd & \rcaeval RE1: 5/3; \openrca: 1/6; \petshop: 4/3 & 55/55 & $-$0.082 & $-$0.013 & $-$0.011 & $-$0.046 & 1.000 \\
\baro vs \cdmin & \rcaeval RE1: 5/3; \openrca: 6/1; \petshop: 3/4 & 55/55 & $+$0.050 & $+$0.020 & $+$0.027 & $+$0.042 & 1.000 \\
\baro vs \ediag & \rcaeval RE1: 8/0; \openrca: 7/0; \petshop: 7/0 & 0/55 & $+$0.272 & $+$0.154 & $+$0.369 & $+$0.252 & 0.340 \\
\baro vs \maxz & \rcaeval RE1: 2/6; \openrca: 4/3; \petshop: 4/3 & 55/55 & $+$0.004 & $-$0.038 & $+$0.058 & $-$0.007 & 1.000 \\
\baro vs \rcd & \rcaeval RE1: 4/4; \openrca: 4/3; \petshop: 6/1 & 55/55 & $+$0.063 & $+$0.040 & $+$0.191 & $+$0.066 & 1.000 \\
\cdmin vs \ediag & \rcaeval RE1: 6/0; \openrca: 5/0; \petshop: 7/0 & 0/55 & $+$0.222 & $+$0.125 & $+$0.342 & $+$0.184 & 0.361 \\
\cdmin vs \maxz & \rcaeval RE1: 2/6; \openrca: 2/5; \petshop: 4/3 & 55/55 & $-$0.046 & $-$0.069 & $+$0.031 & $-$0.053 & 1.000 \\
\cdmin vs \rcd & \rcaeval RE1: 3/4; \openrca: 3/3; \petshop: 6/1 & 55/55 & $+$0.013 & $+$0.076 & $+$0.163 & $+$0.011 & 1.000 \\
\ediag vs \maxz & \rcaeval RE1: 0/8; \openrca: 0/7; \petshop: 0/7 & 0/55 & $-$0.268 & $-$0.192 & $-$0.311 & $-$0.252 & 0.086 \\
\ediag vs \rcd & \rcaeval RE1: 0/7; \openrca: 0/6; \petshop: 0/7 & 0/55 & $-$0.209 & $-$0.147 & $-$0.178 & $-$0.186 & 0.469 \\
\maxz vs \rcd & \rcaeval RE1: 5/2; \openrca: 3/3; \petshop: 6/1 & 55/55 & $+$0.059 & $+$0.039 & $+$0.132 & $+$0.060 & 0.997 \\
\bottomrule
\end{tabular}
\end{table}


\section{Survey of Reporting Practice}
\label{supp:survey-sec}
\textbf{Selection.} We listed 26 microservice RCA method and benchmark papers from 2021 to 2026: the papers of the audited benchmarks and methods, and papers found by keyword search in software-engineering, systems, and machine-learning venues, with at most five papers per research group. This is a convenience sample, not a systematic review. Full text was obtained for \SurveyN papers; the other three (MicroRank, Nezha, and MRCA) are listed in \cref{supp:s11-survey} but excluded from every count in the paper and in this document.

\textbf{Per-system reporting.} For each of the \SurveyN papers we recorded whether it evaluates more than one system and, if so, whether it reports a score for each system. \SurveyMulti papers evaluate more than one system or dataset (simulation datasets of different sizes count as datasets), and all \SurveyMulti report a score per system; \SurveySingle evaluate one system, among them \petshop, whose three traffic scenarios are one application. \Cref{supp:s11-survey} gives the result and the table, figure, or passage checked for each paper.

\textbf{Release check.} On 2 and 3 October 2026 we listed the files of every repository and archive linked from each paper, at the current commit, and recorded whether the release contains per-case predictions or scores of any method, which methods and cases they cover, and whether they suffice to recompute every pairwise comparison of the paper's main results. \SurveyRelAny papers release per-case outputs of at least one method. For \SurveyRelDemo (\baro, \rcaeval, and the study howfar2024, which shares \rcaeval's repository) the only such outputs are the ones a tutorial notebook saved for a single case, and they appear in the current repository, not necessarily in the one at publication. \SurveyRelOwn release the outputs of the authors' own method on their test cases: \openrca's repository holds one archived run of its agent on all 335 queries, whose files record neither the model nor the setting; CausalRCA's repository holds the rankings of the authors' method as text files per fault type and setting; and the archive behind CHASE's link, which we downloaded, holds the predictions of the authors' method on the 939 test cases of its split. \circa's README links an archive with the rankings of all eleven methods of its simulation study on every case of every run, from which its Table 1 can be recomputed, but only summary files for its production dataset. None of these releases includes outputs of the other methods of the paper's main comparison on every dataset. For \SurveyRelUndet papers a linked release had expired or could not be reached (TraceRCA's data archive and the anonymous review links of one paper). For six papers the full text gives no link to a release of the authors' own, which we record as no release linked rather than as no outputs; for the remaining papers we examined the linked repositories and, where linked, the data archives, without an exhaustive scan of nested archives. \Cref{supp:release} lists the result per paper; the notes in the package give every link and the files found.

\textbf{LLM use.} LLM sessions located the tables and passages behind the per-system check and listed the repository contents behind the release check; every statement above was verified by an author against the paper or the repository (paper \S3.5).

\begin{longtable}{lllp{7.6cm}}
\caption{Released per-case outputs in the 23 surveyed papers with full text, checked on 2 October 2026 (UTC). Any: the release contains per-case predictions or scores of any method. All: it contains them for every method of the paper's main comparison. undet.: the linked release had expired or could not be listed. Release: first repository checked, at the commit checked (--: no repository linked, or the link is not a listable repository); the notes in the package give every link and the files found.}\label{supp:release}\\
\toprule
Paper & Any & All & Release \\
\midrule
\endhead
baro2024 & no & no & \texttt{\footnotesize phamquiluan/baro@e35f4ec} \\
circa2022 & no & no & \texttt{\footnotesize NetManAIOps/CIRCA@0215e18} \\
rcd2022 & no & no & \texttt{\footnotesize azamikram/rcd@373882c} \\
tracerca2021 & no & no & \texttt{\footnotesize NetManAIOps/TraceRCA@b9087c7} \\
rcaeval2025 & no & no & \texttt{\footnotesize phamquiluan/RCAEval@259ea41} \\
fang2026 & no & no & \texttt{\footnotesize LGU-SE-Internal/Simple-RCA@1be3911} \\
openrca2025 & yes & no & \texttt{\footnotesize microsoft/OpenRCA@c1bd4af} \\
petshop2024 & no & no & \texttt{\footnotesize amazon-science/petshop-root-cause-analysis@2e96f93} \\
eadro2023 & no & no & \texttt{\footnotesize BEbillionaireUSD/Eadro@82ff9c9} \\
mabc2024 & no & no & \texttt{\footnotesize zwpride/mABC@79272ce} \\
mulan2024 & no & no & \texttt{\footnotesize --} \\
ocean2024 & no & no & \texttt{\footnotesize --} \\
tvdiag2024 & no & no & \texttt{\footnotesize WHU-AISE/TVDiag@cf6f334} \\
chase2024 & undet. & undet. & \texttt{\footnotesize --} \\
howfar2024 & no & no & \texttt{\footnotesize phamquiluan/RCAEval@259ea41} \\
lemmarca2024 & no & no & \texttt{\footnotesize zhengzhangchen/LEMMA\_rca\_baselines@9b56afb} \\
sparserca2024 & no & no & \texttt{\footnotesize --} \\
tracecontrast2024 & no & no & \texttt{\footnotesize --} \\
tracediag2023 & no & no & \texttt{\footnotesize --} \\
run2024 & no & no & \texttt{\footnotesize zmlin1998/RUN@e2c97d7} \\
dynacausal2025 & no & no & \texttt{\footnotesize --} \\
causalrca2023 & yes & no & \texttt{\footnotesize AXinx/CausalRCA\_code@21f4eb0} \\
toomanycooks2025 & undet. & undet. & \texttt{\footnotesize --} \\
\bottomrule
\end{longtable}


\section{Complete Result Tables}
\label{sec:full}

% Auto-generated by paper/scripts/make_supp_tables.py. Do not edit.
\subsection*{All pairs, main analysis}
Every seed-averaged method pair, with both pooled differences, reversal counts, dispersion, and the Holm-adjusted likelihood-ratio test. Source: \texttt{paper/data/rq1/pairs.csv}.

{\scriptsize\setlength{\tabcolsep}{3pt}
\begin{longtable}{llrrrrrrrlc}
\caption{All method pairs in the main analysis (seed-averaged). Pooled $\Delta$ (equal) weights subsystems equally and is shown with its 95\% CI; the case-weighted pooled $\Delta$ weights subsystems by case count. \#rev is the number of subsystems whose $\Delta$ has the opposite sign of the equal-weight pooled $\Delta$; \#rev CI counts those whose CI excludes 0. Below line marks $|\mathrm{pooled}\,\Delta| < \mathrm{SD}$ of subsystem $\Delta$. The 95\% prediction interval is reported only where the random-effects fit supplies one. LRT $p$ is Holm-adjusted. Sep.\ marks complete separation.}\\\label{supp:s1-pairs}\\
\toprule
Pair & $k$ & \shortstack{Pooled $\Delta$\\(equal) [CI]} & \shortstack{Pooled $\Delta$\\(case-wtd.)} & \#rev & \shortstack{\#rev\\CI} & \shortstack{SD of\\$\Delta_s$} & \shortstack{Below\\line} & 95\% PI & \shortstack{LRT $p$\\(Holm)} & Sep. \\
\midrule
\endfirsthead
\caption{All method pairs in the main analysis (seed-averaged). Pooled $\Delta$ (equal) weights subsystems equally and is shown with its 95\% CI; the case-weighted pooled $\Delta$ weights subsystems by case count. \#rev is the number of subsystems whose $\Delta$ has the opposite sign of the equal-weight pooled $\Delta$; \#rev CI counts those whose CI excludes 0. Below line marks $|\mathrm{pooled}\,\Delta| < \mathrm{SD}$ of subsystem $\Delta$. The 95\% prediction interval is reported only where the random-effects fit supplies one. LRT $p$ is Holm-adjusted. Sep.\ marks complete separation.}\\
\toprule
Pair & $k$ & \shortstack{Pooled $\Delta$\\(equal) [CI]} & \shortstack{Pooled $\Delta$\\(case-wtd.)} & \#rev & \shortstack{\#rev\\CI} & \shortstack{SD of\\$\Delta_s$} & \shortstack{Below\\line} & 95\% PI & \shortstack{LRT $p$\\(Holm)} & Sep. \\
\midrule
\endhead
\midrule
\multicolumn{11}{r}{\emph{Continued on next page}} \\
\bottomrule
\endfoot
\bottomrule
\endlastfoot
\multicolumn{11}{l}{\emph{RCAEval RE1, Published}} \\
BARO vs CausalRCA & 3 & +0.405 [+0.357, +0.453] & +0.405 & 0 & 0 & 0.276 &  & -- & 9.53$\times 10^{-9}$ &  \\
BARO vs CIRCA & 3 & +0.147 [+0.096, +0.200] & +0.147 & 0 & 0 & 0.084 &  & -- & 0.204 &  \\
BARO vs $\epsilon$-Diagnosis & 3 & +0.656 [+0.608, +0.704] & +0.656 & 0 & 0 & 0.083 &  & -- & 0.054 & \checkmark \\
BARO vs RCD & 3 & +0.375 [+0.319, +0.430] & +0.375 & 0 & 0 & 0.156 &  & -- & 0.001 &  \\
CausalRCA vs CIRCA & 3 & $-$0.259 [$-$0.308, $-$0.207] & $-$0.259 & 0 & 0 & 0.243 &  & -- & 2.30$\times 10^{-6}$ &  \\
CausalRCA vs $\epsilon$-Diagnosis & 3 & +0.251 [+0.211, +0.291] & +0.251 & 0 & 0 & 0.293 & \checkmark & -- & 8.12$\times 10^{-10}$ & \checkmark \\
CausalRCA vs RCD & 3 & $-$0.030 [$-$0.077, +0.017] & $-$0.030 & 2 & 0 & 0.159 & \checkmark & -- & 0.001 &  \\
CIRCA vs $\epsilon$-Diagnosis & 3 & +0.509 [+0.461, +0.560] & +0.509 & 0 & 0 & 0.162 &  & -- & 0.017 & \checkmark \\
CIRCA vs RCD & 3 & +0.228 [+0.167, +0.286] & +0.228 & 0 & 0 & 0.093 &  & -- & 0.057 &  \\
$\epsilon$-Diagnosis vs RCD & 3 & $-$0.281 [$-$0.329, $-$0.234] & $-$0.281 & 0 & 0 & 0.213 &  & -- & 0.001 & \checkmark \\
\addlinespace
\multicolumn{11}{l}{\emph{RCAEval RE1, All}} \\
Alert-count vs BARO & 3 & $-$0.443 [$-$0.501, $-$0.381] & $-$0.443 & 0 & 0 & 0.088 &  & -- & 0.019 &  \\
Alert-count vs CausalRCA & 3 & $-$0.037 [$-$0.088, +0.012] & $-$0.037 & 1 & 1 & 0.216 & \checkmark & -- & 6.98$\times 10^{-6}$ &  \\
Alert-count vs CD-1min & 3 & $-$0.373 [$-$0.432, $-$0.315] & $-$0.373 & 0 & 0 & 0.163 &  & -- & 0.019 &  \\
Alert-count vs CIRCA & 3 & $-$0.296 [$-$0.357, $-$0.232] & $-$0.296 & 0 & 0 & 0.028 &  & -- & 0.406 &  \\
Alert-count vs $\epsilon$-Diagnosis & 3 & +0.213 [+0.165, +0.261] & +0.213 & 0 & 0 & 0.157 &  & -- & 0.053 & \checkmark \\
Alert-count vs max-$|Z|$ & 3 & $-$0.285 [$-$0.347, $-$0.221] & $-$0.285 & 0 & 0 & 0.150 &  & -- & 1.40$\times 10^{-4}$ &  \\
Alert-count vs RCD & 3 & $-$0.068 [$-$0.123, $-$0.012] & $-$0.068 & 0 & 0 & 0.072 & \checkmark & -- & 0.826 &  \\
BARO vs CausalRCA & 3 & +0.405 [+0.357, +0.453] & +0.405 & 0 & 0 & 0.276 &  & -- & 2.65$\times 10^{-8}$ &  \\
BARO vs CD-1min & 3 & +0.069 [+0.016, +0.120] & +0.069 & 1 & 1 & 0.188 & \checkmark & -- & 0.001 &  \\
BARO vs CIRCA & 3 & +0.147 [+0.096, +0.200] & +0.147 & 0 & 0 & 0.084 &  & -- & 0.818 &  \\
BARO vs $\epsilon$-Diagnosis & 3 & +0.656 [+0.608, +0.704] & +0.656 & 0 & 0 & 0.083 &  & -- & 0.131 & \checkmark \\
BARO vs max-$|Z|$ & 3 & +0.157 [+0.107, +0.208] & +0.157 & 0 & 0 & 0.067 &  & -- & 0.880 &  \\
BARO vs RCD & 3 & +0.375 [+0.319, +0.430] & +0.375 & 0 & 0 & 0.156 &  & -- & 0.002 &  \\
CausalRCA vs CD-1min & 3 & $-$0.336 [$-$0.383, $-$0.290] & $-$0.336 & 0 & 0 & 0.361 & \checkmark & -- & 1.87$\times 10^{-14}$ &  \\
CausalRCA vs CIRCA & 3 & $-$0.259 [$-$0.308, $-$0.207] & $-$0.259 & 0 & 0 & 0.243 &  & -- & 6.60$\times 10^{-6}$ &  \\
CausalRCA vs $\epsilon$-Diagnosis & 3 & +0.251 [+0.211, +0.291] & +0.251 & 0 & 0 & 0.293 & \checkmark & -- & 2.11$\times 10^{-9}$ & \checkmark \\
CausalRCA vs max-$|Z|$ & 3 & $-$0.248 [$-$0.304, $-$0.192] & $-$0.248 & 1 & 1 & 0.342 & \checkmark & -- & 1.32$\times 10^{-13}$ &  \\
CausalRCA vs RCD & 3 & $-$0.030 [$-$0.077, +0.017] & $-$0.030 & 2 & 0 & 0.159 & \checkmark & -- & 0.004 &  \\
CD-1min vs CIRCA & 3 & +0.077 [+0.021, +0.133] & +0.077 & 1 & 0 & 0.137 & \checkmark & -- & 0.002 &  \\
CD-1min vs $\epsilon$-Diagnosis & 3 & +0.587 [+0.544, +0.632] & +0.587 & 0 & 0 & 0.269 &  & -- & 0.880 & \checkmark \\
CD-1min vs max-$|Z|$ & 3 & +0.088 [+0.032, +0.147] & +0.088 & 1 & 1 & 0.202 & \checkmark & -- & 1.68$\times 10^{-5}$ &  \\
CD-1min vs RCD & 3 & +0.306 [+0.250, +0.362] & +0.306 & 0 & 0 & 0.202 &  & -- & 9.51$\times 10^{-5}$ &  \\
CIRCA vs $\epsilon$-Diagnosis & 3 & +0.509 [+0.461, +0.560] & +0.509 & 0 & 0 & 0.162 &  & -- & 0.047 & \checkmark \\
CIRCA vs max-$|Z|$ & 3 & +0.011 [$-$0.051, +0.072] & +0.011 & 2 & 0 & 0.139 & \checkmark & -- & 0.055 &  \\
CIRCA vs RCD & 3 & +0.228 [+0.167, +0.286] & +0.228 & 0 & 0 & 0.093 &  & -- & 0.172 &  \\
$\epsilon$-Diagnosis vs max-$|Z|$ & 3 & $-$0.499 [$-$0.552, $-$0.443] & $-$0.499 & 0 & 0 & 0.091 &  & -- & 0.131 & \checkmark \\
$\epsilon$-Diagnosis vs RCD & 3 & $-$0.281 [$-$0.329, $-$0.234] & $-$0.281 & 0 & 0 & 0.213 &  & -- & 0.002 & \checkmark \\
max-$|Z|$ vs RCD & 3 & +0.218 [+0.158, +0.275] & +0.218 & 1 & 0 & 0.221 & \checkmark & -- & 1.40$\times 10^{-6}$ &  \\
\addlinespace
\multicolumn{11}{l}{\emph{RCAEval RE2, Published}} \\
BARO vs CausalRCA & 3 & +0.105 [+0.047, +0.163] & +0.105 & 2 & 1 & 0.307 & \checkmark & -- & 4.11$\times 10^{-7}$ &  \\
BARO vs CIRCA & 3 & $-$0.370 [$-$0.433, $-$0.307] & $-$0.370 & 1 & 1 & 0.437 & \checkmark & -- & 0 &  \\
BARO vs $\epsilon$-Diagnosis & 3 & +0.230 [+0.174, +0.285] & +0.230 & 1 & 0 & 0.391 & \checkmark & -- & 3.33$\times 10^{-14}$ & \checkmark \\
BARO vs MicroRank & 2 & +0.344 [+0.267, +0.422] & +0.344 & 0 & 0 & 0.283 &  & -- & 1.000 & \checkmark \\
BARO vs BARO (multi-source) & 3 & $-$0.015 [$-$0.041, +0.007] & $-$0.015 & 0 & 0 & 0.013 &  & -- & 1.000 &  \\
BARO vs SimpleRCA & 3 & $-$0.211 [$-$0.270, $-$0.156] & $-$0.211 & 0 & 0 & 0.099 &  & -- & 1.000 &  \\
BARO vs TraceRCA & 2 & +0.089 [+0.000, +0.172] & +0.089 & 0 & 0 & 0.063 &  & -- & 1.000 &  \\
CausalRCA vs CIRCA & 3 & $-$0.475 [$-$0.538, $-$0.412] & $-$0.475 & 0 & 0 & 0.147 &  & -- & 0.152 &  \\
CausalRCA vs $\epsilon$-Diagnosis & 3 & +0.125 [+0.078, +0.172] & +0.125 & 1 & 0 & 0.116 &  & -- & 2.63$\times 10^{-4}$ & \checkmark \\
CausalRCA vs MicroRank & 2 & +0.146 [+0.089, +0.204] & +0.146 & 0 & 0 & 0.086 &  & -- & 0.004 & \checkmark \\
CausalRCA vs BARO (multi-source) & 3 & $-$0.120 [$-$0.179, $-$0.062] & $-$0.120 & 2 & 1 & 0.314 & \checkmark & -- & 3.29$\times 10^{-7}$ &  \\
CausalRCA vs SimpleRCA & 3 & $-$0.316 [$-$0.375, $-$0.256] & $-$0.316 & 0 & 0 & 0.253 &  & -- & 0.002 &  \\
CausalRCA vs TraceRCA & 2 & $-$0.109 [$-$0.183, $-$0.035] & $-$0.109 & 1 & 1 & 0.306 & \checkmark & -- & 2.30$\times 10^{-4}$ &  \\
CIRCA vs $\epsilon$-Diagnosis & 3 & +0.600 [+0.537, +0.663] & +0.600 & 0 & 0 & 0.048 &  & -- & 0.152 & \checkmark \\
CIRCA vs MicroRank & 2 & +0.544 [+0.467, +0.622] & +0.544 & 0 & 0 & 0.173 &  & -- & 2.63$\times 10^{-4}$ & \checkmark \\
CIRCA vs BARO (multi-source) & 3 & +0.356 [+0.289, +0.422] & +0.356 & 1 & 1 & 0.446 & \checkmark & -- & 0 &  \\
CIRCA vs SimpleRCA & 3 & +0.159 [+0.089, +0.230] & +0.159 & 1 & 1 & 0.396 & \checkmark & -- & 1.70$\times 10^{-12}$ &  \\
CIRCA vs TraceRCA & 2 & +0.289 [+0.189, +0.389] & +0.289 & 0 & 0 & 0.393 & \checkmark & -- & 6.98$\times 10^{-8}$ &  \\
$\epsilon$-Diagnosis vs MicroRank & 2 & $-$0.044 [$-$0.083, $-$0.006] & $-$0.044 & 1 & 0 & 0.110 & \checkmark & -- & 0.002 & \checkmark \\
$\epsilon$-Diagnosis vs BARO (multi-source) & 3 & $-$0.244 [$-$0.300, $-$0.189] & $-$0.244 & 1 & 0 & 0.401 & \checkmark & -- & 1.15$\times 10^{-14}$ & \checkmark \\
$\epsilon$-Diagnosis vs SimpleRCA & 3 & $-$0.441 [$-$0.500, $-$0.381] & $-$0.441 & 0 & 0 & 0.356 &  & -- & 9.09$\times 10^{-12}$ & \checkmark \\
$\epsilon$-Diagnosis vs TraceRCA & 2 & $-$0.300 [$-$0.361, $-$0.239] & $-$0.300 & 0 & 0 & 0.330 & \checkmark & -- & 0.003 & \checkmark \\
MicroRank vs BARO (multi-source) & 2 & $-$0.367 [$-$0.444, $-$0.289] & $-$0.367 & 0 & 0 & 0.283 &  & -- & 1.000 & \checkmark \\
MicroRank vs SimpleRCA & 2 & $-$0.572 [$-$0.650, $-$0.494] & $-$0.572 & 0 & 0 & 0.149 &  & -- & 0.370 & \checkmark \\
MicroRank vs TraceRCA & 2 & $-$0.256 [$-$0.322, $-$0.189] & $-$0.256 & 0 & 0 & 0.220 &  & -- & 1.000 & \checkmark \\
BARO (multi-source) vs SimpleRCA & 3 & $-$0.196 [$-$0.256, $-$0.137] & $-$0.196 & 0 & 0 & 0.096 &  & -- & 1.000 &  \\
BARO (multi-source) vs TraceRCA & 2 & +0.111 [+0.022, +0.200] & +0.111 & 0 & 0 & 0.063 &  & -- & 1.000 &  \\
SimpleRCA vs TraceRCA & 2 & +0.317 [+0.228, +0.411] & +0.317 & 0 & 0 & 0.071 &  & -- & 1.000 &  \\
\addlinespace
\multicolumn{11}{l}{\emph{RCAEval RE2, All}} \\
Alert-count vs BARO & 3 & $-$0.022 [$-$0.089, +0.044] & $-$0.022 & 2 & 1 & 0.396 & \checkmark & -- & 4.17$\times 10^{-14}$ &  \\
Alert-count vs CausalRCA & 3 & +0.083 [+0.022, +0.147] & +0.083 & 0 & 0 & 0.137 & \checkmark & -- & 0.258 &  \\
Alert-count vs CD-1min & 3 & $-$0.552 [$-$0.615, $-$0.485] & $-$0.552 & 0 & 0 & 0.028 &  & -- & 1.000 &  \\
Alert-count vs CIRCA & 3 & $-$0.393 [$-$0.467, $-$0.319] & $-$0.393 & 0 & 0 & 0.061 &  & -- & 1.000 &  \\
Alert-count vs $\epsilon$-Diagnosis & 3 & +0.207 [+0.152, +0.267] & +0.207 & 0 & 0 & 0.028 &  & -- & 0.182 & \checkmark \\
Alert-count vs max-$|Z|$ & 3 & $-$0.407 [$-$0.481, $-$0.330] & $-$0.407 & 0 & 0 & 0.188 &  & -- & 0.068 &  \\
Alert-count vs MicroRank & 2 & +0.150 [+0.089, +0.217] & +0.150 & 0 & 0 & 0.086 &  & -- & 0.008 & \checkmark \\
Alert-count vs BARO (multi-source) & 3 & $-$0.037 [$-$0.104, +0.033] & $-$0.037 & 2 & 1 & 0.406 & \checkmark & -- & 1.11$\times 10^{-14}$ &  \\
Alert-count vs SimpleRCA & 3 & $-$0.233 [$-$0.304, $-$0.159] & $-$0.233 & 1 & 1 & 0.367 & \checkmark & -- & 2.07$\times 10^{-10}$ &  \\
Alert-count vs TraceRCA & 2 & $-$0.106 [$-$0.183, $-$0.028] & $-$0.106 & 1 & 1 & 0.306 & \checkmark & -- & 3.82$\times 10^{-4}$ &  \\
BARO vs CausalRCA & 3 & +0.105 [+0.047, +0.163] & +0.105 & 2 & 1 & 0.307 & \checkmark & -- & 7.40$\times 10^{-7}$ &  \\
BARO vs CD-1min & 3 & $-$0.530 [$-$0.585, $-$0.474] & $-$0.530 & 0 & 0 & 0.381 &  & -- & 1.83$\times 10^{-12}$ &  \\
BARO vs CIRCA & 3 & $-$0.370 [$-$0.433, $-$0.307] & $-$0.370 & 1 & 1 & 0.437 & \checkmark & -- & 0 &  \\
BARO vs $\epsilon$-Diagnosis & 3 & +0.230 [+0.174, +0.285] & +0.230 & 1 & 0 & 0.391 & \checkmark & -- & 6.13$\times 10^{-14}$ & \checkmark \\
BARO vs max-$|Z|$ & 3 & $-$0.385 [$-$0.444, $-$0.322] & $-$0.385 & 1 & 1 & 0.560 & \checkmark & -- & 0 &  \\
BARO vs MicroRank & 2 & +0.344 [+0.267, +0.422] & +0.344 & 0 & 0 & 0.283 &  & -- & 1.000 & \checkmark \\
BARO vs BARO (multi-source) & 3 & $-$0.015 [$-$0.041, +0.007] & $-$0.015 & 0 & 0 & 0.013 &  & -- & 1.000 &  \\
BARO vs SimpleRCA & 3 & $-$0.211 [$-$0.270, $-$0.156] & $-$0.211 & 0 & 0 & 0.099 &  & -- & 1.000 &  \\
BARO vs TraceRCA & 2 & +0.089 [+0.000, +0.172] & +0.089 & 0 & 0 & 0.063 &  & -- & 1.000 &  \\
CausalRCA vs CD-1min & 3 & $-$0.635 [$-$0.686, $-$0.581] & $-$0.635 & 0 & 0 & 0.143 &  & -- & 0.014 &  \\
CausalRCA vs CIRCA & 3 & $-$0.475 [$-$0.538, $-$0.412] & $-$0.475 & 0 & 0 & 0.147 &  & -- & 0.240 &  \\
CausalRCA vs $\epsilon$-Diagnosis & 3 & +0.125 [+0.078, +0.172] & +0.125 & 1 & 0 & 0.116 &  & -- & 4.67$\times 10^{-4}$ & \checkmark \\
CausalRCA vs max-$|Z|$ & 3 & $-$0.490 [$-$0.549, $-$0.431] & $-$0.490 & 0 & 0 & 0.256 &  & -- & 3.43$\times 10^{-4}$ &  \\
CausalRCA vs MicroRank & 2 & +0.146 [+0.089, +0.204] & +0.146 & 0 & 0 & 0.086 &  & -- & 0.008 & \checkmark \\
CausalRCA vs BARO (multi-source) & 3 & $-$0.120 [$-$0.179, $-$0.062] & $-$0.120 & 2 & 1 & 0.314 & \checkmark & -- & 5.80$\times 10^{-7}$ &  \\
CausalRCA vs SimpleRCA & 3 & $-$0.316 [$-$0.375, $-$0.256] & $-$0.316 & 0 & 0 & 0.253 &  & -- & 0.004 &  \\
CausalRCA vs TraceRCA & 2 & $-$0.109 [$-$0.183, $-$0.035] & $-$0.109 & 1 & 1 & 0.306 & \checkmark & -- & 3.99$\times 10^{-4}$ &  \\
CD-1min vs CIRCA & 3 & +0.159 [+0.100, +0.219] & +0.159 & 0 & 0 & 0.089 &  & -- & 0.833 &  \\
CD-1min vs $\epsilon$-Diagnosis & 3 & +0.759 [+0.704, +0.811] & +0.759 & 0 & 0 & 0.051 &  & -- & 0.832 & \checkmark \\
CD-1min vs max-$|Z|$ & 3 & +0.144 [+0.089, +0.200] & +0.144 & 1 & 0 & 0.215 & \checkmark & -- & 0.007 &  \\
CD-1min vs MicroRank & 2 & +0.700 [+0.622, +0.778] & +0.700 & 0 & 0 & 0.047 &  & -- & 0.023 & \checkmark \\
CD-1min vs BARO (multi-source) & 3 & +0.515 [+0.456, +0.570] & +0.515 & 0 & 0 & 0.391 &  & -- & 5.50$\times 10^{-13}$ &  \\
CD-1min vs SimpleRCA & 3 & +0.319 [+0.259, +0.378] & +0.319 & 1 & 0 & 0.358 & \checkmark & -- & 2.07$\times 10^{-10}$ &  \\
CD-1min vs TraceRCA & 2 & +0.444 [+0.356, +0.533] & +0.444 & 0 & 0 & 0.267 &  & -- & 0.003 &  \\
CIRCA vs $\epsilon$-Diagnosis & 3 & +0.600 [+0.537, +0.663] & +0.600 & 0 & 0 & 0.048 &  & -- & 0.240 & \checkmark \\
CIRCA vs max-$|Z|$ & 3 & $-$0.015 [$-$0.078, +0.048] & $-$0.015 & 1 & 0 & 0.130 & \checkmark & -- & 0.445 &  \\
CIRCA vs MicroRank & 2 & +0.544 [+0.467, +0.622] & +0.544 & 0 & 0 & 0.173 &  & -- & 4.67$\times 10^{-4}$ & \checkmark \\
CIRCA vs BARO (multi-source) & 3 & +0.356 [+0.289, +0.422] & +0.356 & 1 & 1 & 0.446 & \checkmark & -- & 0 &  \\
CIRCA vs SimpleRCA & 3 & +0.159 [+0.089, +0.230] & +0.159 & 1 & 1 & 0.396 & \checkmark & -- & 3.05$\times 10^{-12}$ &  \\
CIRCA vs TraceRCA & 2 & +0.289 [+0.189, +0.389] & +0.289 & 0 & 0 & 0.393 & \checkmark & -- & 1.24$\times 10^{-7}$ &  \\
$\epsilon$-Diagnosis vs max-$|Z|$ & 3 & $-$0.615 [$-$0.674, $-$0.552] & $-$0.615 & 0 & 0 & 0.179 &  & -- & 0.833 & \checkmark \\
$\epsilon$-Diagnosis vs MicroRank & 2 & $-$0.044 [$-$0.083, $-$0.006] & $-$0.044 & 1 & 0 & 0.110 & \checkmark & -- & 0.003 & \checkmark \\
$\epsilon$-Diagnosis vs BARO (multi-source) & 3 & $-$0.244 [$-$0.300, $-$0.189] & $-$0.244 & 1 & 0 & 0.401 & \checkmark & -- & 2.13$\times 10^{-14}$ & \checkmark \\
$\epsilon$-Diagnosis vs SimpleRCA & 3 & $-$0.441 [$-$0.500, $-$0.381] & $-$0.441 & 0 & 0 & 0.356 &  & -- & 1.66$\times 10^{-11}$ & \checkmark \\
$\epsilon$-Diagnosis vs TraceRCA & 2 & $-$0.300 [$-$0.361, $-$0.239] & $-$0.300 & 0 & 0 & 0.330 & \checkmark & -- & 0.007 & \checkmark \\
max-$|Z|$ vs MicroRank & 2 & +0.533 [+0.456, +0.611] & +0.533 & 0 & 0 & 0.346 &  & -- & 4.14$\times 10^{-7}$ & \checkmark \\
max-$|Z|$ vs BARO (multi-source) & 3 & +0.370 [+0.307, +0.433] & +0.370 & 1 & 1 & 0.567 & \checkmark & -- & 0 &  \\
max-$|Z|$ vs SimpleRCA & 3 & +0.174 [+0.111, +0.237] & +0.174 & 1 & 1 & 0.508 & \checkmark & -- & 0 &  \\
max-$|Z|$ vs TraceRCA & 2 & +0.278 [+0.183, +0.372] & +0.278 & 1 & 0 & 0.566 & \checkmark & -- & 1.11$\times 10^{-14}$ &  \\
MicroRank vs BARO (multi-source) & 2 & $-$0.367 [$-$0.444, $-$0.289] & $-$0.367 & 0 & 0 & 0.283 &  & -- & 1.000 & \checkmark \\
MicroRank vs SimpleRCA & 2 & $-$0.572 [$-$0.650, $-$0.494] & $-$0.572 & 0 & 0 & 0.149 &  & -- & 0.555 & \checkmark \\
MicroRank vs TraceRCA & 2 & $-$0.256 [$-$0.322, $-$0.189] & $-$0.256 & 0 & 0 & 0.220 &  & -- & 1.000 & \checkmark \\
BARO (multi-source) vs SimpleRCA & 3 & $-$0.196 [$-$0.256, $-$0.137] & $-$0.196 & 0 & 0 & 0.096 &  & -- & 1.000 &  \\
BARO (multi-source) vs TraceRCA & 2 & +0.111 [+0.022, +0.200] & +0.111 & 0 & 0 & 0.063 &  & -- & 1.000 &  \\
SimpleRCA vs TraceRCA & 2 & +0.317 [+0.228, +0.411] & +0.317 & 0 & 0 & 0.071 &  & -- & 1.000 &  \\
\addlinespace
\multicolumn{11}{l}{\emph{OpenRCA, Published}} \\
BARO vs $\epsilon$-Diagnosis & 4 & +0.086 [+0.054, +0.118] & +0.090 & 0 & 0 & 0.055 &  & [$-$0.153, +0.311] & 1.000 &  \\
BARO vs RCD & 4 & +0.027 [+0.002, +0.051] & +0.038 & 1 & 0 & 0.036 & \checkmark & [$-$0.109, +0.157] & 1.000 &  \\
$\epsilon$-Diagnosis vs RCD & 4 & $-$0.059 [$-$0.089, $-$0.031] & $-$0.052 & 0 & 0 & 0.061 & \checkmark & [$-$0.281, +0.181] & 1.000 &  \\
\addlinespace
\multicolumn{11}{l}{\emph{OpenRCA, All}} \\
Alert-count vs BARO & 4 & +0.062 [+0.030, +0.094] & +0.053 & 0 & 0 & 0.036 &  & [$-$0.052, +0.171] & 1.000 &  \\
Alert-count vs CD-1min & 4 & +0.018 [$-$0.029, +0.064] & +0.012 & 1 & 0 & 0.042 & \checkmark & [$-$0.087, +0.097] & 1.000 &  \\
Alert-count vs $\epsilon$-Diagnosis & 4 & +0.148 [+0.113, +0.185] & +0.143 & 0 & 0 & 0.038 &  & [+0.064, +0.205] & 1.000 &  \\
Alert-count vs max-$|Z|$ & 4 & +0.009 [$-$0.028, +0.046] & +0.016 & 2 & 0 & 0.027 & \checkmark & [$-$0.060, +0.082] & 1.000 &  \\
Alert-count vs RCD & 4 & +0.089 [+0.055, +0.122] & +0.091 & 0 & 0 & 0.033 &  & [+0.017, +0.158] & 1.000 &  \\
BARO vs CD-1min & 4 & $-$0.044 [$-$0.087, $-$0.003] & $-$0.041 & 1 & 0 & 0.061 & \checkmark & [$-$0.270, +0.169] & 1.000 &  \\
BARO vs $\epsilon$-Diagnosis & 4 & +0.086 [+0.054, +0.118] & +0.090 & 0 & 0 & 0.055 &  & [$-$0.153, +0.311] & 1.000 &  \\
BARO vs max-$|Z|$ & 4 & $-$0.053 [$-$0.090, $-$0.018] & $-$0.037 & 1 & 0 & 0.063 & \checkmark & [$-$0.324, +0.221] & 1.000 &  \\
BARO vs RCD & 4 & +0.027 [+0.002, +0.051] & +0.038 & 1 & 0 & 0.036 & \checkmark & [$-$0.109, +0.157] & 1.000 &  \\
CD-1min vs $\epsilon$-Diagnosis & 4 & +0.130 [+0.090, +0.171] & +0.131 & 0 & 0 & 0.014 &  & [+0.050, +0.210] & 1.000 &  \\
CD-1min vs max-$|Z|$ & 4 & $-$0.009 [$-$0.055, +0.037] & +0.003 & 2 & 0 & 0.044 & \checkmark & [$-$0.088, +0.096] & 1.000 &  \\
CD-1min vs RCD & 4 & +0.071 [+0.033, +0.110] & +0.079 & 1 & 0 & 0.070 &  & [$-$0.153, +0.319] & 1.000 &  \\
$\epsilon$-Diagnosis vs max-$|Z|$ & 4 & $-$0.139 [$-$0.177, $-$0.103] & $-$0.127 & 0 & 0 & 0.041 &  & [$-$0.230, $-$0.020] & 1.000 &  \\
$\epsilon$-Diagnosis vs RCD & 4 & $-$0.059 [$-$0.089, $-$0.031] & $-$0.052 & 0 & 0 & 0.061 & \checkmark & [$-$0.281, +0.181] & 1.000 &  \\
max-$|Z|$ vs RCD & 4 & +0.080 [+0.043, +0.116] & +0.075 & 0 & 0 & 0.049 &  & [$-$0.096, +0.251] & 1.000 &  \\
\addlinespace
\multicolumn{11}{l}{\emph{PetShop, Published}} \\
BARO vs CIRCA & 4 & $-$0.243 [$-$0.418, $-$0.058] & $-$0.206 & 0 & 0 & 0.106 &  & [$-$0.442, +0.083] & 1.000 &  \\
BARO vs Counterfactual & 4 & $-$0.007 [$-$0.163, +0.147] & $-$0.044 & 1 & 0 & 0.114 & \checkmark & [$-$0.343, +0.241] & 1.000 &  \\
BARO vs $\epsilon$-Diagnosis & 4 & +0.171 [+0.036, +0.312] & +0.162 & 0 & 0 & 0.055 &  & [$-$0.044, +0.366] & 1.000 & \checkmark \\
BARO vs RCD & 4 & $-$0.135 [$-$0.300, +0.041] & $-$0.074 & 1 & 1 & 0.219 & \checkmark & [$-$0.913, +0.775] & 0.173 & \checkmark \\
CIRCA vs Counterfactual & 4 & +0.236 [+0.055, +0.404] & +0.162 & 1 & 0 & 0.195 &  & [$-$0.497, +0.861] & 1.000 &  \\
CIRCA vs $\epsilon$-Diagnosis & 4 & +0.413 [+0.274, +0.550] & +0.368 & 0 & 0 & 0.110 &  & [+0.103, +0.609] & 1.000 & \checkmark \\
CIRCA vs RCD & 4 & +0.108 [$-$0.089, +0.298] & +0.132 & 1 & 0 & 0.187 & \checkmark & [$-$0.153, +0.523] & 0.233 & \checkmark \\
Counterfactual vs $\epsilon$-Diagnosis & 4 & +0.178 [+0.087, +0.288] & +0.206 & 0 & 0 & 0.088 &  & [$-$0.019, +0.392] & 1.000 & \checkmark \\
Counterfactual vs RCD & 4 & $-$0.127 [$-$0.300, +0.058] & $-$0.029 & 1 & 1 & 0.333 & \checkmark & [$-$1.493, +1.338] & 0.002 & \checkmark \\
$\epsilon$-Diagnosis vs RCD & 4 & $-$0.305 [$-$0.452, $-$0.151] & $-$0.235 & 0 & 0 & 0.258 &  & -- & 1.000 & \checkmark \\
\addlinespace
\multicolumn{11}{l}{\emph{PetShop, All}} \\
Alert-count vs BARO & 4 & $-$0.130 [$-$0.262, $-$0.005] & $-$0.132 & 0 & 0 & 0.103 &  & [$-$0.354, +0.077] & 1.000 & \checkmark \\
Alert-count vs Traversal & 4 & $-$0.543 [$-$0.675, $-$0.401] & $-$0.500 & 0 & 0 & 0.133 &  & [$-$0.870, $-$0.156] & 1.000 & \checkmark \\
Alert-count vs CD-1min & 4 & +0.000 [$-$0.125, +0.115] & +0.000 & 0 & 0 & 0.107 & \checkmark & [$-$0.142, +0.133] & 1.000 & \checkmark \\
Alert-count vs CIRCA & 4 & $-$0.373 [$-$0.529, $-$0.202] & $-$0.338 & 0 & 0 & 0.110 &  & [$-$0.610, $-$0.062] & 1.000 & \checkmark \\
Alert-count vs Counterfactual & 4 & $-$0.137 [$-$0.291, +0.010] & $-$0.176 & 0 & 0 & 0.111 &  & [$-$0.398, +0.054] & 1.000 & \checkmark \\
Alert-count vs $\epsilon$-Diagnosis & 4 & +0.041 [$-$0.062, +0.144] & +0.029 & 0 & 0 & 0.059 & \checkmark & -- & 1.000 & \checkmark \\
Alert-count vs max-$|Z|$ & 4 & $-$0.183 [$-$0.327, $-$0.038] & $-$0.147 & 0 & 0 & 0.143 &  & [$-$0.509, +0.236] & 1.000 & \checkmark \\
Alert-count vs Random walk & 4 & +0.041 [$-$0.062, +0.135] & +0.029 & 1 & 0 & 0.156 & \checkmark & -- & 1.000 & \checkmark \\
Alert-count vs Ranked correlation & 4 & $-$0.567 [$-$0.692, $-$0.440] & $-$0.603 & 0 & 0 & 0.223 &  & [$-$1.392, +0.234] & 1.000 & \checkmark \\
Alert-count vs RCD & 4 & $-$0.264 [$-$0.428, $-$0.096] & $-$0.206 & 1 & 0 & 0.287 & \checkmark & [$-$1.494, +0.998] & 0.509 & \checkmark \\
BARO vs Traversal & 4 & $-$0.413 [$-$0.555, $-$0.269] & $-$0.368 & 0 & 0 & 0.163 &  & [$-$0.769, $-$0.008] & 1.000 &  \\
BARO vs CD-1min & 4 & +0.130 [+0.048, +0.231] & +0.132 & 0 & 0 & 0.017 &  & [$-$0.050, +0.312] & 1.000 & \checkmark \\
BARO vs CIRCA & 4 & $-$0.243 [$-$0.418, $-$0.058] & $-$0.206 & 0 & 0 & 0.106 &  & [$-$0.442, +0.083] & 1.000 &  \\
BARO vs Counterfactual & 4 & $-$0.007 [$-$0.163, +0.147] & $-$0.044 & 1 & 0 & 0.114 & \checkmark & [$-$0.343, +0.241] & 1.000 &  \\
BARO vs $\epsilon$-Diagnosis & 4 & +0.171 [+0.036, +0.312] & +0.162 & 0 & 0 & 0.055 &  & [$-$0.044, +0.366] & 1.000 & \checkmark \\
BARO vs max-$|Z|$ & 4 & $-$0.053 [$-$0.147, +0.019] & $-$0.015 & 1 & 0 & 0.096 & \checkmark & [$-$0.306, +0.272] & 1.000 &  \\
BARO vs Random walk & 4 & +0.171 [+0.036, +0.312] & +0.162 & 0 & 0 & 0.055 &  & [$-$0.044, +0.366] & 1.000 & \checkmark \\
BARO vs Ranked correlation & 4 & $-$0.438 [$-$0.579, $-$0.298] & $-$0.471 & 0 & 0 & 0.140 &  & [$-$0.749, $-$0.149] & 1.000 &  \\
BARO vs RCD & 4 & $-$0.135 [$-$0.300, +0.041] & $-$0.074 & 1 & 1 & 0.219 & \checkmark & [$-$0.913, +0.775] & 0.997 & \checkmark \\
Traversal vs CD-1min & 4 & +0.543 [+0.406, +0.673] & +0.500 & 0 & 0 & 0.153 &  & [+0.087, +0.938] & 1.000 & \checkmark \\
Traversal vs CIRCA & 4 & +0.171 [+0.017, +0.322] & +0.162 & 0 & 0 & 0.063 &  & [$-$0.137, +0.428] & 1.000 &  \\
Traversal vs Counterfactual & 4 & +0.406 [+0.255, +0.548] & +0.324 & 0 & 0 & 0.235 &  & [$-$0.538, +1.308] & 1.000 &  \\
Traversal vs $\epsilon$-Diagnosis & 4 & +0.584 [+0.450, +0.709] & +0.529 & 0 & 0 & 0.152 &  & [+0.057, +1.036] & 1.000 & \checkmark \\
Traversal vs max-$|Z|$ & 4 & +0.361 [+0.221, +0.500] & +0.353 & 0 & 0 & 0.108 &  & [+0.057, +0.646] & 1.000 &  \\
Traversal vs Random walk & 4 & +0.584 [+0.459, +0.704] & +0.529 & 0 & 0 & 0.209 &  & [$-$0.288, +1.453] & 1.000 & \checkmark \\
Traversal vs Ranked correlation & 4 & $-$0.024 [$-$0.185, +0.137] & $-$0.103 & 1 & 1 & 0.302 & \checkmark & [$-$1.298, +1.231] & 1.000 &  \\
Traversal vs RCD & 4 & +0.279 [+0.079, +0.464] & +0.294 & 0 & 0 & 0.216 &  & [+0.005, +0.638] & 1.000 & \checkmark \\
CD-1min vs CIRCA & 4 & $-$0.373 [$-$0.548, $-$0.178] & $-$0.338 & 0 & 0 & 0.096 &  & [$-$0.595, $-$0.037] & 1.000 & \checkmark \\
CD-1min vs Counterfactual & 4 & $-$0.137 [$-$0.264, $-$0.024] & $-$0.176 & 0 & 0 & 0.127 &  & [$-$0.490, +0.150] & 1.000 & \checkmark \\
CD-1min vs $\epsilon$-Diagnosis & 4 & +0.041 [$-$0.062, +0.154] & +0.029 & 0 & 0 & 0.059 & \checkmark & -- & 1.000 & \checkmark \\
CD-1min vs max-$|Z|$ & 4 & $-$0.183 [$-$0.298, $-$0.079] & $-$0.147 & 0 & 0 & 0.084 &  & [$-$0.295, +0.051] & 1.000 & \checkmark \\
CD-1min vs Random walk & 4 & +0.041 [$-$0.062, +0.154] & +0.029 & 0 & 0 & 0.059 & \checkmark & -- & 1.000 & \checkmark \\
CD-1min vs Ranked correlation & 4 & $-$0.567 [$-$0.702, $-$0.433] & $-$0.603 & 0 & 0 & 0.150 &  & [$-$1.021, $-$0.189] & 1.000 & \checkmark \\
CD-1min vs RCD & 4 & $-$0.264 [$-$0.397, $-$0.139] & $-$0.206 & 0 & 0 & 0.206 &  & -- & 1.000 & \checkmark \\
CIRCA vs Counterfactual & 4 & +0.236 [+0.055, +0.404] & +0.162 & 1 & 0 & 0.195 &  & [$-$0.497, +0.861] & 1.000 &  \\
CIRCA vs $\epsilon$-Diagnosis & 4 & +0.413 [+0.274, +0.550] & +0.368 & 0 & 0 & 0.110 &  & [+0.103, +0.609] & 1.000 & \checkmark \\
CIRCA vs max-$|Z|$ & 4 & +0.190 [+0.019, +0.363] & +0.191 & 0 & 0 & 0.051 &  & [$-$0.097, +0.478] & 1.000 &  \\
CIRCA vs Random walk & 4 & +0.413 [+0.281, +0.546] & +0.368 & 0 & 0 & 0.150 &  & [+0.057, +0.665] & 1.000 & \checkmark \\
CIRCA vs Ranked correlation & 4 & $-$0.195 [$-$0.399, +0.007] & $-$0.265 & 1 & 0 & 0.243 & \checkmark & [$-$0.948, +0.400] & 1.000 &  \\
CIRCA vs RCD & 4 & +0.108 [$-$0.089, +0.298] & +0.132 & 1 & 0 & 0.187 & \checkmark & [$-$0.153, +0.523] & 1.000 & \checkmark \\
Counterfactual vs $\epsilon$-Diagnosis & 4 & +0.178 [+0.087, +0.288] & +0.206 & 0 & 0 & 0.088 &  & [$-$0.019, +0.392] & 1.000 & \checkmark \\
Counterfactual vs max-$|Z|$ & 4 & $-$0.046 [$-$0.212, +0.132] & +0.029 & 1 & 1 & 0.204 & \checkmark & [$-$0.681, +0.729] & 1.000 &  \\
Counterfactual vs Random walk & 4 & +0.178 [+0.053, +0.312] & +0.206 & 0 & 0 & 0.135 &  & [$-$0.022, +0.422] & 1.000 & \checkmark \\
Counterfactual vs Ranked correlation & 4 & $-$0.430 [$-$0.606, $-$0.236] & $-$0.426 & 0 & 0 & 0.153 &  & [$-$0.731, $-$0.116] & 1.000 &  \\
Counterfactual vs RCD & 4 & $-$0.127 [$-$0.300, +0.058] & $-$0.029 & 1 & 1 & 0.333 & \checkmark & [$-$1.493, +1.338] & 0.010 & \checkmark \\
$\epsilon$-Diagnosis vs max-$|Z|$ & 4 & $-$0.224 [$-$0.370, $-$0.070] & $-$0.176 & 0 & 0 & 0.124 &  & [$-$0.420, +0.135] & 1.000 & \checkmark \\
$\epsilon$-Diagnosis vs Random walk & 4 & +0.000 [$-$0.094, +0.094] & +0.000 & 0 & 0 & 0.102 & \checkmark & -- & 1.000 & \checkmark \\
$\epsilon$-Diagnosis vs Ranked correlation & 4 & $-$0.608 [$-$0.757, $-$0.445] & $-$0.632 & 0 & 0 & 0.174 &  & [$-$0.916, $-$0.402] & 1.000 & \checkmark \\
$\epsilon$-Diagnosis vs RCD & 4 & $-$0.305 [$-$0.452, $-$0.151] & $-$0.235 & 0 & 0 & 0.258 &  & -- & 1.000 & \checkmark \\
max-$|Z|$ vs Random walk & 4 & +0.224 [+0.077, +0.368] & +0.176 & 0 & 0 & 0.124 &  & [$-$0.135, +0.420] & 1.000 & \checkmark \\
max-$|Z|$ vs Ranked correlation & 4 & $-$0.385 [$-$0.514, $-$0.264] & $-$0.456 & 0 & 0 & 0.216 &  & [$-$1.307, +0.526] & 1.000 &  \\
max-$|Z|$ vs RCD & 4 & $-$0.082 [$-$0.236, +0.072] & $-$0.059 & 1 & 0 & 0.148 & \checkmark & [$-$0.444, +0.412] & 1.000 & \checkmark \\
Random walk vs Ranked correlation & 4 & $-$0.608 [$-$0.743, $-$0.466] & $-$0.632 & 0 & 0 & 0.097 &  & [$-$0.896, $-$0.389] & 1.000 & \checkmark \\
Random walk vs RCD & 4 & $-$0.305 [$-$0.438, $-$0.180] & $-$0.235 & 0 & 0 & 0.214 &  & -- & 1.000 & \checkmark \\
Ranked correlation vs RCD & 4 & +0.303 [+0.139, +0.466] & +0.397 & 0 & 0 & 0.290 &  & [$-$0.940, +1.631] & 0.024 & \checkmark \\
\end{longtable}
}


\subsection*{Per-seed pair results}
The same pair contrasts computed separately on each of the three seeds. Source: \texttt{paper/data/rq1/pairs.csv}.

{\scriptsize\setlength{\tabcolsep}{3pt}
\begin{longtable}{llllrr}
\caption{Per-seed pair results (seeds 0, 1, and 2). Pooled $\Delta$ is the equal-subsystem-weight difference with its 95\% CI. \#rev is the number of subsystems reversing the pooled sign; the parenthetical count is the number of those reversals whose CI excludes 0.}\\\label{supp:s2-seeds}\\
\toprule
Family & Layer & Pair & Seed & Pooled $\Delta$ [CI] & \#rev (\#rev CI) \\
\midrule
\endfirsthead
\caption{Per-seed pair results (seeds 0, 1, and 2). Pooled $\Delta$ is the equal-subsystem-weight difference with its 95\% CI. \#rev is the number of subsystems reversing the pooled sign; the parenthetical count is the number of those reversals whose CI excludes 0.}\\
\toprule
Family & Layer & Pair & Seed & Pooled $\Delta$ [CI] & \#rev (\#rev CI) \\
\midrule
\endhead
\midrule
\multicolumn{6}{r}{\emph{Continued on next page}} \\
\bottomrule
\endfoot
\bottomrule
\endlastfoot
RCAEval RE1 & Published & BARO vs CausalRCA & 0 & +0.363 [+0.307, +0.419] & 0 (0) \\
RCAEval RE1 & Published & BARO vs CausalRCA & 1 & +0.405 [+0.349, +0.459] & 0 (0) \\
RCAEval RE1 & Published & BARO vs CausalRCA & 2 & +0.448 [+0.397, +0.499] & 0 (0) \\
RCAEval RE1 & Published & BARO vs RCD & 0 & +0.389 [+0.333, +0.445] & 0 (0) \\
RCAEval RE1 & Published & BARO vs RCD & 1 & +0.371 [+0.312, +0.429] & 0 (0) \\
RCAEval RE1 & Published & BARO vs RCD & 2 & +0.365 [+0.307, +0.424] & 0 (0) \\
RCAEval RE1 & Published & CausalRCA vs CIRCA & 0 & $-$0.216 [$-$0.275, $-$0.157] & 0 (0) \\
RCAEval RE1 & Published & CausalRCA vs CIRCA & 1 & $-$0.259 [$-$0.315, $-$0.203] & 0 (0) \\
RCAEval RE1 & Published & CausalRCA vs CIRCA & 2 & $-$0.301 [$-$0.352, $-$0.248] & 0 (0) \\
RCAEval RE1 & Published & CausalRCA vs $\epsilon$-Diagnosis & 0 & +0.293 [+0.243, +0.344] & 0 (0) \\
RCAEval RE1 & Published & CausalRCA vs $\epsilon$-Diagnosis & 1 & +0.251 [+0.205, +0.299] & 0 (0) \\
RCAEval RE1 & Published & CausalRCA vs $\epsilon$-Diagnosis & 2 & +0.208 [+0.165, +0.251] & 1 (0) \\
RCAEval RE1 & Published & CausalRCA vs RCD & 0 & +0.027 [$-$0.029, +0.083] & 1 (1) \\
RCAEval RE1 & Published & CausalRCA vs RCD & 1 & $-$0.035 [$-$0.091, +0.021] & 1 (1) \\
RCAEval RE1 & Published & CausalRCA vs RCD & 2 & $-$0.083 [$-$0.136, $-$0.029] & 1 (0) \\
RCAEval RE1 & Published & CIRCA vs RCD & 0 & +0.243 [+0.181, +0.304] & 0 (0) \\
RCAEval RE1 & Published & CIRCA vs RCD & 1 & +0.224 [+0.160, +0.285] & 0 (0) \\
RCAEval RE1 & Published & CIRCA vs RCD & 2 & +0.219 [+0.155, +0.283] & 0 (0) \\
RCAEval RE1 & Published & $\epsilon$-Diagnosis vs RCD & 0 & $-$0.267 [$-$0.317, $-$0.219] & 0 (0) \\
RCAEval RE1 & Published & $\epsilon$-Diagnosis vs RCD & 1 & $-$0.285 [$-$0.336, $-$0.237] & 0 (0) \\
RCAEval RE1 & Published & $\epsilon$-Diagnosis vs RCD & 2 & $-$0.291 [$-$0.341, $-$0.243] & 0 (0) \\
RCAEval RE1 & All & Alert-count vs CausalRCA & 0 & $-$0.080 [$-$0.139, $-$0.021] & 1 (1) \\
RCAEval RE1 & All & Alert-count vs CausalRCA & 1 & $-$0.037 [$-$0.093, +0.019] & 1 (1) \\
RCAEval RE1 & All & Alert-count vs CausalRCA & 2 & +0.005 [$-$0.045, +0.056] & 1 (1) \\
RCAEval RE1 & All & Alert-count vs RCD & 0 & $-$0.053 [$-$0.112, +0.005] & 0 (0) \\
RCAEval RE1 & All & Alert-count vs RCD & 1 & $-$0.072 [$-$0.128, $-$0.016] & 0 (0) \\
RCAEval RE1 & All & Alert-count vs RCD & 2 & $-$0.077 [$-$0.136, $-$0.019] & 0 (0) \\
RCAEval RE1 & All & BARO vs CausalRCA & 0 & +0.363 [+0.307, +0.419] & 0 (0) \\
RCAEval RE1 & All & BARO vs CausalRCA & 1 & +0.405 [+0.349, +0.459] & 0 (0) \\
RCAEval RE1 & All & BARO vs CausalRCA & 2 & +0.448 [+0.397, +0.499] & 0 (0) \\
RCAEval RE1 & All & BARO vs RCD & 0 & +0.389 [+0.333, +0.445] & 0 (0) \\
RCAEval RE1 & All & BARO vs RCD & 1 & +0.371 [+0.312, +0.429] & 0 (0) \\
RCAEval RE1 & All & BARO vs RCD & 2 & +0.365 [+0.307, +0.424] & 0 (0) \\
RCAEval RE1 & All & CausalRCA vs CD-1min & 0 & $-$0.293 [$-$0.349, $-$0.240] & 0 (0) \\
RCAEval RE1 & All & CausalRCA vs CD-1min & 1 & $-$0.336 [$-$0.392, $-$0.283] & 0 (0) \\
RCAEval RE1 & All & CausalRCA vs CD-1min & 2 & $-$0.379 [$-$0.427, $-$0.331] & 0 (0) \\
RCAEval RE1 & All & CausalRCA vs CIRCA & 0 & $-$0.216 [$-$0.275, $-$0.157] & 0 (0) \\
RCAEval RE1 & All & CausalRCA vs CIRCA & 1 & $-$0.259 [$-$0.315, $-$0.203] & 0 (0) \\
RCAEval RE1 & All & CausalRCA vs CIRCA & 2 & $-$0.301 [$-$0.352, $-$0.248] & 0 (0) \\
RCAEval RE1 & All & CausalRCA vs $\epsilon$-Diagnosis & 0 & +0.293 [+0.243, +0.344] & 0 (0) \\
RCAEval RE1 & All & CausalRCA vs $\epsilon$-Diagnosis & 1 & +0.251 [+0.205, +0.299] & 0 (0) \\
RCAEval RE1 & All & CausalRCA vs $\epsilon$-Diagnosis & 2 & +0.208 [+0.165, +0.251] & 1 (0) \\
RCAEval RE1 & All & CausalRCA vs max-$|Z|$ & 0 & $-$0.205 [$-$0.267, $-$0.144] & 1 (1) \\
RCAEval RE1 & All & CausalRCA vs max-$|Z|$ & 1 & $-$0.248 [$-$0.309, $-$0.184] & 1 (0) \\
RCAEval RE1 & All & CausalRCA vs max-$|Z|$ & 2 & $-$0.291 [$-$0.347, $-$0.235] & 1 (1) \\
RCAEval RE1 & All & CausalRCA vs RCD & 0 & +0.027 [$-$0.029, +0.083] & 1 (1) \\
RCAEval RE1 & All & CausalRCA vs RCD & 1 & $-$0.035 [$-$0.091, +0.021] & 1 (1) \\
RCAEval RE1 & All & CausalRCA vs RCD & 2 & $-$0.083 [$-$0.136, $-$0.029] & 1 (0) \\
RCAEval RE1 & All & CD-1min vs RCD & 0 & +0.320 [+0.261, +0.379] & 0 (0) \\
RCAEval RE1 & All & CD-1min vs RCD & 1 & +0.301 [+0.243, +0.360] & 0 (0) \\
RCAEval RE1 & All & CD-1min vs RCD & 2 & +0.296 [+0.235, +0.357] & 0 (0) \\
RCAEval RE1 & All & CIRCA vs RCD & 0 & +0.243 [+0.181, +0.304] & 0 (0) \\
RCAEval RE1 & All & CIRCA vs RCD & 1 & +0.224 [+0.160, +0.285] & 0 (0) \\
RCAEval RE1 & All & CIRCA vs RCD & 2 & +0.219 [+0.155, +0.283] & 0 (0) \\
RCAEval RE1 & All & $\epsilon$-Diagnosis vs RCD & 0 & $-$0.267 [$-$0.317, $-$0.219] & 0 (0) \\
RCAEval RE1 & All & $\epsilon$-Diagnosis vs RCD & 1 & $-$0.285 [$-$0.336, $-$0.237] & 0 (0) \\
RCAEval RE1 & All & $\epsilon$-Diagnosis vs RCD & 2 & $-$0.291 [$-$0.341, $-$0.243] & 0 (0) \\
RCAEval RE1 & All & max-$|Z|$ vs RCD & 0 & +0.232 [+0.171, +0.291] & 1 (0) \\
RCAEval RE1 & All & max-$|Z|$ vs RCD & 1 & +0.213 [+0.149, +0.275] & 1 (0) \\
RCAEval RE1 & All & max-$|Z|$ vs RCD & 2 & +0.208 [+0.144, +0.269] & 1 (0) \\
RCAEval RE2 & Published & BARO vs CausalRCA & 0 & +0.063 [$-$0.004, +0.130] & 2 (1) \\
RCAEval RE2 & Published & BARO vs CausalRCA & 1 & +0.137 [+0.074, +0.200] & 2 (1) \\
RCAEval RE2 & Published & BARO vs CausalRCA & 2 & +0.115 [+0.052, +0.178] & 2 (0) \\
RCAEval RE2 & Published & CausalRCA vs CIRCA & 0 & $-$0.433 [$-$0.504, $-$0.363] & 0 (0) \\
RCAEval RE2 & Published & CausalRCA vs CIRCA & 1 & $-$0.507 [$-$0.574, $-$0.441] & 0 (0) \\
RCAEval RE2 & Published & CausalRCA vs CIRCA & 2 & $-$0.485 [$-$0.556, $-$0.415] & 0 (0) \\
RCAEval RE2 & Published & CausalRCA vs $\epsilon$-Diagnosis & 0 & +0.167 [+0.111, +0.222] & 1 (0) \\
RCAEval RE2 & Published & CausalRCA vs $\epsilon$-Diagnosis & 1 & +0.093 [+0.041, +0.144] & 0 (0) \\
RCAEval RE2 & Published & CausalRCA vs $\epsilon$-Diagnosis & 2 & +0.115 [+0.059, +0.170] & 1 (0) \\
RCAEval RE2 & Published & CausalRCA vs MicroRank & 0 & +0.211 [+0.139, +0.283] & 0 (0) \\
RCAEval RE2 & Published & CausalRCA vs MicroRank & 1 & +0.083 [+0.022, +0.150] & 0 (0) \\
RCAEval RE2 & Published & CausalRCA vs MicroRank & 2 & +0.144 [+0.078, +0.211] & 0 (0) \\
RCAEval RE2 & Published & CausalRCA vs BARO (multi-source) & 0 & $-$0.078 [$-$0.144, $-$0.011] & 2 (0) \\
RCAEval RE2 & Published & CausalRCA vs BARO (multi-source) & 1 & $-$0.152 [$-$0.215, $-$0.089] & 1 (1) \\
RCAEval RE2 & Published & CausalRCA vs BARO (multi-source) & 2 & $-$0.130 [$-$0.193, $-$0.067] & 2 (0) \\
RCAEval RE2 & Published & CausalRCA vs SimpleRCA & 0 & $-$0.274 [$-$0.344, $-$0.204] & 0 (0) \\
RCAEval RE2 & Published & CausalRCA vs SimpleRCA & 1 & $-$0.348 [$-$0.411, $-$0.281] & 0 (0) \\
RCAEval RE2 & Published & CausalRCA vs SimpleRCA & 2 & $-$0.326 [$-$0.393, $-$0.263] & 0 (0) \\
RCAEval RE2 & Published & CausalRCA vs TraceRCA & 0 & $-$0.044 [$-$0.128, +0.039] & 1 (1) \\
RCAEval RE2 & Published & CausalRCA vs TraceRCA & 1 & $-$0.172 [$-$0.256, $-$0.089] & 1 (0) \\
RCAEval RE2 & Published & CausalRCA vs TraceRCA & 2 & $-$0.111 [$-$0.189, $-$0.033] & 1 (0) \\
RCAEval RE2 & All & Alert-count vs CausalRCA & 0 & +0.041 [$-$0.030, +0.115] & 2 (0) \\
RCAEval RE2 & All & Alert-count vs CausalRCA & 1 & +0.115 [+0.052, +0.181] & 0 (0) \\
RCAEval RE2 & All & Alert-count vs CausalRCA & 2 & +0.093 [+0.026, +0.163] & 1 (0) \\
RCAEval RE2 & All & BARO vs CausalRCA & 0 & +0.063 [$-$0.004, +0.130] & 2 (1) \\
RCAEval RE2 & All & BARO vs CausalRCA & 1 & +0.137 [+0.074, +0.200] & 2 (1) \\
RCAEval RE2 & All & BARO vs CausalRCA & 2 & +0.115 [+0.052, +0.178] & 2 (0) \\
RCAEval RE2 & All & CausalRCA vs CD-1min & 0 & $-$0.593 [$-$0.652, $-$0.533] & 0 (0) \\
RCAEval RE2 & All & CausalRCA vs CD-1min & 1 & $-$0.667 [$-$0.726, $-$0.607] & 0 (0) \\
RCAEval RE2 & All & CausalRCA vs CD-1min & 2 & $-$0.644 [$-$0.700, $-$0.585] & 0 (0) \\
RCAEval RE2 & All & CausalRCA vs CIRCA & 0 & $-$0.433 [$-$0.504, $-$0.363] & 0 (0) \\
RCAEval RE2 & All & CausalRCA vs CIRCA & 1 & $-$0.507 [$-$0.574, $-$0.441] & 0 (0) \\
RCAEval RE2 & All & CausalRCA vs CIRCA & 2 & $-$0.485 [$-$0.556, $-$0.415] & 0 (0) \\
RCAEval RE2 & All & CausalRCA vs $\epsilon$-Diagnosis & 0 & +0.167 [+0.111, +0.222] & 1 (0) \\
RCAEval RE2 & All & CausalRCA vs $\epsilon$-Diagnosis & 1 & +0.093 [+0.041, +0.144] & 0 (0) \\
RCAEval RE2 & All & CausalRCA vs $\epsilon$-Diagnosis & 2 & +0.115 [+0.059, +0.170] & 1 (0) \\
RCAEval RE2 & All & CausalRCA vs max-$|Z|$ & 0 & $-$0.448 [$-$0.519, $-$0.378] & 0 (0) \\
RCAEval RE2 & All & CausalRCA vs max-$|Z|$ & 1 & $-$0.522 [$-$0.585, $-$0.459] & 0 (0) \\
RCAEval RE2 & All & CausalRCA vs max-$|Z|$ & 2 & $-$0.500 [$-$0.563, $-$0.437] & 0 (0) \\
RCAEval RE2 & All & CausalRCA vs MicroRank & 0 & +0.211 [+0.139, +0.283] & 0 (0) \\
RCAEval RE2 & All & CausalRCA vs MicroRank & 1 & +0.083 [+0.022, +0.150] & 0 (0) \\
RCAEval RE2 & All & CausalRCA vs MicroRank & 2 & +0.144 [+0.078, +0.211] & 0 (0) \\
RCAEval RE2 & All & CausalRCA vs BARO (multi-source) & 0 & $-$0.078 [$-$0.144, $-$0.011] & 2 (0) \\
RCAEval RE2 & All & CausalRCA vs BARO (multi-source) & 1 & $-$0.152 [$-$0.215, $-$0.089] & 1 (1) \\
RCAEval RE2 & All & CausalRCA vs BARO (multi-source) & 2 & $-$0.130 [$-$0.193, $-$0.067] & 2 (0) \\
RCAEval RE2 & All & CausalRCA vs SimpleRCA & 0 & $-$0.274 [$-$0.344, $-$0.204] & 0 (0) \\
RCAEval RE2 & All & CausalRCA vs SimpleRCA & 1 & $-$0.348 [$-$0.411, $-$0.281] & 0 (0) \\
RCAEval RE2 & All & CausalRCA vs SimpleRCA & 2 & $-$0.326 [$-$0.393, $-$0.263] & 0 (0) \\
RCAEval RE2 & All & CausalRCA vs TraceRCA & 0 & $-$0.044 [$-$0.128, +0.039] & 1 (1) \\
RCAEval RE2 & All & CausalRCA vs TraceRCA & 1 & $-$0.172 [$-$0.256, $-$0.089] & 1 (0) \\
RCAEval RE2 & All & CausalRCA vs TraceRCA & 2 & $-$0.111 [$-$0.189, $-$0.033] & 1 (0) \\
OpenRCA & Published & BARO vs RCD & 0 & +0.025 [+0.001, +0.049] & 0 (0) \\
OpenRCA & Published & BARO vs RCD & 1 & +0.027 [$-$0.004, +0.057] & 1 (0) \\
OpenRCA & Published & BARO vs RCD & 2 & +0.029 [+0.002, +0.057] & 1 (0) \\
OpenRCA & Published & $\epsilon$-Diagnosis vs RCD & 0 & $-$0.061 [$-$0.092, $-$0.032] & 0 (0) \\
OpenRCA & Published & $\epsilon$-Diagnosis vs RCD & 1 & $-$0.059 [$-$0.094, $-$0.026] & 1 (0) \\
OpenRCA & Published & $\epsilon$-Diagnosis vs RCD & 2 & $-$0.056 [$-$0.088, $-$0.027] & 0 (0) \\
OpenRCA & All & Alert-count vs RCD & 0 & +0.087 [+0.054, +0.120] & 0 (0) \\
OpenRCA & All & Alert-count vs RCD & 1 & +0.089 [+0.050, +0.127] & 0 (0) \\
OpenRCA & All & Alert-count vs RCD & 2 & +0.091 [+0.055, +0.127] & 0 (0) \\
OpenRCA & All & BARO vs RCD & 0 & +0.025 [+0.001, +0.049] & 0 (0) \\
OpenRCA & All & BARO vs RCD & 1 & +0.027 [$-$0.004, +0.057] & 1 (0) \\
OpenRCA & All & BARO vs RCD & 2 & +0.029 [+0.002, +0.057] & 1 (0) \\
OpenRCA & All & CD-1min vs RCD & 0 & +0.068 [+0.029, +0.108] & 1 (0) \\
OpenRCA & All & CD-1min vs RCD & 1 & +0.070 [+0.029, +0.113] & 1 (0) \\
OpenRCA & All & CD-1min vs RCD & 2 & +0.073 [+0.032, +0.115] & 1 (0) \\
OpenRCA & All & $\epsilon$-Diagnosis vs RCD & 0 & $-$0.061 [$-$0.092, $-$0.032] & 0 (0) \\
OpenRCA & All & $\epsilon$-Diagnosis vs RCD & 1 & $-$0.059 [$-$0.094, $-$0.026] & 1 (0) \\
OpenRCA & All & $\epsilon$-Diagnosis vs RCD & 2 & $-$0.056 [$-$0.088, $-$0.027] & 0 (0) \\
OpenRCA & All & max-$|Z|$ vs RCD & 0 & +0.078 [+0.040, +0.114] & 0 (0) \\
OpenRCA & All & max-$|Z|$ vs RCD & 1 & +0.080 [+0.040, +0.119] & 0 (0) \\
OpenRCA & All & max-$|Z|$ vs RCD & 2 & +0.083 [+0.044, +0.121] & 0 (0) \\
\end{longtable}
}


\subsection*{Per-subsystem differences}
Accuracy of each published method on each subsystem, and the per-subsystem difference, for published-method pairs. Source: \texttt{paper/data/rq1/pair\_subsystem\_deltas.csv}.

{\scriptsize\setlength{\tabcolsep}{3pt}
\begin{longtable}{llllrrr}
\caption{Per-subsystem accuracy differences for published-method pairs (seed-averaged). $\Delta$ is accuracy of A minus accuracy of B, with its 95\% CI. A row marked $^\ast$ has a $\Delta$ whose sign is opposite to that pair's equal-subsystem-weight pooled $\Delta$.}\\\label{supp:s3-subsystem-deltas}\\
\toprule
Family & Pair & Subsystem & $n$ & Acc.\ A & Acc.\ B & $\Delta$ [CI] \\
\midrule
\endfirsthead
\caption{Per-subsystem accuracy differences for published-method pairs (seed-averaged). $\Delta$ is accuracy of A minus accuracy of B, with its 95\% CI. A row marked $^\ast$ has a $\Delta$ whose sign is opposite to that pair's equal-subsystem-weight pooled $\Delta$.}\\
\toprule
Family & Pair & Subsystem & $n$ & Acc.\ A & Acc.\ B & $\Delta$ [CI] \\
\midrule
\endhead
\midrule
\multicolumn{7}{r}{\emph{Continued on next page}} \\
\bottomrule
\endfoot
\bottomrule
\endlastfoot
RCAEval RE1 & BARO vs CausalRCA & RE1-OB & 125 & 0.736 & 0.616 & +0.120 [+0.043, +0.195] \\
RCAEval RE1 & BARO vs CausalRCA & RE1-SS & 125 & 0.848 & 0.176 & +0.672 [+0.592, +0.752] \\
RCAEval RE1 & BARO vs CausalRCA & RE1-TT & 125 & 0.560 & 0.136 & +0.424 [+0.331, +0.515] \\
RCAEval RE1 & BARO vs CIRCA & RE1-OB & 125 & 0.736 & 0.672 & +0.064 [$-$0.032, +0.160] \\
RCAEval RE1 & BARO vs CIRCA & RE1-SS & 125 & 0.848 & 0.704 & +0.144 [+0.064, +0.232] \\
RCAEval RE1 & BARO vs CIRCA & RE1-TT & 125 & 0.560 & 0.328 & +0.232 [+0.136, +0.328] \\
RCAEval RE1 & BARO vs $\epsilon$-Diagnosis & RE1-OB & 125 & 0.736 & 0.032 & +0.704 [+0.624, +0.784] \\
RCAEval RE1 & BARO vs $\epsilon$-Diagnosis & RE1-SS & 125 & 0.848 & 0.144 & +0.704 [+0.616, +0.784] \\
RCAEval RE1 & BARO vs $\epsilon$-Diagnosis & RE1-TT & 125 & 0.560 & 0.000 & +0.560 [+0.472, +0.640] \\
RCAEval RE1 & BARO vs RCD & RE1-OB & 125 & 0.736 & 0.541 & +0.195 [+0.099, +0.288] \\
RCAEval RE1 & BARO vs RCD & RE1-SS & 125 & 0.848 & 0.389 & +0.459 [+0.355, +0.560] \\
RCAEval RE1 & BARO vs RCD & RE1-TT & 125 & 0.560 & 0.088 & +0.472 [+0.381, +0.560] \\
RCAEval RE1 & CausalRCA vs CIRCA & RE1-OB & 125 & 0.616 & 0.672 & $-$0.056 [$-$0.144, +0.035] \\
RCAEval RE1 & CausalRCA vs CIRCA & RE1-SS & 125 & 0.176 & 0.704 & $-$0.528 [$-$0.613, $-$0.440] \\
RCAEval RE1 & CausalRCA vs CIRCA & RE1-TT & 125 & 0.136 & 0.328 & $-$0.192 [$-$0.280, $-$0.104] \\
RCAEval RE1 & CausalRCA vs $\epsilon$-Diagnosis & RE1-OB & 125 & 0.616 & 0.032 & +0.584 [+0.512, +0.659] \\
RCAEval RE1 & CausalRCA vs $\epsilon$-Diagnosis & RE1-SS & 125 & 0.176 & 0.144 & +0.032 [$-$0.051, +0.117] \\
RCAEval RE1 & CausalRCA vs $\epsilon$-Diagnosis & RE1-TT & 125 & 0.136 & 0.000 & +0.136 [+0.096, +0.179] \\
RCAEval RE1 & CausalRCA vs RCD & RE1-OB & 125 & 0.616 & 0.541 & +0.075 [$-$0.019, +0.165]$^\ast$ \\
RCAEval RE1 & CausalRCA vs RCD & RE1-SS & 125 & 0.176 & 0.389 & $-$0.213 [$-$0.299, $-$0.131] \\
RCAEval RE1 & CausalRCA vs RCD & RE1-TT & 125 & 0.136 & 0.088 & +0.048 [$-$0.019, +0.112]$^\ast$ \\
RCAEval RE1 & CIRCA vs $\epsilon$-Diagnosis & RE1-OB & 125 & 0.672 & 0.032 & +0.640 [+0.552, +0.720] \\
RCAEval RE1 & CIRCA vs $\epsilon$-Diagnosis & RE1-SS & 125 & 0.704 & 0.144 & +0.560 [+0.464, +0.648] \\
RCAEval RE1 & CIRCA vs $\epsilon$-Diagnosis & RE1-TT & 125 & 0.328 & 0.000 & +0.328 [+0.248, +0.408] \\
RCAEval RE1 & CIRCA vs RCD & RE1-OB & 125 & 0.672 & 0.541 & +0.131 [+0.021, +0.237] \\
RCAEval RE1 & CIRCA vs RCD & RE1-SS & 125 & 0.704 & 0.389 & +0.315 [+0.195, +0.432] \\
RCAEval RE1 & CIRCA vs RCD & RE1-TT & 125 & 0.328 & 0.088 & +0.240 [+0.149, +0.328] \\
RCAEval RE1 & $\epsilon$-Diagnosis vs RCD & RE1-OB & 125 & 0.032 & 0.541 & $-$0.509 [$-$0.595, $-$0.424] \\
RCAEval RE1 & $\epsilon$-Diagnosis vs RCD & RE1-SS & 125 & 0.144 & 0.389 & $-$0.245 [$-$0.352, $-$0.141] \\
RCAEval RE1 & $\epsilon$-Diagnosis vs RCD & RE1-TT & 125 & 0.000 & 0.088 & $-$0.088 [$-$0.136, $-$0.048] \\
RCAEval RE2 & BARO vs CausalRCA & RE2-OB & 90 & 0.144 & 0.207 & $-$0.063 [$-$0.159, +0.041]$^\ast$ \\
RCAEval RE2 & BARO vs CausalRCA & RE2-SS & 90 & 0.067 & 0.148 & $-$0.081 [$-$0.163, $-$0.004]$^\ast$ \\
RCAEval RE2 & BARO vs CausalRCA & RE2-TT & 90 & 0.667 & 0.207 & +0.459 [+0.341, +0.574] \\
RCAEval RE2 & BARO vs CIRCA & RE2-OB & 90 & 0.144 & 0.667 & $-$0.522 [$-$0.644, $-$0.400] \\
RCAEval RE2 & BARO vs CIRCA & RE2-SS & 90 & 0.067 & 0.778 & $-$0.711 [$-$0.800, $-$0.611] \\
RCAEval RE2 & BARO vs CIRCA & RE2-TT & 90 & 0.667 & 0.544 & +0.122 [+0.011, +0.233]$^\ast$ \\
RCAEval RE2 & BARO vs $\epsilon$-Diagnosis & RE2-OB & 90 & 0.144 & 0.033 & +0.111 [+0.022, +0.200] \\
RCAEval RE2 & BARO vs $\epsilon$-Diagnosis & RE2-SS & 90 & 0.067 & 0.156 & $-$0.089 [$-$0.189, +0.000]$^\ast$ \\
RCAEval RE2 & BARO vs $\epsilon$-Diagnosis & RE2-TT & 90 & 0.667 & 0.000 & +0.667 [+0.567, +0.767] \\
RCAEval RE2 & BARO vs MicroRank & RE2-OB & 90 & 0.144 & 0.000 & +0.144 [+0.078, +0.222] \\
RCAEval RE2 & BARO vs MicroRank & RE2-TT & 90 & 0.667 & 0.122 & +0.544 [+0.411, +0.667] \\
RCAEval RE2 & BARO vs BARO (multi-source) & RE2-OB & 90 & 0.144 & 0.167 & $-$0.022 [$-$0.067, +0.022] \\
RCAEval RE2 & BARO vs BARO (multi-source) & RE2-SS & 90 & 0.067 & 0.067 & +0.000 [$-$0.033, +0.033] \\
RCAEval RE2 & BARO vs BARO (multi-source) & RE2-TT & 90 & 0.667 & 0.689 & $-$0.022 [$-$0.078, +0.033] \\
RCAEval RE2 & BARO vs SimpleRCA & RE2-OB & 90 & 0.144 & 0.467 & $-$0.322 [$-$0.433, $-$0.211] \\
RCAEval RE2 & BARO vs SimpleRCA & RE2-SS & 90 & 0.067 & 0.244 & $-$0.178 [$-$0.278, $-$0.089] \\
RCAEval RE2 & BARO vs SimpleRCA & RE2-TT & 90 & 0.667 & 0.800 & $-$0.133 [$-$0.222, $-$0.044] \\
RCAEval RE2 & BARO vs TraceRCA & RE2-OB & 90 & 0.144 & 0.100 & +0.044 [$-$0.056, +0.144] \\
RCAEval RE2 & BARO vs TraceRCA & RE2-TT & 90 & 0.667 & 0.533 & +0.133 [$-$0.011, +0.267] \\
RCAEval RE2 & CausalRCA vs CIRCA & RE2-OB & 90 & 0.207 & 0.667 & $-$0.459 [$-$0.567, $-$0.356] \\
RCAEval RE2 & CausalRCA vs CIRCA & RE2-SS & 90 & 0.148 & 0.778 & $-$0.630 [$-$0.722, $-$0.530] \\
RCAEval RE2 & CausalRCA vs CIRCA & RE2-TT & 90 & 0.207 & 0.544 & $-$0.337 [$-$0.456, $-$0.219] \\
RCAEval RE2 & CausalRCA vs $\epsilon$-Diagnosis & RE2-OB & 90 & 0.207 & 0.033 & +0.174 [+0.089, +0.259] \\
RCAEval RE2 & CausalRCA vs $\epsilon$-Diagnosis & RE2-SS & 90 & 0.148 & 0.156 & $-$0.007 [$-$0.104, +0.089]$^\ast$ \\
RCAEval RE2 & CausalRCA vs $\epsilon$-Diagnosis & RE2-TT & 90 & 0.207 & 0.000 & +0.207 [+0.152, +0.267] \\
RCAEval RE2 & CausalRCA vs MicroRank & RE2-OB & 90 & 0.207 & 0.000 & +0.207 [+0.137, +0.278] \\
RCAEval RE2 & CausalRCA vs MicroRank & RE2-TT & 90 & 0.207 & 0.122 & +0.085 [$-$0.007, +0.174] \\
RCAEval RE2 & CausalRCA vs BARO (multi-source) & RE2-OB & 90 & 0.207 & 0.167 & +0.041 [$-$0.067, +0.141]$^\ast$ \\
RCAEval RE2 & CausalRCA vs BARO (multi-source) & RE2-SS & 90 & 0.148 & 0.067 & +0.081 [+0.004, +0.163]$^\ast$ \\
RCAEval RE2 & CausalRCA vs BARO (multi-source) & RE2-TT & 90 & 0.207 & 0.689 & $-$0.481 [$-$0.596, $-$0.363] \\
RCAEval RE2 & CausalRCA vs SimpleRCA & RE2-OB & 90 & 0.207 & 0.467 & $-$0.259 [$-$0.374, $-$0.148] \\
RCAEval RE2 & CausalRCA vs SimpleRCA & RE2-SS & 90 & 0.148 & 0.244 & $-$0.096 [$-$0.204, +0.007] \\
RCAEval RE2 & CausalRCA vs SimpleRCA & RE2-TT & 90 & 0.207 & 0.800 & $-$0.593 [$-$0.681, $-$0.496] \\
RCAEval RE2 & CausalRCA vs TraceRCA & RE2-OB & 90 & 0.207 & 0.100 & +0.107 [+0.011, +0.204]$^\ast$ \\
RCAEval RE2 & CausalRCA vs TraceRCA & RE2-TT & 90 & 0.207 & 0.533 & $-$0.326 [$-$0.437, $-$0.215] \\
RCAEval RE2 & CIRCA vs $\epsilon$-Diagnosis & RE2-OB & 90 & 0.667 & 0.033 & +0.633 [+0.533, +0.744] \\
RCAEval RE2 & CIRCA vs $\epsilon$-Diagnosis & RE2-SS & 90 & 0.778 & 0.156 & +0.622 [+0.500, +0.733] \\
RCAEval RE2 & CIRCA vs $\epsilon$-Diagnosis & RE2-TT & 90 & 0.544 & 0.000 & +0.544 [+0.444, +0.644] \\
RCAEval RE2 & CIRCA vs MicroRank & RE2-OB & 90 & 0.667 & 0.000 & +0.667 [+0.567, +0.767] \\
RCAEval RE2 & CIRCA vs MicroRank & RE2-TT & 90 & 0.544 & 0.122 & +0.422 [+0.300, +0.544] \\
RCAEval RE2 & CIRCA vs BARO (multi-source) & RE2-OB & 90 & 0.667 & 0.167 & +0.500 [+0.378, +0.622] \\
RCAEval RE2 & CIRCA vs BARO (multi-source) & RE2-SS & 90 & 0.778 & 0.067 & +0.711 [+0.622, +0.800] \\
RCAEval RE2 & CIRCA vs BARO (multi-source) & RE2-TT & 90 & 0.544 & 0.689 & $-$0.144 [$-$0.267, $-$0.022]$^\ast$ \\
RCAEval RE2 & CIRCA vs SimpleRCA & RE2-OB & 90 & 0.667 & 0.467 & +0.200 [+0.078, +0.322] \\
RCAEval RE2 & CIRCA vs SimpleRCA & RE2-SS & 90 & 0.778 & 0.244 & +0.533 [+0.422, +0.644] \\
RCAEval RE2 & CIRCA vs SimpleRCA & RE2-TT & 90 & 0.544 & 0.800 & $-$0.256 [$-$0.378, $-$0.133]$^\ast$ \\
RCAEval RE2 & CIRCA vs TraceRCA & RE2-OB & 90 & 0.667 & 0.100 & +0.567 [+0.456, +0.678] \\
RCAEval RE2 & CIRCA vs TraceRCA & RE2-TT & 90 & 0.544 & 0.533 & +0.011 [$-$0.144, +0.167] \\
RCAEval RE2 & $\epsilon$-Diagnosis vs MicroRank & RE2-OB & 90 & 0.033 & 0.000 & +0.033 [+0.000, +0.078]$^\ast$ \\
RCAEval RE2 & $\epsilon$-Diagnosis vs MicroRank & RE2-TT & 90 & 0.000 & 0.122 & $-$0.122 [$-$0.189, $-$0.056] \\
RCAEval RE2 & $\epsilon$-Diagnosis vs BARO (multi-source) & RE2-OB & 90 & 0.033 & 0.167 & $-$0.133 [$-$0.222, $-$0.044] \\
RCAEval RE2 & $\epsilon$-Diagnosis vs BARO (multi-source) & RE2-SS & 90 & 0.156 & 0.067 & +0.089 [$-$0.000, +0.189]$^\ast$ \\
RCAEval RE2 & $\epsilon$-Diagnosis vs BARO (multi-source) & RE2-TT & 90 & 0.000 & 0.689 & $-$0.689 [$-$0.789, $-$0.589] \\
RCAEval RE2 & $\epsilon$-Diagnosis vs SimpleRCA & RE2-OB & 90 & 0.033 & 0.467 & $-$0.433 [$-$0.544, $-$0.322] \\
RCAEval RE2 & $\epsilon$-Diagnosis vs SimpleRCA & RE2-SS & 90 & 0.156 & 0.244 & $-$0.089 [$-$0.211, +0.033] \\
RCAEval RE2 & $\epsilon$-Diagnosis vs SimpleRCA & RE2-TT & 90 & 0.000 & 0.800 & $-$0.800 [$-$0.878, $-$0.711] \\
RCAEval RE2 & $\epsilon$-Diagnosis vs TraceRCA & RE2-OB & 90 & 0.033 & 0.100 & $-$0.067 [$-$0.144, +0.011] \\
RCAEval RE2 & $\epsilon$-Diagnosis vs TraceRCA & RE2-TT & 90 & 0.000 & 0.533 & $-$0.533 [$-$0.633, $-$0.433] \\
RCAEval RE2 & MicroRank vs BARO (multi-source) & RE2-OB & 90 & 0.000 & 0.167 & $-$0.167 [$-$0.244, $-$0.089] \\
RCAEval RE2 & MicroRank vs BARO (multi-source) & RE2-TT & 90 & 0.122 & 0.689 & $-$0.567 [$-$0.689, $-$0.433] \\
RCAEval RE2 & MicroRank vs SimpleRCA & RE2-OB & 90 & 0.000 & 0.467 & $-$0.467 [$-$0.567, $-$0.367] \\
RCAEval RE2 & MicroRank vs SimpleRCA & RE2-TT & 90 & 0.122 & 0.800 & $-$0.678 [$-$0.789, $-$0.544] \\
RCAEval RE2 & MicroRank vs TraceRCA & RE2-OB & 90 & 0.000 & 0.100 & $-$0.100 [$-$0.167, $-$0.044] \\
RCAEval RE2 & MicroRank vs TraceRCA & RE2-TT & 90 & 0.122 & 0.533 & $-$0.411 [$-$0.522, $-$0.300] \\
RCAEval RE2 & BARO (multi-source) vs SimpleRCA & RE2-OB & 90 & 0.167 & 0.467 & $-$0.300 [$-$0.411, $-$0.178] \\
RCAEval RE2 & BARO (multi-source) vs SimpleRCA & RE2-SS & 90 & 0.067 & 0.244 & $-$0.178 [$-$0.278, $-$0.078] \\
RCAEval RE2 & BARO (multi-source) vs SimpleRCA & RE2-TT & 90 & 0.689 & 0.800 & $-$0.111 [$-$0.200, $-$0.033] \\
RCAEval RE2 & BARO (multi-source) vs TraceRCA & RE2-OB & 90 & 0.167 & 0.100 & +0.067 [$-$0.044, +0.178] \\
RCAEval RE2 & BARO (multi-source) vs TraceRCA & RE2-TT & 90 & 0.689 & 0.533 & +0.156 [+0.011, +0.300] \\
RCAEval RE2 & SimpleRCA vs TraceRCA & RE2-OB & 90 & 0.467 & 0.100 & +0.367 [+0.244, +0.489] \\
RCAEval RE2 & SimpleRCA vs TraceRCA & RE2-TT & 90 & 0.800 & 0.533 & +0.267 [+0.122, +0.400] \\
OpenRCA & BARO vs $\epsilon$-Diagnosis & Bank & 136 & 0.160 & 0.036 & +0.125 [+0.067, +0.184] \\
OpenRCA & BARO vs $\epsilon$-Diagnosis & Market-1 & 70 & 0.057 & 0.038 & +0.019 [$-$0.020, +0.059] \\
OpenRCA & BARO vs $\epsilon$-Diagnosis & Market-2 & 78 & 0.106 & 0.044 & +0.062 [+0.005, +0.122] \\
OpenRCA & BARO vs $\epsilon$-Diagnosis & Telecom & 51 & 0.183 & 0.046 & +0.137 [+0.052, +0.232] \\
OpenRCA & BARO vs RCD & Bank & 136 & 0.160 & 0.089 & +0.071 [+0.018, +0.127] \\
OpenRCA & BARO vs RCD & Market-1 & 70 & 0.057 & 0.051 & +0.006 [$-$0.035, +0.049] \\
OpenRCA & BARO vs RCD & Market-2 & 78 & 0.106 & 0.066 & +0.040 [$-$0.013, +0.093] \\
OpenRCA & BARO vs RCD & Telecom & 51 & 0.183 & 0.193 & $-$0.010 [$-$0.058, +0.035]$^\ast$ \\
OpenRCA & $\epsilon$-Diagnosis vs RCD & Bank & 136 & 0.036 & 0.089 & $-$0.054 [$-$0.091, $-$0.018] \\
OpenRCA & $\epsilon$-Diagnosis vs RCD & Market-1 & 70 & 0.038 & 0.051 & $-$0.013 [$-$0.054, +0.027] \\
OpenRCA & $\epsilon$-Diagnosis vs RCD & Market-2 & 78 & 0.044 & 0.066 & $-$0.023 [$-$0.076, +0.025] \\
OpenRCA & $\epsilon$-Diagnosis vs RCD & Telecom & 51 & 0.046 & 0.193 & $-$0.147 [$-$0.241, $-$0.062] \\
PetShop & BARO vs CIRCA & High & 26 & 0.154 & 0.269 & $-$0.115 [$-$0.269, +0.038] \\
PetShop & BARO vs CIRCA & Low & 26 & 0.154 & 0.385 & $-$0.231 [$-$0.423, $-$0.038] \\
PetShop & BARO vs CIRCA & Temporal-1 & 8 & 0.250 & 0.500 & $-$0.250 [$-$0.750, +0.250] \\
PetShop & BARO vs CIRCA & Temporal-2 & 8 & 0.250 & 0.625 & $-$0.375 [$-$0.875, +0.125] \\
PetShop & BARO vs Counterfactual & High & 26 & 0.154 & 0.308 & $-$0.154 [$-$0.346, +0.038] \\
PetShop & BARO vs Counterfactual & Low & 26 & 0.154 & 0.154 & +0.000 [$-$0.231, +0.192] \\
PetShop & BARO vs Counterfactual & Temporal-1 & 8 & 0.250 & 0.125 & +0.125 [$-$0.250, +0.500]$^\ast$ \\
PetShop & BARO vs Counterfactual & Temporal-2 & 8 & 0.250 & 0.250 & +0.000 [$-$0.375, +0.375] \\
PetShop & BARO vs $\epsilon$-Diagnosis & High & 26 & 0.154 & 0.000 & +0.154 [+0.038, +0.308] \\
PetShop & BARO vs $\epsilon$-Diagnosis & Low & 26 & 0.154 & 0.000 & +0.154 [+0.038, +0.308] \\
PetShop & BARO vs $\epsilon$-Diagnosis & Temporal-1 & 8 & 0.250 & 0.000 & +0.250 [+0.000, +0.625] \\
PetShop & BARO vs $\epsilon$-Diagnosis & Temporal-2 & 8 & 0.250 & 0.125 & +0.125 [$-$0.250, +0.500] \\
PetShop & BARO vs RCD & High & 26 & 0.154 & 0.000 & +0.154 [+0.038, +0.308]$^\ast$ \\
PetShop & BARO vs RCD & Low & 26 & 0.154 & 0.346 & $-$0.192 [$-$0.423, +0.038] \\
PetShop & BARO vs RCD & Temporal-1 & 8 & 0.250 & 0.625 & $-$0.375 [$-$0.875, +0.125] \\
PetShop & BARO vs RCD & Temporal-2 & 8 & 0.250 & 0.375 & $-$0.125 [$-$0.500, +0.250] \\
PetShop & CIRCA vs Counterfactual & High & 26 & 0.269 & 0.308 & $-$0.038 [$-$0.269, +0.192]$^\ast$ \\
PetShop & CIRCA vs Counterfactual & Low & 26 & 0.385 & 0.154 & +0.231 [$-$0.038, +0.500] \\
PetShop & CIRCA vs Counterfactual & Temporal-1 & 8 & 0.500 & 0.125 & +0.375 [+0.125, +0.750] \\
PetShop & CIRCA vs Counterfactual & Temporal-2 & 8 & 0.625 & 0.250 & +0.375 [$-$0.125, +0.875] \\
PetShop & CIRCA vs $\epsilon$-Diagnosis & High & 26 & 0.269 & 0.000 & +0.269 [+0.115, +0.462] \\
PetShop & CIRCA vs $\epsilon$-Diagnosis & Low & 26 & 0.385 & 0.000 & +0.385 [+0.192, +0.577] \\
PetShop & CIRCA vs $\epsilon$-Diagnosis & Temporal-1 & 8 & 0.500 & 0.000 & +0.500 [+0.125, +0.875] \\
PetShop & CIRCA vs $\epsilon$-Diagnosis & Temporal-2 & 8 & 0.625 & 0.125 & +0.500 [+0.125, +0.875] \\
PetShop & CIRCA vs RCD & High & 26 & 0.269 & 0.000 & +0.269 [+0.115, +0.462] \\
PetShop & CIRCA vs RCD & Low & 26 & 0.385 & 0.346 & +0.038 [$-$0.231, +0.308] \\
PetShop & CIRCA vs RCD & Temporal-1 & 8 & 0.500 & 0.625 & $-$0.125 [$-$0.625, +0.375]$^\ast$ \\
PetShop & CIRCA vs RCD & Temporal-2 & 8 & 0.625 & 0.375 & +0.250 [$-$0.250, +0.750] \\
PetShop & Counterfactual vs $\epsilon$-Diagnosis & High & 26 & 0.308 & 0.000 & +0.308 [+0.154, +0.500] \\
PetShop & Counterfactual vs $\epsilon$-Diagnosis & Low & 26 & 0.154 & 0.000 & +0.154 [+0.038, +0.308] \\
PetShop & Counterfactual vs $\epsilon$-Diagnosis & Temporal-1 & 8 & 0.125 & 0.000 & +0.125 [+0.000, +0.375] \\
PetShop & Counterfactual vs $\epsilon$-Diagnosis & Temporal-2 & 8 & 0.250 & 0.125 & +0.125 [+0.000, +0.375] \\
PetShop & Counterfactual vs RCD & High & 26 & 0.308 & 0.000 & +0.308 [+0.154, +0.500]$^\ast$ \\
PetShop & Counterfactual vs RCD & Low & 26 & 0.154 & 0.346 & $-$0.192 [$-$0.423, +0.038] \\
PetShop & Counterfactual vs RCD & Temporal-1 & 8 & 0.125 & 0.625 & $-$0.500 [$-$0.875, +0.000] \\
PetShop & Counterfactual vs RCD & Temporal-2 & 8 & 0.250 & 0.375 & $-$0.125 [$-$0.500, +0.250] \\
PetShop & $\epsilon$-Diagnosis vs RCD & High & 26 & 0.000 & 0.000 & +0.000 [+0.000, +0.000] \\
PetShop & $\epsilon$-Diagnosis vs RCD & Low & 26 & 0.000 & 0.346 & $-$0.346 [$-$0.538, $-$0.154] \\
PetShop & $\epsilon$-Diagnosis vs RCD & Temporal-1 & 8 & 0.000 & 0.625 & $-$0.625 [$-$0.875, $-$0.250] \\
PetShop & $\epsilon$-Diagnosis vs RCD & Temporal-2 & 8 & 0.125 & 0.375 & $-$0.250 [$-$0.750, +0.250] \\
\end{longtable}
}


\subsection*{Small-sample sensitivity}
How many pairs still reverse when subsystems below a minimum case count are set aside. Source: \texttt{paper/data/rq1/small\_sample.csv}.

{\scriptsize\setlength{\tabcolsep}{3pt}
\begin{longtable}{lllrrr}
\caption{Small-sample sensitivity. For each minimum case count $n_{\min}$, family, and layer, the number of pairs, the number of pairs with at least one subsystem reversal, and the number of pairs with at least one CI-supported reversal.}\\\label{supp:s4-small-sample}\\
\toprule
$n_{\min}$ & Family & Layer & Pairs & \shortstack{Pairs with\\$\geq 1$ reversal} & \shortstack{With CI-supported\\reversal} \\
\midrule
\endfirsthead
\caption{Small-sample sensitivity. For each minimum case count $n_{\min}$, family, and layer, the number of pairs, the number of pairs with at least one subsystem reversal, and the number of pairs with at least one CI-supported reversal.}\\
\toprule
$n_{\min}$ & Family & Layer & Pairs & \shortstack{Pairs with\\$\geq 1$ reversal} & \shortstack{With CI-supported\\reversal} \\
\midrule
\endhead
\midrule
\multicolumn{6}{r}{\emph{Continued on next page}} \\
\bottomrule
\endfoot
\bottomrule
\endlastfoot
0 & RCAEval RE1 & Published & 10 & 1 & 0 \\
0 & RCAEval RE1 & All & 28 & 8 & 4 \\
0 & OpenRCA & Published & 3 & 1 & 0 \\
0 & OpenRCA & All & 15 & 7 & 0 \\
0 & PetShop & Published & 10 & 5 & 2 \\
0 & PetShop & All & 55 & 12 & 4 \\
25 & RCAEval RE1 & Published & 10 & 1 & 0 \\
25 & RCAEval RE1 & All & 28 & 8 & 4 \\
25 & OpenRCA & Published & 3 & 1 & 0 \\
25 & OpenRCA & All & 15 & 7 & 0 \\
25 & PetShop & Published & 10 & 3 & 1 \\
25 & PetShop & All & 55 & 7 & 1 \\
50 & RCAEval RE1 & Published & 10 & 1 & 0 \\
50 & RCAEval RE1 & All & 28 & 8 & 4 \\
50 & OpenRCA & Published & 3 & 1 & 0 \\
50 & OpenRCA & All & 15 & 7 & 0 \\
\end{longtable}
}


\subsection*{Coverage and accuracy}
Valid-output coverage and accuracy under both denominators, for every subsystem, method, and seed. Source: \texttt{paper/data/rq2/coverage.csv}.

{\scriptsize\setlength{\tabcolsep}{3pt}
\begin{longtable}{llrrrrrr}
\caption{Output coverage and accuracy for every subsystem, method, and seed. Valid is the number of cases with a usable ranking. Accuracy (all cases) uses every case as the denominator (an unusable output counts as a miss); accuracy (valid only) uses only cases with a usable ranking. A seed of -- means the method was not re-run across seeds.}\\\label{supp:s5-coverage}\\
\toprule
Subsystem & Method & Seed & $n$ & Valid & Coverage & Acc.\ (all cases) & Acc.\ (valid only) \\
\midrule
\endfirsthead
\caption{Output coverage and accuracy for every subsystem, method, and seed. Valid is the number of cases with a usable ranking. Accuracy (all cases) uses every case as the denominator (an unusable output counts as a miss); accuracy (valid only) uses only cases with a usable ranking. A seed of -- means the method was not re-run across seeds.}\\
\toprule
Subsystem & Method & Seed & $n$ & Valid & Coverage & Acc.\ (all cases) & Acc.\ (valid only) \\
\midrule
\endhead
\midrule
\multicolumn{8}{r}{\emph{Continued on next page}} \\
\bottomrule
\endfoot
\bottomrule
\endlastfoot
RE1-OB & Alert-count & -- & 125 & 125 & 1.000 & 0.392 & 0.392 \\
RE1-OB & BARO & -- & 125 & 125 & 1.000 & 0.736 & 0.736 \\
RE1-OB & CausalRCA & 0 & 125 & 125 & 1.000 & 0.640 & 0.640 \\
RE1-OB & CausalRCA & 1 & 125 & 125 & 1.000 & 0.568 & 0.568 \\
RE1-OB & CausalRCA & 2 & 125 & 125 & 1.000 & 0.640 & 0.640 \\
RE1-OB & CD-1min & -- & 125 & 125 & 1.000 & 0.728 & 0.728 \\
RE1-OB & CIRCA & -- & 125 & 125 & 1.000 & 0.672 & 0.672 \\
RE1-OB & $\epsilon$-Diagnosis & -- & 125 & 125 & 1.000 & 0.032 & 0.032 \\
RE1-OB & max-$|Z|$ & -- & 125 & 125 & 1.000 & 0.504 & 0.504 \\
RE1-OB & RCD & 0 & 125 & 125 & 1.000 & 0.528 & 0.528 \\
RE1-OB & RCD & 1 & 125 & 125 & 1.000 & 0.568 & 0.568 \\
RE1-OB & RCD & 2 & 125 & 125 & 1.000 & 0.528 & 0.528 \\
RE1-SS & Alert-count & -- & 125 & 125 & 1.000 & 0.376 & 0.376 \\
RE1-SS & BARO & -- & 125 & 125 & 1.000 & 0.848 & 0.848 \\
RE1-SS & CausalRCA & 0 & 125 & 125 & 1.000 & 0.256 & 0.256 \\
RE1-SS & CausalRCA & 1 & 125 & 125 & 1.000 & 0.152 & 0.152 \\
RE1-SS & CausalRCA & 2 & 125 & 125 & 1.000 & 0.120 & 0.120 \\
RE1-SS & CD-1min & -- & 125 & 125 & 1.000 & 0.928 & 0.928 \\
RE1-SS & CIRCA & -- & 125 & 125 & 1.000 & 0.704 & 0.704 \\
RE1-SS & $\epsilon$-Diagnosis & -- & 125 & 125 & 1.000 & 0.144 & 0.144 \\
RE1-SS & max-$|Z|$ & -- & 125 & 125 & 1.000 & 0.744 & 0.744 \\
RE1-SS & RCD & 0 & 125 & 125 & 1.000 & 0.376 & 0.376 \\
RE1-SS & RCD & 1 & 125 & 125 & 1.000 & 0.384 & 0.384 \\
RE1-SS & RCD & 2 & 125 & 125 & 1.000 & 0.408 & 0.408 \\
RE1-TT & Alert-count & -- & 125 & 125 & 1.000 & 0.048 & 0.048 \\
RE1-TT & BARO & -- & 125 & 125 & 1.000 & 0.560 & 0.560 \\
RE1-TT & CausalRCA & 0 & 125 & 125 & 1.000 & 0.160 & 0.160 \\
RE1-TT & CausalRCA & 1 & 125 & 125 & 1.000 & 0.208 & 0.208 \\
RE1-TT & CausalRCA & 2 & 125 & 125 & 1.000 & 0.040 & 0.040 \\
RE1-TT & CD-1min & -- & 125 & 125 & 1.000 & 0.280 & 0.280 \\
RE1-TT & CIRCA & -- & 125 & 125 & 1.000 & 0.328 & 0.328 \\
RE1-TT & $\epsilon$-Diagnosis & -- & 125 & 125 & 1.000 & 0.000 & 0.000 \\
RE1-TT & max-$|Z|$ & -- & 125 & 125 & 1.000 & 0.424 & 0.424 \\
RE1-TT & RCD & 0 & 125 & 125 & 1.000 & 0.072 & 0.072 \\
RE1-TT & RCD & 1 & 125 & 125 & 1.000 & 0.080 & 0.080 \\
RE1-TT & RCD & 2 & 125 & 125 & 1.000 & 0.112 & 0.112 \\
RE2-OB & Alert-count & -- & 90 & 90 & 1.000 & 0.211 & 0.211 \\
RE2-OB & BARO & -- & 90 & 90 & 1.000 & 0.144 & 0.144 \\
RE2-OB & CausalRCA & 0 & 90 & 90 & 1.000 & 0.278 & 0.278 \\
RE2-OB & CausalRCA & 1 & 90 & 90 & 1.000 & 0.156 & 0.156 \\
RE2-OB & CausalRCA & 2 & 90 & 90 & 1.000 & 0.189 & 0.189 \\
RE2-OB & CD-1min & -- & 90 & 90 & 1.000 & 0.733 & 0.733 \\
RE2-OB & CIRCA & -- & 90 & 90 & 1.000 & 0.667 & 0.667 \\
RE2-OB & $\epsilon$-Diagnosis & -- & 90 & 90 & 1.000 & 0.033 & 0.033 \\
RE2-OB & max-$|Z|$ & -- & 90 & 90 & 1.000 & 0.778 & 0.778 \\
RE2-OB & MicroRank & -- & 90 & 90 & 1.000 & 0.000 & 0.000 \\
RE2-OB & BARO (multi-source) & -- & 90 & 90 & 1.000 & 0.167 & 0.167 \\
RE2-OB & SimpleRCA & -- & 90 & 90 & 1.000 & 0.467 & 0.467 \\
RE2-OB & TraceRCA & -- & 90 & 90 & 1.000 & 0.100 & 0.100 \\
RE2-SS & Alert-count & -- & 90 & 90 & 1.000 & 0.389 & 0.389 \\
RE2-SS & BARO & -- & 90 & 90 & 1.000 & 0.067 & 0.067 \\
RE2-SS & CausalRCA & 0 & 90 & 90 & 1.000 & 0.144 & 0.144 \\
RE2-SS & CausalRCA & 1 & 90 & 90 & 1.000 & 0.178 & 0.178 \\
RE2-SS & CausalRCA & 2 & 90 & 90 & 1.000 & 0.122 & 0.122 \\
RE2-SS & CD-1min & -- & 90 & 90 & 1.000 & 0.944 & 0.944 \\
RE2-SS & CIRCA & -- & 90 & 90 & 1.000 & 0.778 & 0.778 \\
RE2-SS & $\epsilon$-Diagnosis & -- & 90 & 90 & 1.000 & 0.156 & 0.156 \\
RE2-SS & max-$|Z|$ & -- & 90 & 90 & 1.000 & 0.844 & 0.844 \\
RE2-SS & BARO (multi-source) & -- & 90 & 90 & 1.000 & 0.067 & 0.067 \\
RE2-SS & SimpleRCA & -- & 90 & 90 & 1.000 & 0.244 & 0.244 \\
RE2-TT & Alert-count & -- & 90 & 90 & 1.000 & 0.211 & 0.211 \\
RE2-TT & BARO & -- & 90 & 90 & 1.000 & 0.667 & 0.667 \\
RE2-TT & CausalRCA & 0 & 90 & 90 & 1.000 & 0.267 & 0.267 \\
RE2-TT & CausalRCA & 1 & 90 & 47 & 0.522 & 0.133 & 0.255 \\
RE2-TT & CausalRCA & 2 & 90 & 90 & 1.000 & 0.222 & 0.222 \\
RE2-TT & CD-1min & -- & 90 & 90 & 1.000 & 0.789 & 0.789 \\
RE2-TT & CIRCA & -- & 90 & 90 & 1.000 & 0.544 & 0.544 \\
RE2-TT & $\epsilon$-Diagnosis & -- & 90 & 90 & 1.000 & 0.000 & 0.000 \\
RE2-TT & max-$|Z|$ & -- & 90 & 90 & 1.000 & 0.411 & 0.411 \\
RE2-TT & MicroRank & -- & 90 & 73 & 0.811 & 0.122 & 0.151 \\
RE2-TT & BARO (multi-source) & -- & 90 & 89 & 0.989 & 0.689 & 0.697 \\
RE2-TT & SimpleRCA & -- & 90 & 90 & 1.000 & 0.800 & 0.800 \\
RE2-TT & TraceRCA & -- & 90 & 73 & 0.811 & 0.533 & 0.658 \\
Bank & Alert-count & -- & 136 & 136 & 1.000 & 0.172 & 0.172 \\
Bank & BARO & -- & 136 & 136 & 1.000 & 0.160 & 0.160 \\
Bank & CD-1min & -- & 136 & 132 & 0.971 & 0.165 & 0.170 \\
Bank & $\epsilon$-Diagnosis & -- & 136 & 136 & 1.000 & 0.036 & 0.036 \\
Bank & max-$|Z|$ & -- & 136 & 136 & 1.000 & 0.127 & 0.127 \\
Bank & RCD & 0 & 136 & 136 & 1.000 & 0.094 & 0.094 \\
Bank & RCD & 1 & 136 & 136 & 1.000 & 0.087 & 0.087 \\
Bank & RCD & 2 & 136 & 136 & 1.000 & 0.086 & 0.086 \\
Market-1 & Alert-count & -- & 70 & 70 & 1.000 & 0.144 & 0.144 \\
Market-1 & BARO & -- & 70 & 70 & 1.000 & 0.057 & 0.057 \\
Market-1 & CD-1min & -- & 70 & 70 & 1.000 & 0.161 & 0.161 \\
Market-1 & $\epsilon$-Diagnosis & -- & 70 & 70 & 1.000 & 0.038 & 0.038 \\
Market-1 & max-$|Z|$ & -- & 70 & 70 & 1.000 & 0.161 & 0.161 \\
Market-1 & RCD & 0 & 70 & 63 & 0.900 & 0.057 & 0.063 \\
Market-1 & RCD & 1 & 70 & 62 & 0.886 & 0.036 & 0.040 \\
Market-1 & RCD & 2 & 70 & 63 & 0.900 & 0.060 & 0.066 \\
Market-2 & Alert-count & -- & 78 & 78 & 1.000 & 0.197 & 0.197 \\
Market-2 & BARO & -- & 78 & 78 & 1.000 & 0.106 & 0.106 \\
Market-2 & CD-1min & -- & 78 & 78 & 1.000 & 0.192 & 0.192 \\
Market-2 & $\epsilon$-Diagnosis & -- & 78 & 78 & 1.000 & 0.044 & 0.044 \\
Market-2 & max-$|Z|$ & -- & 78 & 78 & 1.000 & 0.200 & 0.200 \\
Market-2 & RCD & 0 & 78 & 74 & 0.949 & 0.074 & 0.078 \\
Market-2 & RCD & 1 & 78 & 73 & 0.936 & 0.058 & 0.062 \\
Market-2 & RCD & 2 & 78 & 74 & 0.949 & 0.067 & 0.071 \\
Telecom & Alert-count & -- & 51 & 51 & 1.000 & 0.242 & 0.242 \\
Telecom & BARO & -- & 51 & 51 & 1.000 & 0.183 & 0.183 \\
Telecom & CD-1min & -- & 51 & 47 & 0.922 & 0.163 & 0.177 \\
Telecom & $\epsilon$-Diagnosis & -- & 51 & 19 & 0.373 & 0.046 & 0.123 \\
Telecom & max-$|Z|$ & -- & 51 & 51 & 1.000 & 0.232 & 0.232 \\
Telecom & RCD & 0 & 51 & 50 & 0.980 & 0.183 & 0.186 \\
Telecom & RCD & 1 & 51 & 50 & 0.980 & 0.219 & 0.223 \\
Telecom & RCD & 2 & 51 & 49 & 0.961 & 0.176 & 0.183 \\
High & Alert-count & -- & 26 & 26 & 1.000 & 0.038 & 0.038 \\
High & BARO & -- & 26 & 26 & 1.000 & 0.154 & 0.154 \\
High & Traversal & -- & 26 & 24 & 0.923 & 0.385 & 0.417 \\
High & CD-1min & -- & 26 & 26 & 1.000 & 0.000 & 0.000 \\
High & CIRCA & -- & 26 & 26 & 1.000 & 0.269 & 0.269 \\
High & Counterfactual & -- & 26 & 26 & 1.000 & 0.308 & 0.308 \\
High & $\epsilon$-Diagnosis & -- & 26 & 16 & 0.615 & 0.000 & 0.000 \\
High & max-$|Z|$ & -- & 26 & 26 & 1.000 & 0.077 & 0.077 \\
High & Random walk & -- & 26 & 24 & 0.923 & 0.000 & 0.000 \\
High & Ranked correlation & -- & 26 & 24 & 0.923 & 0.731 & 0.792 \\
High & RCD & -- & 26 & 26 & 1.000 & 0.000 & 0.000 \\
Low & Alert-count & -- & 26 & 26 & 1.000 & 0.000 & 0.000 \\
Low & BARO & -- & 26 & 26 & 1.000 & 0.154 & 0.154 \\
Low & Traversal & -- & 26 & 23 & 0.885 & 0.577 & 0.652 \\
Low & CD-1min & -- & 26 & 26 & 1.000 & 0.038 & 0.038 \\
Low & CIRCA & -- & 26 & 26 & 1.000 & 0.385 & 0.385 \\
Low & Counterfactual & -- & 26 & 26 & 1.000 & 0.154 & 0.154 \\
Low & $\epsilon$-Diagnosis & -- & 26 & 17 & 0.654 & 0.000 & 0.000 \\
Low & max-$|Z|$ & -- & 26 & 26 & 1.000 & 0.192 & 0.192 \\
Low & Random walk & -- & 26 & 23 & 0.885 & 0.000 & 0.000 \\
Low & Ranked correlation & -- & 26 & 23 & 0.885 & 0.577 & 0.652 \\
Low & RCD & -- & 26 & 26 & 1.000 & 0.346 & 0.346 \\
Temporal-1 & Alert-count & -- & 8 & 8 & 1.000 & 0.000 & 0.000 \\
Temporal-1 & BARO & -- & 8 & 8 & 1.000 & 0.250 & 0.250 \\
Temporal-1 & Traversal & -- & 8 & 8 & 1.000 & 0.625 & 0.625 \\
Temporal-1 & CD-1min & -- & 8 & 8 & 1.000 & 0.125 & 0.125 \\
Temporal-1 & CIRCA & -- & 8 & 8 & 1.000 & 0.500 & 0.500 \\
Temporal-1 & Counterfactual & -- & 8 & 8 & 1.000 & 0.125 & 0.125 \\
Temporal-1 & $\epsilon$-Diagnosis & -- & 8 & 6 & 0.750 & 0.000 & 0.000 \\
Temporal-1 & max-$|Z|$ & -- & 8 & 8 & 1.000 & 0.375 & 0.375 \\
Temporal-1 & Random walk & -- & 8 & 8 & 1.000 & 0.125 & 0.125 \\
Temporal-1 & Ranked correlation & -- & 8 & 8 & 1.000 & 0.750 & 0.750 \\
Temporal-1 & RCD & -- & 8 & 8 & 1.000 & 0.625 & 0.625 \\
Temporal-2 & Alert-count & -- & 8 & 8 & 1.000 & 0.250 & 0.250 \\
Temporal-2 & BARO & -- & 8 & 8 & 1.000 & 0.250 & 0.250 \\
Temporal-2 & Traversal & -- & 8 & 8 & 1.000 & 0.875 & 0.875 \\
Temporal-2 & CD-1min & -- & 8 & 8 & 1.000 & 0.125 & 0.125 \\
Temporal-2 & CIRCA & -- & 8 & 8 & 1.000 & 0.625 & 0.625 \\
Temporal-2 & Counterfactual & -- & 8 & 8 & 1.000 & 0.250 & 0.250 \\
Temporal-2 & $\epsilon$-Diagnosis & -- & 8 & 8 & 1.000 & 0.125 & 0.125 \\
Temporal-2 & max-$|Z|$ & -- & 8 & 8 & 1.000 & 0.375 & 0.375 \\
Temporal-2 & Random walk & -- & 8 & 8 & 1.000 & 0.000 & 0.000 \\
Temporal-2 & Ranked correlation & -- & 8 & 8 & 1.000 & 0.500 & 0.500 \\
Temporal-2 & RCD & -- & 8 & 8 & 1.000 & 0.375 & 0.375 \\
\end{longtable}
}


\subsection*{Input representation}
Accuracy on the reduced input file and on the raw export, with the file names. Source: \texttt{paper/data/rq2/input\_accuracy.csv}.

{\scriptsize\setlength{\tabcolsep}{3pt}
\begin{longtable}{llrrl}
\caption{Accuracy under the reduced (simple) input released for analysis and under the raw export. The Files column names the reduced file and the raw file.}\\\label{supp:s6-input}\\
\toprule
Subsystem & Method & \shortstack{Acc.\ reduced\\input} & \shortstack{Acc.\ raw\\input} & Files (reduced / raw) \\
\midrule
\endfirsthead
\caption{Accuracy under the reduced (simple) input released for analysis and under the raw export. The Files column names the reduced file and the raw file.}\\
\toprule
Subsystem & Method & \shortstack{Acc.\ reduced\\input} & \shortstack{Acc.\ raw\\input} & Files (reduced / raw) \\
\midrule
\endhead
\midrule
\multicolumn{5}{r}{\emph{Continued on next page}} \\
\bottomrule
\endfoot
\bottomrule
\endlastfoot
RE1-SS & Alert-count & 0.376 & 0.432 & simple\_data.csv / data.csv \\
RE1-SS & BARO & 0.848 & 0.200 & simple\_data.csv / data.csv \\
RE1-SS & CD-1min & 0.928 & 0.736 & simple\_data.csv / data.csv \\
RE1-SS & CIRCA & 0.704 & 0.200 & simple\_data.csv / data.csv \\
RE1-SS & $\epsilon$-Diagnosis & 0.144 & 0.168 & simple\_data.csv / data.csv \\
RE1-SS & max-$|Z|$ & 0.744 & 0.608 & simple\_data.csv / data.csv \\
RE1-SS & RCD & 0.389 & 0.296 & simple\_data.csv / data.csv \\
RE1-TT & Alert-count & 0.048 & 0.192 & simple\_data.csv / data.csv \\
RE1-TT & BARO & 0.560 & 0.120 & simple\_data.csv / data.csv \\
RE1-TT & CD-1min & 0.280 & 0.456 & simple\_data.csv / data.csv \\
RE1-TT & CIRCA & 0.328 & 0.000 & simple\_data.csv / data.csv \\
RE1-TT & $\epsilon$-Diagnosis & 0.000 & 0.000 & simple\_data.csv / data.csv \\
RE1-TT & max-$|Z|$ & 0.424 & 0.448 & simple\_data.csv / data.csv \\
RE1-TT & RCD & 0.088 & 0.091 & simple\_data.csv / data.csv \\
RE2-OB & Alert-count & 0.211 & 0.289 & simple\_metrics.csv / metrics.csv \\
RE2-OB & CD-1min & 0.733 & 0.400 & simple\_metrics.csv / metrics.csv \\
RE2-OB & max-$|Z|$ & 0.778 & 0.522 & simple\_metrics.csv / metrics.csv \\
RE2-OB & BARO (multi-source) & 0.167 & 0.233 & simple\_metrics.csv + logts.csv + tracets\_{lat,err}.csv / metrics.csv + logts.csv + tracets\_{lat,err}.csv \\
RE2-SS & Alert-count & 0.389 & 0.356 & simple\_metrics.csv / metrics.csv \\
RE2-SS & CD-1min & 0.944 & 0.533 & simple\_metrics.csv / metrics.csv \\
RE2-SS & max-$|Z|$ & 0.844 & 0.556 & simple\_metrics.csv / metrics.csv \\
RE2-SS & BARO (multi-source) & 0.067 & 0.267 & simple\_metrics.csv + logts.csv + tracets\_{lat,err}.csv / metrics.csv + logts.csv + tracets\_{lat,err}.csv \\
RE2-TT & Alert-count & 0.211 & 0.278 & simple\_metrics.csv / metrics.csv \\
RE2-TT & CD-1min & 0.789 & 0.244 & simple\_metrics.csv / metrics.csv \\
RE2-TT & max-$|Z|$ & 0.411 & 0.222 & simple\_metrics.csv / metrics.csv \\
RE2-TT & BARO (multi-source) & 0.689 & 0.144 & simple\_metrics.csv + logts.csv + tracets\_{lat,err}.csv / metrics.csv + logts.csv + tracets\_{lat,err}.csv \\
\end{longtable}
}


\subsection*{OpenRCA scoring variants}
Partial-credit, exact-match, and at-least-0.5 accuracy on each OpenRCA subsystem. Source: \texttt{paper/data/rq2/openrca\_scoring\_accuracy.csv}.

{\scriptsize\setlength{\tabcolsep}{3pt}
\begin{longtable}{llrrr}
\caption{OpenRCA scoring variants. Partial credit is the released fractional score; strict counts only a score of exactly 1; $\geq 0.5$ counts a score of at least 0.5.}\\\label{supp:s7-openrca-scoring}\\
\toprule
Subsystem & Method & Partial credit & Strict ($=1$) & $\geq 0.5$ \\
\midrule
\endfirsthead
\caption{OpenRCA scoring variants. Partial credit is the released fractional score; strict counts only a score of exactly 1; $\geq 0.5$ counts a score of at least 0.5.}\\
\toprule
Subsystem & Method & Partial credit & Strict ($=1$) & $\geq 0.5$ \\
\midrule
\endhead
\midrule
\multicolumn{5}{r}{\emph{Continued on next page}} \\
\bottomrule
\endfoot
\bottomrule
\endlastfoot
Bank & Alert-count & 0.172 & 0.081 & 0.235 \\
Bank & BARO & 0.160 & 0.096 & 0.169 \\
Bank & CD-1min & 0.165 & 0.096 & 0.191 \\
Bank & $\epsilon$-Diagnosis & 0.036 & 0.022 & 0.044 \\
Bank & max-$|Z|$ & 0.127 & 0.059 & 0.162 \\
Bank & RCD & 0.089 & 0.049 & 0.081 \\
Market-1 & Alert-count & 0.144 & 0.014 & 0.214 \\
Market-1 & BARO & 0.057 & 0.000 & 0.071 \\
Market-1 & CD-1min & 0.161 & 0.057 & 0.214 \\
Market-1 & $\epsilon$-Diagnosis & 0.038 & 0.000 & 0.043 \\
Market-1 & max-$|Z|$ & 0.161 & 0.029 & 0.229 \\
Market-1 & RCD & 0.051 & 0.000 & 0.043 \\
Market-2 & Alert-count & 0.197 & 0.090 & 0.269 \\
Market-2 & BARO & 0.106 & 0.026 & 0.167 \\
Market-2 & CD-1min & 0.192 & 0.115 & 0.231 \\
Market-2 & $\epsilon$-Diagnosis & 0.044 & 0.000 & 0.051 \\
Market-2 & max-$|Z|$ & 0.200 & 0.077 & 0.282 \\
Market-2 & RCD & 0.066 & 0.026 & 0.051 \\
Telecom & Alert-count & 0.242 & 0.157 & 0.255 \\
Telecom & BARO & 0.183 & 0.118 & 0.196 \\
Telecom & CD-1min & 0.163 & 0.118 & 0.157 \\
Telecom & $\epsilon$-Diagnosis & 0.046 & 0.020 & 0.059 \\
Telecom & max-$|Z|$ & 0.232 & 0.137 & 0.275 \\
Telecom & RCD & 0.193 & 0.124 & 0.216 \\
\end{longtable}
}


\subsection*{Duplicated header: fixed vs as shipped}
Accuracy with the RE1-OB duplicate header fixed and as shipped, and the pair contrasts whose sign or reversal status changes. Source: \texttt{paper/data/rq2/release\_defect\_accuracy.csv}.

{\scriptsize\setlength{\tabcolsep}{3pt}
\begin{longtable}{lrrrrr}
\caption{Accuracy with the released RE1-OB duplicate-header silent failure left as shipped, and with that header fixed. Affected columns restrict both accuracies to the cases whose input row is changed by the fix.}\\\label{supp:s8-release-defect}\\
\toprule
Method & \shortstack{Acc.\\fixed} & \shortstack{Acc.\\as shipped} & \shortstack{Acc.\ fixed\\(affected)} & \shortstack{Acc.\ as shipped\\(affected)} & \shortstack{$n$\\affected} \\
\midrule
\endfirsthead
\caption{Accuracy with the released RE1-OB duplicate-header silent failure left as shipped, and with that header fixed. Affected columns restrict both accuracies to the cases whose input row is changed by the fix.}\\
\toprule
Method & \shortstack{Acc.\\fixed} & \shortstack{Acc.\\as shipped} & \shortstack{Acc.\ fixed\\(affected)} & \shortstack{Acc.\ as shipped\\(affected)} & \shortstack{$n$\\affected} \\
\midrule
\endhead
\midrule
\multicolumn{6}{r}{\emph{Continued on next page}} \\
\bottomrule
\endfoot
\bottomrule
\endlastfoot
CIRCA & 0.672 & 0.432 & 0.820 & 0.220 & 50 \\
$\epsilon$-Diagnosis & 0.032 & 0.032 & 0.000 & 0.000 & 50 \\
RCD & 0.541 & 0.547 & 0.807 & 0.820 & 50 \\
\end{longtable}
}

{\scriptsize\setlength{\tabcolsep}{3pt}
\begin{longtable}{lllll}
\caption{Pair--subsystem comparisons whose sign, or whose reversal status, changes between the fixed release and the release as shipped. $\Delta$ fixed is the reference (header repaired); $\Delta$ as shipped is the alternative.}\\\label{supp:s8-release-defect-changes}\\
\toprule
Layer & Pair & Subsystem & $\Delta$ fixed [CI] & $\Delta$ as shipped [CI] \\
\midrule
\endfirsthead
\caption{Pair--subsystem comparisons whose sign, or whose reversal status, changes between the fixed release and the release as shipped. $\Delta$ fixed is the reference (header repaired); $\Delta$ as shipped is the alternative.}\\
\toprule
Layer & Pair & Subsystem & $\Delta$ fixed [CI] & $\Delta$ as shipped [CI] \\
\midrule
\endhead
\midrule
\multicolumn{5}{r}{\emph{Continued on next page}} \\
\bottomrule
\endfoot
\bottomrule
\endlastfoot
Published & CausalRCA vs CIRCA & RE1-OB & $-$0.056 [$-$0.144, +0.035] & +0.184 [+0.077, +0.291] \\
Published & CIRCA vs RCD & RE1-OB & +0.131 [+0.021, +0.237] & $-$0.115 [$-$0.245, +0.016] \\
All & CausalRCA vs CIRCA & RE1-OB & $-$0.056 [$-$0.144, +0.035] & +0.184 [+0.077, +0.291] \\
All & CIRCA vs max-$|Z|$ & RE1-OB & +0.168 [+0.064, +0.280] & $-$0.072 [$-$0.192, +0.048] \\
All & CIRCA vs max-$|Z|$ & RE1-SS & $-$0.040 [$-$0.144, +0.064] & $-$0.040 [$-$0.144, +0.064] \\
All & CIRCA vs max-$|Z|$ & RE1-TT & $-$0.096 [$-$0.192, +0.008] & $-$0.096 [$-$0.192, +0.008] \\
All & CIRCA vs RCD & RE1-OB & +0.131 [+0.021, +0.237] & $-$0.115 [$-$0.245, +0.016] \\
\end{longtable}
}


\subsection*{Published tables}
Pooled margin against between-stratum dispersion for every pair inside two published tables. Source: \texttt{paper/data/rq1/published\_tables\_margin\_dispersion.csv}.

{\scriptsize\setlength{\tabcolsep}{3pt}
\begin{longtable}{llrrrrrc}
\caption{Margin and between-stratum dispersion inside two published tables (RCAEval Table 6 and PetShop Table 8). Below line marks $|\mathrm{pooled\ margin}| < \mathrm{SD}$ across strata. \#strata reversing is the number of strata whose sign disagrees with the pooled margin.}\\\label{supp:s9-published-tables}\\
\toprule
Table & Pair & $k$ & Pooled & $|\mathrm{margin}|$ & SD & \shortstack{Below\\line} & \shortstack{\#strata\\reversing} \\
\midrule
\endfirsthead
\caption{Margin and between-stratum dispersion inside two published tables (RCAEval Table 6 and PetShop Table 8). Below line marks $|\mathrm{pooled\ margin}| < \mathrm{SD}$ across strata. \#strata reversing is the number of strata whose sign disagrees with the pooled margin.}\\
\toprule
Table & Pair & $k$ & Pooled & $|\mathrm{margin}|$ & SD & \shortstack{Below\\line} & \shortstack{\#strata\\reversing} \\
\midrule
\endhead
\midrule
\multicolumn{8}{r}{\emph{Continued on next page}} \\
\bottomrule
\endfoot
\bottomrule
\endlastfoot
RCAEval Table 6 (RE2-TT, AC@1, 6 fault types) & BARO (metric) vs CausalRCA (metric) & 6 & +0.447 & 0.447 & 0.255 &  & 0 \\
RCAEval Table 6 (RE2-TT, AC@1, 6 fault types) & BARO (metric) vs CIRCA (metric) & 6 & +0.343 & 0.343 & 0.106 &  & 0 \\
RCAEval Table 6 (RE2-TT, AC@1, 6 fault types) & BARO (metric) vs MicroCause (metric) & 6 & +0.568 & 0.568 & 0.213 &  & 0 \\
RCAEval Table 6 (RE2-TT, AC@1, 6 fault types) & BARO (metric) vs RCD (metric) & 6 & +0.578 & 0.578 & 0.282 &  & 0 \\
RCAEval Table 6 (RE2-TT, AC@1, 6 fault types) & BARO (metric) vs MicroRank (trace) & 6 & +0.503 & 0.503 & 0.284 &  & 0 \\
RCAEval Table 6 (RE2-TT, AC@1, 6 fault types) & BARO (metric) vs TraceRCA (trace) & 6 & +0.012 & 0.012 & 0.281 & \checkmark & 3 \\
RCAEval Table 6 (RE2-TT, AC@1, 6 fault types) & BARO (metric) vs BARO (multi) & 6 & $-$0.023 & 0.023 & 0.057 & \checkmark & 0 \\
RCAEval Table 6 (RE2-TT, AC@1, 6 fault types) & BARO (metric) vs CIRCA (multi) & 6 & +0.610 & 0.610 & 0.256 &  & 0 \\
RCAEval Table 6 (RE2-TT, AC@1, 6 fault types) & BARO (metric) vs PDiagnose (multi) & 6 & +0.190 & 0.190 & 0.401 & \checkmark & 2 \\
RCAEval Table 6 (RE2-TT, AC@1, 6 fault types) & BARO (metric) vs RCD (multi) & 6 & +0.572 & 0.572 & 0.283 &  & 0 \\
RCAEval Table 6 (RE2-TT, AC@1, 6 fault types) & CausalRCA (metric) vs CIRCA (metric) & 6 & $-$0.103 & 0.103 & 0.164 & \checkmark & 1 \\
RCAEval Table 6 (RE2-TT, AC@1, 6 fault types) & CausalRCA (metric) vs MicroCause (metric) & 6 & +0.122 & 0.122 & 0.086 &  & 0 \\
RCAEval Table 6 (RE2-TT, AC@1, 6 fault types) & CausalRCA (metric) vs RCD (metric) & 6 & +0.132 & 0.132 & 0.182 & \checkmark & 1 \\
RCAEval Table 6 (RE2-TT, AC@1, 6 fault types) & CausalRCA (metric) vs MicroRank (trace) & 6 & +0.057 & 0.057 & 0.221 & \checkmark & 3 \\
RCAEval Table 6 (RE2-TT, AC@1, 6 fault types) & CausalRCA (metric) vs TraceRCA (trace) & 6 & $-$0.435 & 0.435 & 0.197 &  & 0 \\
RCAEval Table 6 (RE2-TT, AC@1, 6 fault types) & CausalRCA (metric) vs BARO (multi) & 6 & $-$0.470 & 0.470 & 0.267 &  & 0 \\
RCAEval Table 6 (RE2-TT, AC@1, 6 fault types) & CausalRCA (metric) vs CIRCA (multi) & 6 & +0.163 & 0.163 & 0.212 & \checkmark & 1 \\
RCAEval Table 6 (RE2-TT, AC@1, 6 fault types) & CausalRCA (metric) vs PDiagnose (multi) & 6 & $-$0.257 & 0.257 & 0.280 & \checkmark & 1 \\
RCAEval Table 6 (RE2-TT, AC@1, 6 fault types) & CausalRCA (metric) vs RCD (multi) & 6 & +0.125 & 0.125 & 0.146 & \checkmark & 1 \\
RCAEval Table 6 (RE2-TT, AC@1, 6 fault types) & CIRCA (metric) vs MicroCause (metric) & 6 & +0.225 & 0.225 & 0.137 &  & 0 \\
RCAEval Table 6 (RE2-TT, AC@1, 6 fault types) & CIRCA (metric) vs RCD (metric) & 6 & +0.235 & 0.235 & 0.185 &  & 0 \\
RCAEval Table 6 (RE2-TT, AC@1, 6 fault types) & CIRCA (metric) vs MicroRank (trace) & 6 & +0.160 & 0.160 & 0.199 & \checkmark & 1 \\
RCAEval Table 6 (RE2-TT, AC@1, 6 fault types) & CIRCA (metric) vs TraceRCA (trace) & 6 & $-$0.332 & 0.332 & 0.191 &  & 0 \\
RCAEval Table 6 (RE2-TT, AC@1, 6 fault types) & CIRCA (metric) vs BARO (multi) & 6 & $-$0.367 & 0.367 & 0.117 &  & 0 \\
RCAEval Table 6 (RE2-TT, AC@1, 6 fault types) & CIRCA (metric) vs CIRCA (multi) & 6 & +0.267 & 0.267 & 0.173 &  & 0 \\
RCAEval Table 6 (RE2-TT, AC@1, 6 fault types) & CIRCA (metric) vs PDiagnose (multi) & 6 & $-$0.153 & 0.153 & 0.311 & \checkmark & 2 \\
RCAEval Table 6 (RE2-TT, AC@1, 6 fault types) & CIRCA (metric) vs RCD (multi) & 6 & +0.228 & 0.228 & 0.186 &  & 0 \\
RCAEval Table 6 (RE2-TT, AC@1, 6 fault types) & MicroCause (metric) vs RCD (metric) & 6 & +0.010 & 0.010 & 0.203 & \checkmark & 4 \\
RCAEval Table 6 (RE2-TT, AC@1, 6 fault types) & MicroCause (metric) vs MicroRank (trace) & 6 & $-$0.065 & 0.065 & 0.250 & \checkmark & 1 \\
RCAEval Table 6 (RE2-TT, AC@1, 6 fault types) & MicroCause (metric) vs TraceRCA (trace) & 6 & $-$0.557 & 0.557 & 0.203 &  & 0 \\
RCAEval Table 6 (RE2-TT, AC@1, 6 fault types) & MicroCause (metric) vs BARO (multi) & 6 & $-$0.592 & 0.592 & 0.216 &  & 0 \\
RCAEval Table 6 (RE2-TT, AC@1, 6 fault types) & MicroCause (metric) vs CIRCA (multi) & 6 & +0.042 & 0.042 & 0.208 & \checkmark & 4 \\
RCAEval Table 6 (RE2-TT, AC@1, 6 fault types) & MicroCause (metric) vs PDiagnose (multi) & 6 & $-$0.378 & 0.378 & 0.300 &  & 1 \\
RCAEval Table 6 (RE2-TT, AC@1, 6 fault types) & MicroCause (metric) vs RCD (multi) & 6 & +0.003 & 0.003 & 0.179 & \checkmark & 3 \\
RCAEval Table 6 (RE2-TT, AC@1, 6 fault types) & RCD (metric) vs MicroRank (trace) & 6 & $-$0.075 & 0.075 & 0.091 & \checkmark & 1 \\
RCAEval Table 6 (RE2-TT, AC@1, 6 fault types) & RCD (metric) vs TraceRCA (trace) & 6 & $-$0.567 & 0.567 & 0.096 &  & 0 \\
RCAEval Table 6 (RE2-TT, AC@1, 6 fault types) & RCD (metric) vs BARO (multi) & 6 & $-$0.602 & 0.602 & 0.276 &  & 0 \\
RCAEval Table 6 (RE2-TT, AC@1, 6 fault types) & RCD (metric) vs CIRCA (multi) & 6 & +0.032 & 0.032 & 0.066 & \checkmark & 1 \\
RCAEval Table 6 (RE2-TT, AC@1, 6 fault types) & RCD (metric) vs PDiagnose (multi) & 6 & $-$0.388 & 0.388 & 0.194 &  & 0 \\
RCAEval Table 6 (RE2-TT, AC@1, 6 fault types) & RCD (metric) vs RCD (multi) & 6 & $-$0.007 & 0.007 & 0.069 & \checkmark & 2 \\
RCAEval Table 6 (RE2-TT, AC@1, 6 fault types) & MicroRank (trace) vs TraceRCA (trace) & 6 & $-$0.492 & 0.492 & 0.177 &  & 0 \\
RCAEval Table 6 (RE2-TT, AC@1, 6 fault types) & MicroRank (trace) vs BARO (multi) & 6 & $-$0.527 & 0.527 & 0.279 &  & 0 \\
RCAEval Table 6 (RE2-TT, AC@1, 6 fault types) & MicroRank (trace) vs CIRCA (multi) & 6 & +0.107 & 0.107 & 0.111 & \checkmark & 0 \\
RCAEval Table 6 (RE2-TT, AC@1, 6 fault types) & MicroRank (trace) vs PDiagnose (multi) & 6 & $-$0.313 & 0.313 & 0.267 &  & 0 \\
RCAEval Table 6 (RE2-TT, AC@1, 6 fault types) & MicroRank (trace) vs RCD (multi) & 6 & +0.068 & 0.068 & 0.105 & \checkmark & 1 \\
RCAEval Table 6 (RE2-TT, AC@1, 6 fault types) & TraceRCA (trace) vs BARO (multi) & 6 & $-$0.035 & 0.035 & 0.281 & \checkmark & 2 \\
RCAEval Table 6 (RE2-TT, AC@1, 6 fault types) & TraceRCA (trace) vs CIRCA (multi) & 6 & +0.598 & 0.598 & 0.116 &  & 0 \\
RCAEval Table 6 (RE2-TT, AC@1, 6 fault types) & TraceRCA (trace) vs PDiagnose (multi) & 6 & +0.178 & 0.178 & 0.135 &  & 1 \\
RCAEval Table 6 (RE2-TT, AC@1, 6 fault types) & TraceRCA (trace) vs RCD (multi) & 6 & +0.560 & 0.560 & 0.144 &  & 0 \\
RCAEval Table 6 (RE2-TT, AC@1, 6 fault types) & BARO (multi) vs CIRCA (multi) & 6 & +0.633 & 0.633 & 0.239 &  & 0 \\
RCAEval Table 6 (RE2-TT, AC@1, 6 fault types) & BARO (multi) vs PDiagnose (multi) & 6 & +0.213 & 0.213 & 0.405 & \checkmark & 2 \\
RCAEval Table 6 (RE2-TT, AC@1, 6 fault types) & BARO (multi) vs RCD (multi) & 6 & +0.595 & 0.595 & 0.277 &  & 0 \\
RCAEval Table 6 (RE2-TT, AC@1, 6 fault types) & CIRCA (multi) vs PDiagnose (multi) & 6 & $-$0.420 & 0.420 & 0.231 &  & 0 \\
RCAEval Table 6 (RE2-TT, AC@1, 6 fault types) & CIRCA (multi) vs RCD (multi) & 6 & $-$0.038 & 0.038 & 0.103 & \checkmark & 3 \\
RCAEval Table 6 (RE2-TT, AC@1, 6 fault types) & PDiagnose (multi) vs RCD (multi) & 6 & +0.382 & 0.382 & 0.242 &  & 0 \\
\addlinespace
PetShop Table 8 (top-1, 6 strata) & Traversal vs CIRCA & 6 & +0.157 & 0.157 & 0.126 &  & 0 \\
PetShop Table 8 (top-1, 6 strata) & Traversal vs Counterfactual & 6 & +0.340 & 0.340 & 0.218 &  & 0 \\
PetShop Table 8 (top-1, 6 strata) & Traversal vs $\epsilon$-Diagnosis & 6 & +0.538 & 0.538 & 0.194 &  & 0 \\
PetShop Table 8 (top-1, 6 strata) & Traversal vs RCD & 6 & +0.325 & 0.325 & 0.356 & \checkmark & 2 \\
PetShop Table 8 (top-1, 6 strata) & Traversal vs Ranked correlation & 6 & $-$0.090 & 0.090 & 0.313 & \checkmark & 2 \\
PetShop Table 8 (top-1, 6 strata) & CIRCA vs Counterfactual & 6 & +0.183 & 0.183 & 0.228 & \checkmark & 1 \\
PetShop Table 8 (top-1, 6 strata) & CIRCA vs $\epsilon$-Diagnosis & 6 & +0.382 & 0.382 & 0.211 &  & 0 \\
PetShop Table 8 (top-1, 6 strata) & CIRCA vs RCD & 6 & +0.168 & 0.168 & 0.301 & \checkmark & 2 \\
PetShop Table 8 (top-1, 6 strata) & CIRCA vs Ranked correlation & 6 & $-$0.247 & 0.247 & 0.326 & \checkmark & 1 \\
PetShop Table 8 (top-1, 6 strata) & Counterfactual vs $\epsilon$-Diagnosis & 6 & +0.198 & 0.198 & 0.239 & \checkmark & 0 \\
PetShop Table 8 (top-1, 6 strata) & Counterfactual vs RCD & 6 & $-$0.015 & 0.015 & 0.450 & \checkmark & 3 \\
PetShop Table 8 (top-1, 6 strata) & Counterfactual vs Ranked correlation & 6 & $-$0.430 & 0.430 & 0.345 &  & 0 \\
PetShop Table 8 (top-1, 6 strata) & $\epsilon$-Diagnosis vs RCD & 6 & $-$0.213 & 0.213 & 0.259 & \checkmark & 0 \\
PetShop Table 8 (top-1, 6 strata) & $\epsilon$-Diagnosis vs Ranked correlation & 6 & $-$0.628 & 0.628 & 0.149 &  & 0 \\
PetShop Table 8 (top-1, 6 strata) & RCD vs Ranked correlation & 6 & $-$0.415 & 0.415 & 0.274 &  & 0 \\
\end{longtable}
}


\subsection*{Fault composition}
Case counts by fault class, and the comparisons whose sign changes after re-weighting fault classes. Source: \texttt{paper/data/rq2/fault\_composition.csv}.

{\scriptsize\setlength{\tabcolsep}{3pt}
\begin{longtable}{lrrrrrrrrrrr}
\caption{Case counts per fault class and subsystem. Fault classes that are zero in every subsystem are omitted.}\\\label{supp:s10-fault-composition}\\
\toprule
Subsystem & availability & cpu & delay & disk & latency & loss & mem & memory & network & other & socket \\
\midrule
\endfirsthead
\caption{Case counts per fault class and subsystem. Fault classes that are zero in every subsystem are omitted.}\\
\toprule
Subsystem & availability & cpu & delay & disk & latency & loss & mem & memory & network & other & socket \\
\midrule
\endhead
\midrule
\multicolumn{12}{r}{\emph{Continued on next page}} \\
\bottomrule
\endfoot
\bottomrule
\endlastfoot
RE1-OB & 0 & 25 & 25 & 25 & 0 & 25 & 25 & 0 & 0 & 0 & 0 \\
RE1-SS & 0 & 25 & 25 & 25 & 0 & 25 & 25 & 0 & 0 & 0 & 0 \\
RE1-TT & 0 & 25 & 25 & 25 & 0 & 25 & 25 & 0 & 0 & 0 & 0 \\
RE2-OB & 0 & 15 & 15 & 15 & 0 & 15 & 15 & 0 & 0 & 0 & 15 \\
RE2-SS & 0 & 15 & 15 & 15 & 0 & 15 & 15 & 0 & 0 & 0 & 15 \\
RE2-TT & 0 & 15 & 15 & 15 & 0 & 15 & 15 & 0 & 0 & 0 & 15 \\
Bank & 0 & 36 & 0 & 24 & 0 & 0 & 0 & 17 & 59 & 0 & 0 \\
Market-1 & 0 & 13 & 0 & 24 & 0 & 0 & 0 & 12 & 17 & 4 & 0 \\
Market-2 & 0 & 13 & 0 & 27 & 0 & 0 & 0 & 10 & 25 & 3 & 0 \\
Telecom & 0 & 19 & 0 & 0 & 0 & 0 & 0 & 0 & 20 & 12 & 0 \\
High & 12 & 0 & 0 & 0 & 14 & 0 & 0 & 0 & 0 & 0 & 0 \\
Low & 12 & 0 & 0 & 0 & 14 & 0 & 0 & 0 & 0 & 0 & 0 \\
Temporal-1 & 4 & 0 & 0 & 0 & 4 & 0 & 0 & 0 & 0 & 0 & 0 \\
Temporal-2 & 4 & 0 & 0 & 0 & 4 & 0 & 0 & 0 & 0 & 0 & 0 \\
\end{longtable}
}

{\scriptsize\setlength{\tabcolsep}{3pt}
\begin{longtable}{llllrr}
\caption{Comparisons whose sign changes when each fault class is re-weighted to a common mix. $\Delta$ is the original per-subsystem difference; $\Delta$ re-weighted is the re-weighted difference with its 95\% CI.}\\\label{supp:s10-fault-reweighted}\\
\toprule
Family & Layer & Pair & Subsystem & $\Delta$ & $\Delta$ re-weighted [CI] \\
\midrule
\endfirsthead
\caption{Comparisons whose sign changes when each fault class is re-weighted to a common mix. $\Delta$ is the original per-subsystem difference; $\Delta$ re-weighted is the re-weighted difference with its 95\% CI.}\\
\toprule
Family & Layer & Pair & Subsystem & $\Delta$ & $\Delta$ re-weighted [CI] \\
\midrule
\endhead
\midrule
\multicolumn{6}{r}{\emph{Continued on next page}} \\
\bottomrule
\endfoot
\bottomrule
\endlastfoot
OpenRCA & Published & BARO vs $\epsilon$-Diagnosis & Market-1 & +0.019 & $-$0.006 [$-$0.067, +0.054] \\
OpenRCA & Published & $\epsilon$-Diagnosis vs RCD & Market-1 & $-$0.013 & +0.017 [$-$0.052, +0.081] \\
OpenRCA & Published & $\epsilon$-Diagnosis vs RCD & Market-2 & $-$0.023 & +0.039 [$-$0.033, +0.116] \\
OpenRCA & All & Alert-count vs BARO & Telecom & +0.059 & $-$0.004 [$-$0.037, +0.027] \\
OpenRCA & All & Alert-count vs CD-1min & Bank & +0.007 & $-$0.030 [$-$0.121, +0.061] \\
OpenRCA & All & Alert-count vs CD-1min & Market-1 & $-$0.017 & +0.052 [$-$0.025, +0.144] \\
OpenRCA & All & Alert-count vs CD-1min & Market-2 & +0.004 & $-$0.010 [$-$0.146, +0.116] \\
OpenRCA & All & Alert-count vs max-$|Z|$ & Telecom & +0.010 & $-$0.001 [$-$0.090, +0.094] \\
OpenRCA & All & Alert-count vs RCD & Telecom & +0.049 & $-$0.010 [$-$0.066, +0.040] \\
OpenRCA & All & BARO vs $\epsilon$-Diagnosis & Market-1 & +0.019 & $-$0.006 [$-$0.067, +0.054] \\
OpenRCA & All & BARO vs max-$|Z|$ & Telecom & $-$0.049 & +0.003 [$-$0.086, +0.094] \\
OpenRCA & All & CD-1min vs max-$|Z|$ & Market-1 & +0.000 & $-$0.088 [$-$0.214, +0.030] \\
OpenRCA & All & CD-1min vs max-$|Z|$ & Market-2 & $-$0.007 & +0.000 [$-$0.145, +0.160] \\
OpenRCA & All & $\epsilon$-Diagnosis vs RCD & Market-1 & $-$0.013 & +0.017 [$-$0.052, +0.081] \\
OpenRCA & All & $\epsilon$-Diagnosis vs RCD & Market-2 & $-$0.023 & +0.039 [$-$0.033, +0.116] \\
OpenRCA & All & max-$|Z|$ vs RCD & Telecom & +0.039 & $-$0.009 [$-$0.099, +0.082] \\
PetShop & Published & BARO vs Counterfactual & Low & +0.000 & +0.018 [$-$0.173, +0.208] \\
PetShop & All & BARO vs Counterfactual & Low & +0.000 & +0.018 [$-$0.173, +0.208] \\
PetShop & All & Traversal vs Ranked correlation & Low & +0.000 & $-$0.012 [$-$0.238, +0.208] \\
\end{longtable}
}


\subsection*{Survey coding}
Per-paper answers to the four report-practice items, joined to the paper list. Source: \texttt{paper/data/survey/papers.csv}.

{\scriptsize\setlength{\tabcolsep}{3pt}
\begin{longtable}{lllcp{5.2cm}}
\caption{The surveyed papers. Per-system results: Yes if the paper evaluates more than one system and reports a score for each; Single system if it evaluates one system; Excluded if no full text was obtained (not counted in the paper). Location: the table, figure, or passage checked. Venue names longer than 30 characters are truncated. The italic line under each row is the paper title.}\\\label{supp:s11-survey}\\
\toprule
ID & Year & Venue & Per-system results & Location \\
\midrule
\endfirsthead
\caption{The surveyed papers. Per-system results: Yes if the paper evaluates more than one system and reports a score for each; Single system if it evaluates one system; Excluded if no full text was obtained (not counted in the paper). Location: the table, figure, or passage checked. Venue names longer than 30 characters are truncated. The italic line under each row is the paper title.}\\
\toprule
ID & Year & Venue & Per-system results & Location \\
\midrule
\endhead
\midrule
\multicolumn{5}{r}{\emph{Continued on next page}} \\
\bottomrule
\endfoot
\bottomrule
\endlastfoot
baro2024 & 2024 & FSE 2024 (Proc. ACM Softw.\ldots & Yes & Table 2 caption (cf. Table 1/3/4 captions, all "three datasets"/"three benchmark microservice systems") \\
\multicolumn{5}{l}{\scriptsize\textit{BARO: Robust Root Cause Analysis for Microservices via Multivariate Bayesian Online Change Point Detection}} \\
circa2022 & 2022 & KDD 2022 & Yes & Table 1 caption; cf. Table 3 caption "Performance of different methods on D\_O" \\
\multicolumn{5}{l}{\scriptsize\textit{Causal Inference-Based Root Cause Analysis for Online Service Systems with Intervention Recognition}} \\
rcd2022 & 2022 & NeurIPS 2022 & Yes & Table 1 caption; cf. Table 2 caption "summary and top-7 recall of epsilon-Diagnosis and RCD on data collected for three outages from a..." \\
\multicolumn{5}{l}{\scriptsize\textit{Root Cause Analysis of Failures in Microservices through Causal Discovery}} \\
tracerca2021 & 2021 & IWQoS 2021 & Yes & Sec. IV-B (discussion of Table II) \\
\multicolumn{5}{l}{\scriptsize\textit{Practical Root Cause Localization for Microservice Systems via Trace Analysis}} \\
microrank2021 & 2021 & WWW 2021 & Excluded (no full text) & -- \\
\multicolumn{5}{l}{\scriptsize\textit{MicroRank: End-to-End Latency Issue Localization with Extended Spectrum Analysis in Microservice Environments}} \\
rcaeval2025 & 2025 & WWW 2025 (Companion\ldots & Single system & Table 6 caption, Sec 5 (p.4) \\
\multicolumn{5}{l}{\scriptsize\textit{RCAEval: A Benchmark for Root Cause Analysis of Microservice Systems with Telemetry Data}} \\
fang2026 & 2026 & FSE 2026 (Proc. ACM Softw.\ldots & Yes & Sec. 3.2.2 Answering RQ2 \\
\multicolumn{5}{l}{\scriptsize\textit{Rethinking the Evaluation of Microservice RCA with a Fault Propagation-Aware Benchmark}} \\
openrca2025 & 2025 & ICLR 2025 & Yes & Text at the discussion of Figure 5 \\
\multicolumn{5}{l}{\scriptsize\textit{OpenRCA: Can Large Language Models Locate the Root Cause of Software Failures?}} \\
petshop2024 & 2024 & CLeaR 2024 & Yes & Table 3 \\
\multicolumn{5}{l}{\scriptsize\textit{The PetShop Dataset --- Finding Causes of Performance Issues across Microservices}} \\
eadro2023 & 2023 & ICSE 2023 & Yes & Sec. V-E RQ1 (Table II lists TT and SN side by side) \\
\multicolumn{5}{l}{\scriptsize\textit{Eadro: An End-to-End Troubleshooting Framework for Microservices on Multi-source Data}} \\
nezha2023 & 2023 & ESEC/FSE 2023 & Excluded (no full text) & -- \\
\multicolumn{5}{l}{\scriptsize\textit{Nezha: Interpretable Fine-Grained Root Causes Analysis for Microservices on Multi-modal Observability Data}} \\
mabc2024 & 2024 & EMNLP 2024 (Findings) & Yes & Table 4 \\
\multicolumn{5}{l}{\scriptsize\textit{mABC: Multi-Agent Blockchain-inspired Collaboration for Root Cause Analysis in Micro-Services Architecture}} \\
mulan2024 & 2024 & WWW 2024 & Yes & Table 2 caption; cf. Table 3 (Online Boutique)/Table 4 (Train Ticket) captions \\
\multicolumn{5}{l}{\scriptsize\textit{MULAN: Multi-modal Causal Structure Learning and Root Cause Analysis}} \\
ocean2024 & 2025 & IEEE BigData 2025 & Yes & Table II caption (Table III: Online Boutique; Table V: Train Ticket, same format) \\
\multicolumn{5}{l}{\scriptsize\textit{Online Multi-modal Root Cause Identification in Microservice Systems (OCEAN)}} \\
tvdiag2024 & 2025 & ACM TOSEM & Yes & Discussion of Table 4 \\
\multicolumn{5}{l}{\scriptsize\textit{TVDiag: A Task-oriented and View-invariant Failure Diagnosis Framework for Microservice-based Systems with Multimodal Data}} \\
mrca2024 & 2024 & ASE 2024 & Excluded (no full text) & -- \\
\multicolumn{5}{l}{\scriptsize\textit{MRCA: Metric-level Root Cause Analysis for Microservices via Multi-Modal Data}} \\
chase2024 & 2024 & arXiv preprint & Yes & Text at the discussion of Table II (overall performance comparison, one row per dataset) \\
\multicolumn{5}{l}{\scriptsize\textit{CHASE: A Causal Hypergraph based Framework for Root Cause Analysis in Multimodal Microservice Systems}} \\
howfar2024 & 2024 & ASE 2024 & Yes & Table 5 caption \\
\multicolumn{5}{l}{\scriptsize\textit{Root Cause Analysis for Microservice System based on Causal Inference: How Far Are We?}} \\
lemmarca2024 & 2024 & OpenReview & Yes & Table 4 caption (cf. Table 3: multi-modality Product Review/Cloud Computing results; Table 5: WADI; Table 6: caption line 793) \\
\multicolumn{5}{l}{\scriptsize\textit{LEMMA-RCA: A Large Multi-modal Multi-domain Dataset for Root Cause Analysis}} \\
sparserca2024 & 2024 & ISSRE 2024 & Single system & Sec. VI.A.1 Datasets, p.9 \\
\multicolumn{5}{l}{\scriptsize\textit{SparseRCA: Unsupervised Root Cause Analysis in Sparse Microservice Testing Traces}} \\
tracecontrast2024 & 2024 & ICSE 2024 & Single system & Sec. 4.1.1 Dataset, p.7 \\
\multicolumn{5}{l}{\scriptsize\textit{Trace-based Multi-Dimensional Root Cause Localization of Performance Issues in Microservice Systems}} \\
tracediag2023 & 2023 & ESEC/FSE 2023 & Single system & Sec. 5.1.1 Dataset, p.13 \\
\multicolumn{5}{l}{\scriptsize\textit{TraceDiag: Adaptive, Interpretable, and Efficient Root Cause Analysis on Large-Scale Microservice Systems}} \\
run2024 & 2024 & AAAI 2024 & Yes & Text (Tables 2 and 3: sock-shop; Figure 6: synthetic dataset, reported separately) \\
\multicolumn{5}{l}{\scriptsize\textit{Root Cause Analysis In Microservice Using Neural Granger Causal Discovery}} \\
dynacausal2025 & 2025 & arXiv preprint & Yes & Table 2 caption (datasets D1 and D2) \\
\multicolumn{5}{l}{\scriptsize\textit{DynaCausal: Dynamic Causality-Aware Root Cause Analysis for Distributed Microservices}} \\
causalrca2023 & 2023 & Journal of Systems and\ldots & Single system & Sec. I Introduction, p.2 (also Sec. Threats to Validity) \\
\multicolumn{5}{l}{\scriptsize\textit{CausalRCA: Causal inference based precise fine-grained root cause localization for microservice applications}} \\
toomanycooks2025 & 2025 & ISSRE 2025 & Yes & Table IV caption \\
\multicolumn{5}{l}{\scriptsize\textit{Too Many Cooks: Assessing the Need for Multi-Source Data in Microservice Failure Diagnosis}} \\
\end{longtable}
}


\subsection*{Cross-release, all pairs}
RE1 versus RE2 differences for every pair, including probes and shipped baselines. Source: \texttt{paper/data/rq1/cross\_release.csv}.

{\scriptsize\setlength{\tabcolsep}{3pt}
\begin{longtable}{llrrcc}
\caption{Cross-release comparison for every pair, including probes (the all-pairs layer). $\Delta$ is accuracy on the RE1 release minus the paired method, and the same contrast on the RE2 release, each with its 95\% CI. Same sign marks agreement of the two differences; CIs opposite marks intervals that lie on opposite sides of zero.}\\\label{supp:s12-cross-release}\\
\toprule
Pair & App & $\Delta$ RE1 [CI] & $\Delta$ RE2 [CI] & Same sign & CIs opposite \\
\midrule
\endfirsthead
\caption{Cross-release comparison for every pair, including probes (the all-pairs layer). $\Delta$ is accuracy on the RE1 release minus the paired method, and the same contrast on the RE2 release, each with its 95\% CI. Same sign marks agreement of the two differences; CIs opposite marks intervals that lie on opposite sides of zero.}\\
\toprule
Pair & App & $\Delta$ RE1 [CI] & $\Delta$ RE2 [CI] & Same sign & CIs opposite \\
\midrule
\endhead
\midrule
\multicolumn{6}{r}{\emph{Continued on next page}} \\
\bottomrule
\endfoot
\bottomrule
\endlastfoot
Alert-count vs BARO & OB & $-$0.344 [$-$0.440, $-$0.248] & +0.067 [$-$0.056, +0.189] &  &  \\
 & SS & $-$0.472 [$-$0.576, $-$0.360] & +0.322 [+0.211, +0.433] &  & \checkmark \\
 & TT & $-$0.512 [$-$0.608, $-$0.408] & $-$0.456 [$-$0.567, $-$0.344] & \checkmark &  \\
Alert-count vs CausalRCA & OB & $-$0.224 [$-$0.320, $-$0.128] & +0.004 [$-$0.111, +0.122] &  &  \\
 & SS & +0.200 [+0.099, +0.301] & +0.241 [+0.130, +0.356] & \checkmark &  \\
 & TT & $-$0.088 [$-$0.144, $-$0.029] & +0.004 [$-$0.089, +0.100] &  &  \\
Alert-count vs CD-1min & OB & $-$0.336 [$-$0.448, $-$0.224] & $-$0.522 [$-$0.644, $-$0.389] & \checkmark &  \\
 & SS & $-$0.552 [$-$0.648, $-$0.448] & $-$0.556 [$-$0.656, $-$0.456] & \checkmark &  \\
 & TT & $-$0.232 [$-$0.328, $-$0.144] & $-$0.578 [$-$0.689, $-$0.467] & \checkmark &  \\
Alert-count vs CIRCA & OB & $-$0.280 [$-$0.392, $-$0.176] & $-$0.456 [$-$0.589, $-$0.322] & \checkmark &  \\
 & SS & $-$0.328 [$-$0.448, $-$0.208] & $-$0.389 [$-$0.511, $-$0.267] & \checkmark &  \\
 & TT & $-$0.280 [$-$0.368, $-$0.192] & $-$0.333 [$-$0.456, $-$0.211] & \checkmark &  \\
Alert-count vs $\epsilon$-Diagnosis & OB & +0.360 [+0.272, +0.448] & +0.178 [+0.089, +0.267] & \checkmark &  \\
 & SS & +0.232 [+0.120, +0.344] & +0.233 [+0.111, +0.356] & \checkmark &  \\
 & TT & +0.048 [+0.016, +0.088] & +0.211 [+0.133, +0.300] & \checkmark &  \\
Alert-count vs max-$|Z|$ & OB & $-$0.112 [$-$0.224, +0.000] & $-$0.567 [$-$0.689, $-$0.433] & \checkmark &  \\
 & SS & $-$0.368 [$-$0.488, $-$0.248] & $-$0.456 [$-$0.589, $-$0.322] & \checkmark &  \\
 & TT & $-$0.376 [$-$0.472, $-$0.280] & $-$0.200 [$-$0.322, $-$0.067] & \checkmark &  \\
BARO vs CausalRCA & OB & +0.120 [+0.043, +0.195] & $-$0.063 [$-$0.159, +0.041] &  &  \\
 & SS & +0.672 [+0.592, +0.752] & $-$0.081 [$-$0.163, $-$0.004] &  & \checkmark \\
 & TT & +0.424 [+0.331, +0.515] & +0.459 [+0.341, +0.574] & \checkmark &  \\
BARO vs CD-1min & OB & +0.008 [$-$0.080, +0.096] & $-$0.589 [$-$0.700, $-$0.478] &  &  \\
 & SS & $-$0.080 [$-$0.136, $-$0.024] & $-$0.878 [$-$0.944, $-$0.811] & \checkmark &  \\
 & TT & +0.280 [+0.160, +0.392] & $-$0.122 [$-$0.233, $-$0.011] &  & \checkmark \\
BARO vs CIRCA & OB & +0.064 [$-$0.032, +0.160] & $-$0.522 [$-$0.644, $-$0.400] &  &  \\
 & SS & +0.144 [+0.064, +0.232] & $-$0.711 [$-$0.800, $-$0.611] &  & \checkmark \\
 & TT & +0.232 [+0.136, +0.328] & +0.122 [+0.011, +0.233] & \checkmark &  \\
BARO vs $\epsilon$-Diagnosis & OB & +0.704 [+0.624, +0.784] & +0.111 [+0.022, +0.200] & \checkmark &  \\
 & SS & +0.704 [+0.616, +0.784] & $-$0.089 [$-$0.189, +0.000] &  &  \\
 & TT & +0.560 [+0.472, +0.640] & +0.667 [+0.567, +0.767] & \checkmark &  \\
BARO vs max-$|Z|$ & OB & +0.232 [+0.136, +0.328] & $-$0.633 [$-$0.733, $-$0.522] &  & \checkmark \\
 & SS & +0.104 [+0.032, +0.176] & $-$0.778 [$-$0.856, $-$0.689] &  & \checkmark \\
 & TT & +0.136 [+0.048, +0.224] & +0.256 [+0.133, +0.378] & \checkmark &  \\
CausalRCA vs CD-1min & OB & $-$0.112 [$-$0.197, $-$0.027] & $-$0.526 [$-$0.619, $-$0.430] & \checkmark &  \\
 & SS & $-$0.752 [$-$0.819, $-$0.680] & $-$0.796 [$-$0.867, $-$0.722] & \checkmark &  \\
 & TT & $-$0.144 [$-$0.235, $-$0.056] & $-$0.581 [$-$0.678, $-$0.481] & \checkmark &  \\
CausalRCA vs CIRCA & OB & $-$0.056 [$-$0.144, +0.035] & $-$0.459 [$-$0.567, $-$0.356] & \checkmark &  \\
 & SS & $-$0.528 [$-$0.613, $-$0.440] & $-$0.630 [$-$0.722, $-$0.530] & \checkmark &  \\
 & TT & $-$0.192 [$-$0.280, $-$0.104] & $-$0.337 [$-$0.456, $-$0.219] & \checkmark &  \\
CausalRCA vs $\epsilon$-Diagnosis & OB & +0.584 [+0.512, +0.659] & +0.174 [+0.089, +0.259] & \checkmark &  \\
 & SS & +0.032 [$-$0.051, +0.117] & $-$0.007 [$-$0.104, +0.089] &  &  \\
 & TT & +0.136 [+0.096, +0.179] & +0.207 [+0.152, +0.267] & \checkmark &  \\
CausalRCA vs max-$|Z|$ & OB & +0.112 [+0.008, +0.216] & $-$0.570 [$-$0.667, $-$0.470] &  & \checkmark \\
 & SS & $-$0.568 [$-$0.661, $-$0.475] & $-$0.696 [$-$0.789, $-$0.596] & \checkmark &  \\
 & TT & $-$0.288 [$-$0.379, $-$0.195] & $-$0.204 [$-$0.311, $-$0.100] & \checkmark &  \\
CD-1min vs CIRCA & OB & +0.056 [$-$0.040, +0.152] & +0.067 [$-$0.044, +0.178] & \checkmark &  \\
 & SS & +0.224 [+0.144, +0.304] & +0.167 [+0.089, +0.256] & \checkmark &  \\
 & TT & $-$0.048 [$-$0.160, +0.072] & +0.244 [+0.133, +0.356] &  &  \\
CD-1min vs $\epsilon$-Diagnosis & OB & +0.696 [+0.616, +0.776] & +0.700 [+0.600, +0.800] & \checkmark &  \\
 & SS & +0.784 [+0.712, +0.856] & +0.789 [+0.689, +0.878] & \checkmark &  \\
 & TT & +0.280 [+0.200, +0.360] & +0.789 [+0.700, +0.867] & \checkmark &  \\
CD-1min vs max-$|Z|$ & OB & +0.224 [+0.120, +0.336] & $-$0.044 [$-$0.122, +0.033] &  &  \\
 & SS & +0.184 [+0.112, +0.256] & +0.100 [+0.022, +0.178] & \checkmark &  \\
 & TT & $-$0.144 [$-$0.248, $-$0.032] & +0.378 [+0.256, +0.500] &  & \checkmark \\
CIRCA vs $\epsilon$-Diagnosis & OB & +0.640 [+0.552, +0.720] & +0.633 [+0.533, +0.744] & \checkmark &  \\
 & SS & +0.560 [+0.464, +0.648] & +0.622 [+0.500, +0.733] & \checkmark &  \\
 & TT & +0.328 [+0.248, +0.408] & +0.544 [+0.444, +0.644] & \checkmark &  \\
CIRCA vs max-$|Z|$ & OB & +0.168 [+0.064, +0.280] & $-$0.111 [$-$0.211, $-$0.011] &  & \checkmark \\
 & SS & $-$0.040 [$-$0.144, +0.064] & $-$0.067 [$-$0.156, +0.011] & \checkmark &  \\
 & TT & $-$0.096 [$-$0.192, +0.008] & +0.133 [+0.000, +0.267] &  &  \\
$\epsilon$-Diagnosis vs max-$|Z|$ & OB & $-$0.472 [$-$0.560, $-$0.376] & $-$0.744 [$-$0.833, $-$0.656] & \checkmark &  \\
 & SS & $-$0.600 [$-$0.696, $-$0.496] & $-$0.689 [$-$0.800, $-$0.567] & \checkmark &  \\
 & TT & $-$0.424 [$-$0.512, $-$0.336] & $-$0.411 [$-$0.511, $-$0.311] & \checkmark &  \\
\end{longtable}
}




\end{document}
